\documentclass[11pt]{jsarticle}
\usepackage[top=20mm,bottom=20mm,left=25mm,right=25mm]{geometry}
\usepackage[compact]{titlesec}
\usepackage{xcolor}
\usepackage{amssymb}  % \lesssim, \gtrsim 等
\titlespacing*{\section}{0pt}{6pt plus 2pt minus 1pt}{3pt plus 1pt}
\titlespacing*{\subsection}{0pt}{4pt plus 1pt minus 1pt}{2pt plus 1pt}

% 各セクションは新しいページから始める
% （titlesec の \sectionbreak がこの環境で効かないため \section を直接上書き）
\let\litsectionorig\section
\renewcommand{\section}{\clearpage\litsectionorig}

% 論文名（Author, year）の見出しの色とマクロ
% 論文の切れ目がパッとわかるよう、色は派手め＋少し大きめにする
\definecolor{citecolor}{HTML}{1565C0}
\newcommand{\litref}[1]{{\large\bfseries\color{citecolor}#1}}

% 文献整理用の作業ファイル．投稿原稿ではない．
%
% 【コメント規約】
%   ==== AI-DRAFT: citekey ==== で始まるブロック
%       → AI 要約．未確認．本文には昇格させない．
%   本文（コメント解除・修正済み）
%       → 人間が原典で確認済み．末尾に確認ページを記す．
%       → 昇格本文は「問い／手法／主要知見／限界・適用範囲／本研究との接点」の
%         箇条書きで記す（リファレンス文書として一覧性を優先）．
%       → 論文名の見出しは \litref{Author (year)} で記す．
%
% 【status の読み方】（手動管理不要．ファイルを横断して導出できる）
%   candidate = .bib にあるが，このファイルに本文記述がない
%   checked   = このファイルの本文にメモがある（AI-DRAFT ブロックは対象外）
%   used      = paper/ 内の原稿で \cite されている
%
% 書誌情報の正本は Zotero．citekey は refs/ 内の .bib と一致させる．
% evidence_table.csv は必要時に生成する出力であり，手動で維持しない．

\title{文献整理マスター}
\author{}
\date{\today}

\begin{document}

\maketitle

% ================================================================
\section{揚水発電と上部調整池の運用}
% ================================================================

\subsection*{主要な問い}

\subsection*{論文メモ}

\noindent\litref{Zhu et al. (2023)}\quad 揚水発電所の鉛直管型流入出口（ゴミ除け格子断面）の流出・流入時の瞬時流速を物理模型で精査．
\begin{itemize}
  \item \textbf{問い}：鉛直管型ポートのゴミ除け格子断面で，流出・流入条件の瞬時流速は時間平均とどの程度異なるか．
  \item \textbf{手法}：幾何縮尺 1:35 の Froude 相似物理模型（水槽 $10 \times 10$~m），ADV（100~Hz），格子孔高さ方向 10 点，14 試験条件．
  \item \textbf{主要知見}：
  \begin{itemize}
    \item 流出時：瞬時流速が激しく変動し，局所最大で時間平均の 3 倍超．格子孔断面は鉛直に 3 ゾーン（上から正 28~\%，正負交番 52~\%，負正交番 20~\%），支配周波数 0.02--0.06~Hz．
    \item 流入時：瞬時流速の変動は小さく，時間平均の 0.9--1.1 倍．
    \item 乱流強度は配管内で逐次増大（トンネル 0.28 →屈曲部 0.61 →格子断面 0.85）．流出経路の 2 回の 90$^\circ$ 屈曲が主因．
    \item 支配周波数は 0.1~Hz 未満で固有振動数を大幅に下回り共振リスクは低いが，交番応力による疲労破壊の可能性を指摘．
  \end{itemize}
  \item \textbf{限界・適用範囲}：揚水発電専用の鉛直管構成．計測対象は格子部の局所構造で，池全体の表面流況とはスケール・形状が異なる．
  \item \textbf{本研究との接点}：流入・流出で流況の非対称性が存在する直接的根拠．流出時の高乱流強度・逆流域が表面流況にも影響しうる．本研究のポートは水平矩形断面で形状は異なるが，流入・流出非対称性の物理的説明として引用できる．
\end{itemize}
\noindent（確認: p.1, p.3--8）

\noindent\litref{Rehman et al. (2015)}\quad（周辺）揚水発電の技術総説（全 13~p 通読）．\texttt{introduction.tex} 冒頭のプレースホルダ \texttt{PSP\_background1--4} を埋めるための背景文献．三階層方針でいう\textbf{周辺}で，背景 1--2 文を \verb|\cite| で支えるだけの用途．
\begin{itemize}
  \item \textbf{内容}：PHES の技術総説．世界の設備容量，立地選定，歴史，hybrid 構成（風力/太陽光/海水），経済性を横断的に整理．純揚水/混合式は EU 統計として 1 箇所出るのみで，\textbf{可変速機の議論はない}．上池の流況・堆砂・取放水構造物には一切触れない．
  \item \textbf{引用できる箇所}：
  \begin{itemize}
    \item \textbf{PHES は世界の大規模蓄電容量の 99~\% 超，約 127~GW}（EPRI, 2012 年時点）＝第 1 文の根拠．
    \item PHES の利点 3 つ（Hino \& Lejeune）：起動停止が柔軟で応答が速い／負荷変動を追随し急変にも適応／周波数を調整し電圧安定性を保つ＝第 2 文（短時間の出力追従・系統安定化）の根拠．
    \item 1890 年代のアルプスが最初で\textbf{1950 年代以降は単一の可逆ポンプ水車が主流}＝本研究が可逆単一ポートを扱うことの実機側の裏づけ．適地が稀になりつつあるとの指摘もあり，既存上池の運用・維持管理が重要という論理にも使える．
  \end{itemize}
  \item \textbf{限界・適用範囲}：「変動性再エネの拡大に伴い運転頻度・起動停止が増える」は本総説には\textbf{明示的にない}．あるのは上記の能力の記述まで．頻度の増加自体を主張するなら Hunt 2020 側か運用寄りの別文献が要る．また統計は 2009--2013 年で 2026 年の原稿には 10 年以上古い．定性的主張（99~\% を占める）は今も妥当だが GW 値を引くなら IEA/IRENA の近年資料へ．本文内でも 127~GW／Table~1 の 104~GW（2009）／「約 100~GW」が基準年混在のまま並び，数値の内的整合性が弱い．
\end{itemize}
\noindent（確認: p.586--598，全 13~p 通読）

\noindent\litref{Hunt et al. (2020)}\quad（周辺）揚水発電の配置形式の総説（全 14~p 通読）．Rehman 2015 と同じく \texttt{PSP\_background1--4} 用．役割分担は Rehman ＝容量・重要性（99~\%，127~GW），本論文＝運転形態（反転頻度・可変速機）．統計も 2019 年時点で Rehman より新しい．
\begin{itemize}
  \item \textbf{内容}：§2 で既存 PHS を 4 軸（貯蔵サイクル／ポンプ水車の速度・型式／用途／配置）で分類．§3 以降は新配置（CCSPHS, CHPHS）の提案とザンベジ流域の GIS 評価で\textbf{本研究とは無関係}．使えるのは §2 のみ．海水・地下・廃坑利用は ``less common configurations'' として 1 文で列挙されるだけで分類の本体ではない．
  \item \textbf{引用できる箇所}：
  \begin{itemize}
    \item \textbf{Hourly PHS は短絡モードで運転でき 1 日 100 回を超える反転が可能}（例：オーストリア Kops II，原典は Goekler \& Meusburger 2009）．＝可逆単一ポートを主題にすることの実機側の動機として最も強い．ただし「再エネ拡大→運転頻度増」という因果は §1（再エネ→柔軟性需要）＋§2.1（HPHS が ancillary service を担う）を繋いだもので，一文で書かれてはいない．引く際は該当の 2 箇所を併記すること．
    \item 貯蔵サイクル分類（Table~1）：Hourly／Daily（現在最も一般的）／Weekly／Seasonal／Pluri-annual．HPHS と DPHS は容量が同じで，区別は運用サイクルによる．
    \item \textbf{可変速ポンプ水車}：固定速は定格でしか運転できないが可変速は 60--100~\% で可変．コスト約 30~\% 高．ternary（Pelton＋Francis，回転反転が不要で切替が速い）・quaternary（Gordon Butte で 20~MW/秒）も．＝Rehman にはない論点．
    \item 電池コストの低下により\textbf{日サイクルの PHS は将来の競争力を失う可能性}があり，PHS は長サイクルか高速 ancillary service の両極へ向かう．＝高頻度反転が今後重要になる論拠．
    \item 日本への言及 2 箇所：多数のダムが可逆水車を備える（原子力の硬直性による．ただし本文の ``reduces the need for daily storage'' は論理が逆で編集上の誤りとみられる）／神流川の pump-back．
  \end{itemize}
  \item \textbf{本研究との接点}：本研究の basin（矩形の人工上池，流入は port のみ）は\textbf{closed-loop} に対応する（大きな水源から離れ流入が限られる小さな人工池．環境影響は小さいが日〜週サイクルに限られる）．ただし\textbf{池の幾何（浅い/深い，平面形）には踏み込まない}ので Methods の prototype-informed の補強には使えない．上池の流況・堆砂も対象外（堆砂は「集水域が小さいので速度が小さい」の 1 行のみ）．
\end{itemize}
\noindent（確認: art.\ 109914，全 14~p 通読）

\noindent\litref{M\"uller (2012)}\quad 本研究のラボスケール直系の先行研究（EPFL TH\`ESE NO 5471，本体 155~p，学位論文）．流況データは本論文の Ch5 にしか存在しない．
\begin{itemize}
  \item \textbf{章構成}（目次・Ch7 で確認済）：Ch3 現地 ADCP＋数値 ＝ M\"uller 2018（会議版 M\"uller 2010），Ch4 濁度監視 ＝ M\"uller 2014，Ch5 ジェット挙動・流況・運動エネルギー ＝ \textbf{論文化されていない}，Ch6 土砂交換率 ＝ M\"uller 2017，Ch7 結論．読解範囲は Ch5（PDF p.71--96）と Ch7（PDF p.121--127）を精読し，Ch3/4/6 は論文版で代替．
  \item \textbf{問い}：揚水発電の流入・流出シーケンスは貯水池の流況と浮遊砂の沈降挙動をどう変えるか．
  \item \textbf{手法}：3 本柱＝現地計測（Ch3）／濁度監視（Ch4）／室内実験（Ch5, Ch6）＋数値（ANSYS-CFD, $k$-$\varepsilon$）．Ch5 の室内実験：main basin $4.0 \times 2.0 \times 1.5$~m（水深 1.15~m）＋mixing tank を可逆回路で接続．前壁に楕円ベルマウス（内径 4.8~cm，$z_i = 0.25/0.5/0.75$~m）．$Q = 0.3$--$1.1$~L/s，$\mathrm{Re}_i = 7560$--$29180$．UVP で水平 2D 流速（ただし 1 象限 $1.0 \times 2.0$~m の 28 点のみを 3 レベルで，線形補間して図示）．Froude 相似．長さは $B_\mathrm{MB}$，時間は $t_m = V/Q$ で正規化．閉鎖系のため水位が最大 0.18~m 変動．なお IN 系列＝池へ流入＝本研究の inflow（同じ向き）．
  \item \textbf{主要知見（Ch5 ＝流況，本研究に直結）}：
  \begin{itemize}
    \item inflow の流況：ジェット＋水平 2 つの再循環セル＋鉛直 1 つのセル．ジェットは軸に沿って微振動．高流量（$Q \geq 0.9$）で再循環セルが計測領域全体を占め，低流量では境界が $s/B = 0.4$ に留まる．→本研究の予備観察「偏向ジェットと二循環構造」と一致．
    \item outflow の流況：取水口近傍のみが取水口を向き，後方は準静止で優先的な流れの向きがない．数値では $v/v_0 < 0.08$（$s/D_j = 1$）．→本研究の予備観察「basin 全体の単一循環」とは異なる．
    \item 運動エネルギーの非対称：$E_\mathrm{kin} = (1/2)m\sum v_n^2$ は本研究の $E(t) = \langle |u|^2 \rangle_\Omega$ と実質同じ量．数値モデルでは outflow は先行 inflow の 60~\%．実験では inflow ピークの 5--10~\% まで落ち，平均で 12.8--29~\%．数値と実験が大きく食い違い未決着（著者はサンプリング体積の制約を挙げる）．
    \item 定常到達は inflow の方が遅い：$t_{P,\mathrm{IN}} > t_{P,\mathrm{OUT}}$（全流量），10--68 分，$t_P/t_m = 0.099$--$0.136$．「ジェットの形成は取水の流れ場より時間がかかる」ため inflow が critical case で，$K_{t_P}$ の基準に採用．
    \item ヒステリシス：outflow の流況はしばしば先行 inflow に影響され，大きな再循環セルは散逸まで持続．高頻度ほど残留セルによりエネルギー入力が増す．散逸所要は $Q = 1.1$ で $0.8\,t_{P,\mathrm{IN}}$，$Q = 0.3$ で $0.2\,t_{P,\mathrm{IN}}$．→本研究が各 run を静水から始め $E(t)$ で残留判定する設計の直接の根拠．
    \item ジェット偏向は $z_i/B = 0.125$ と $0.375$ で発生（$0.25$ では起きない）．数値でも再現され，原因は overflow section による幾何的非対称と説明．
    \item 底層（$z/B = 0.125$）で $Q > 0.7$~L/s のとき取水口へ向かう backflow．
  \end{itemize}
  \item \textbf{主要知見（Ch3, Ch6 ＝現地・土砂．論文版と重複するため要点のみ）}：
  \begin{itemize}
    \item 現地でも同じ非対称：pumping は 30~m で 6~cm/s，150~m で影響消失．turbine は 150~m でも 12~cm/s．
    \item TKE 入力は風の約 25 倍（経験式による推定）．8 ヶ月で約 45{,}000~t が 2 湖間を移動．
    \item 初期 $t/t_m < 0.2$ で約 60~\% が沈降しサイクル規模に依存しない．低流量で SSR 10--40~\% 増，高流量で 50--60~\% 増（高頻度と組合せで最大 80~\%）．ポート位置は底面・水面近傍が中間より 20~\% 高い．
    \item 土砂収支はほぼ中立．Ch7.1.4 は「両槽の濃度が同程度なら」という条件を明示するが，濃度が大きく異なる場合は実験していないと Ch7.1.5 が認めており，条件節は推測であって実証ではない．
  \end{itemize}
  \item \textbf{限界・適用範囲}：浅水槽ではない（$h/B \approx 0.58$，機構は鉛直再循環を含む）．UVP は 1 象限 28 点 $\times$ 3 レベルで空間解像度が粗く，表面 PIV のような面的データではない．目的は排砂・土砂管理であり表面流レジームの分類ではない．ポート位置について Ch5 の運動エネルギー（$0.125$ は低い）と Ch6 の SSR（$0.125$ も 20~\% 高い）で傾向が食い違う点は未解決．数値と実験の不一致が複数箇所．
  \item \textbf{本研究との接点}：
  \begin{itemize}
    \item 「矩形槽＋可逆ポート＋ポート鉛直位置＋2D 流速場」という設定が本研究に最も近い既往研究．
    \item inflow の流況（ジェット＋二循環）は本研究の予備観察と一致するが，outflow は一致しない（彼らは後方停滞，本研究は全体循環）．深さの差（$h/B \approx 0.58$ vs $h/W \approx 0.02$--$0.05$）で，浅い系では吸込流が 2 次元的に広がり全体に及ぶ可能性．→本研究の独自性を主張できる余地．
    \item $E_\mathrm{kin}$ の inflow/outflow 比（実験値 12.8--29~\%）は，本研究の $E(t)$ と $I_\mathrm{asym}$ を比較すべき先行値．
    \item ヒステリシスの観察は本研究の静水開始プロトコルを支持する．
    \item ポートの鉛直位置が効くのは Chagdali 2025 の $Z_\mathrm{jet}/H$ と同じ論点（Chagdali は未確認）．
  \end{itemize}
\end{itemize}
\noindent（確認: Ch5 = p.71--96, Ch7 = p.121--127）

\noindent\litref{M\"uller et al. (2018)}\quad 本研究のモード非対称を現地スケールで裏づける論文（学位論文 Ch3 の論文版）．
\begin{itemize}
  \item \textbf{問い}：揚水発電運転下の実貯水池で，取水/放流口近傍の流況はどうなっているか．ポンプ運転とタービン運転で流れ場はどう違うか．
  \item \textbf{手法}：Grimsel 2 揚水発電所の下池 Lake Grimsel のみ（Oberaar は上池として言及されるだけで計測なし）．円筒形の水中取放水口（有効高 6.25~m，径 21.7~m，10 セクター，案内壁）前面の湖底に ADCP 3 台，2008 年 9 月（5 日）と 11 月（2 週間）．運転データとの FFT 主周波数比較．加えて 3D ANSYS-CFX（$k$-$\varepsilon$，非定常，140 万セル，水位一定・清水・等温，Coriolis 無視）．なお turbine mode ＝池へ吐出＝本研究の inflow，pumping ＝吸込＝本研究の outflow（逆対応）．
  \item \textbf{主要知見}：
  \begin{itemize}
    \item モード非対称（核心）：pumping（吸込）は取水口ごく近傍にしか影響せず，数値モデルでも最寄りのプロファイラ以外に測定可能な流速を生じない．対して turbine（吐出）は backflow 域と大規模循環セルを形成し，吐出口から 150~m 離れても高流速が認められる．
    \item 流速信号と運転サイクルに明瞭な相関．卓越周期 $1/f = 24$~h（日周の揚水サイクル）．セイシュの周波数帯にピークはなく，内部振動は流況計測に影響していない．
    \item 定常流況の発達は連続タービン運転の約 2.5~h 後＝150 分．ただしこれは数値モデルの結果で，現地計測からはタービンモードの時間発展を明確に述べられないと本文が認めている．
    \item TKE 入力は風の約 25 倍（Imboden 1980 と Anderson 2010 の経験式による推定であり計測値ではない）．日平均 $\eta_\mathrm{Turb} = 2\times10^{-3}$~W/m$^2$ 対 $\eta_a = 7.9\times10^{-5}$~W/m$^2$．強風 12~m/s と比べてもなお約 10 倍．→揚水運転の乱れを細粒土砂の沈降遅延に利用しうる，という本研究群全体の動機の定量的根拠．
    \item 無運転時にも深層は連続的に撹拌される．循環セルは発電停止後も回転を続ける（履歴効果）．
    \item 否定的結果：2 つの取水口（Grimsel 2 と Gelmer/Grimsel 1）の流れが繋がるのは 9~h 以上の連続タービン運転後で，周方向流入はよく整った射流を作らず急減衰．2 発電所を使った湖間の浮遊土砂輸送は実行可能な方策ではない．
    \item 温度差 $\pm 0.5^\circ$ では流況は根本的に変わらない．ただし土砂を含む水なら浮力効果がジェット軌道を大きく変えうると留保．
  \end{itemize}
  \item \textbf{限界・適用範囲}：水柱の大半が未計測＝装置から 20--30~m は信号損失，その上の約 30~m は水面までデータ不良（水-空気界面の信号増幅）．スペクトル解析は湖底上 20~m のみ．水平流速の測定誤差 $\pm 12$~mm/s，5 分サンプリング＋5 ステップ移動平均のため乱流量は測っていない．数値モデルの検証は VAW (1982) 物理模型のラック断面に対してのみで ADCP との定量検証ではない．複雑地形の天然湖であり制御された矩形形状ではない．深い湖（最大 100~m）で浅水流ではない．
  \item \textbf{本研究との接点}：
  \begin{itemize}
    \item 「同一ポートでも inflow は大規模循環・outflow は局所」を現地スケールで示した論文．本研究の問い (1) を実施設側から支える引用先．室内側の M\"uller 2017（inflow 系列で SSR が高い＝流入ジェットのエネルギー投入が主因）と機構的に整合する．
    \item ただし本研究の予備観察では outflow でも basin 全体の単一循環が生じており，「局所」とは異なる．深い湖では吸込の影響が距離とともに急減衰するのに対し，極端に浅く側壁で囲まれた basin では吸込流が全体に及びうる．＝深さ・閉鎖度に依存する差の可能性．そのまま「同じ現象」とは書けない．
    \item 「定常到達 2.5~h」は本研究の計測時間・静水開始判定の参照値になりうるが，スケールが桁違い（湖 100~m 深 vs basin 0.03--0.10~m）で直接換算は不可．
  \end{itemize}
\end{itemize}
\noindent（確認: p.109--124）



\noindent\litref{Jenzer Althaus et al. (2015)}\quad 矩形タンク（$W = 2$~m）＋前壁の取水口で，人工ジェットによる細粒土砂の浮遊維持と連続排出を検証．
\begin{itemize}
  \item \textbf{問い}：ダム前面の細粒土砂を人工ジェットで浮遊維持し，取水口から連続排出できるか．最適配置は．
  \item \textbf{手法}：角柱水槽（$4 \times 2 \times 1.5$~m，前壁＝ダム）の水理実験のみ（CFD なし）．水平面に正方配置した 4 本の垂直交差ジェットで回転流を誘起．比較対象はジェットなしの参照実験（配置間比較は Jenzer Althaus 2011 の博士論文）．UVP で 2 鉛直面の流速，濁度計で排砂量，4 時間．
  \item \textbf{主要知見}：
  \begin{itemize}
    \item ジェットで排砂比 ESR がほぼ倍増（上乗せ効率 $\phi = \mathrm{ESR}_j/\mathrm{ESR}_\mathrm{ref} \approx 1.5$--$2$）．
    \item 流動は axial / radial mixer 型に分かれ axial が排砂に有利．境界は底面クリアランス $h_c/W = 0.175$．
    \item 最適は幾何比のみで決まる（$h_c/W = 0.175$, $h_\mathrm{in}/W = 0.25$, $d_\mathrm{in}/W = 0.525$, $d_\mathrm{jet}/W = 0.15$, $\theta = 0$）．ジェットの momentum・Fr・Re は ESR にほぼ効かない＝動的量でなく幾何が支配．
    \item Froude 相似のみ（Re 非一致）．密度比を実機 $s \approx 2.68$ →模型 $s \approx 1.5$ へ意図的に下げ，浮遊砂輸送の相似性を確保（Yalin 1971）．特定実機の縮尺模型ではなく，一般的適用性の検証と位置づけ．
  \end{itemize}
  \item \textbf{限界・適用範囲}：著者自身が「流動パターンは壁に強く支配され，壁がなければセル形状は異なりうる」とし，実機（半無限空間）での成立は未証明の仮定と明記．$h_0/W = 0.6$ で浅水流ではない（本研究は $h/W \approx 0.02$--$0.05$ と 1--2 桁違い，鉛直循環セルが主役の流れ場）．駆動は人工ジェットでポート流入ジェットではない．
  \item \textbf{本研究との接点}：Fig.~8（$h_\mathrm{in}/W$--$h_c/W$ 平面を axial/radial に分割，判別不能域は白抜き）＋Fig.~9（同一軸の ESR 等高線）は，$h/a$ と流量でレジームマップを描く際の作法の先行例．幾何量を槽幅 $W$ で正規化する流儀（$h_\mathrm{in}/W$）は本研究の $h/a$（開口高さ基準）と対比可能．目的は排砂であり流況モードではない．§6（土砂）と接続．
\end{itemize}
\noindent（確認: art.\ 04014078）

\noindent\litref{Berm\'udez et al. (2017)}\quad 実施設（計画段階）の取水/放流口の配置・形状設計を，2D 浅水・3D RANS・物理模型の組合せで検討．
\begin{itemize}
  \item \textbf{問い}：狭く浅い貯水池において，揚水発電の取水/放流構造物の位置と形状をどう決めるべきか．
  \item \textbf{手法}：スペイン北西部 Belesar III（計画段階）の下部取水/放流口．(i) 遠方場：2D 浅水（Iber）で取水口位置の妥当性を評価．(ii) 近傍場：3D RANS（CFX，$k$-$\varepsilon$，rigid-lid，定常）＋1:40 Froude 物理模型（ADV）で潜没条件と形状代替案（開口 6→5 門）を比較．位置の代替案は振っていない．なお pumping ＝池から吸込＝本研究の outflow，generating ＝池へ吐出＝本研究の inflow（逆対応）．
  \item \textbf{主要知見}：
  \begin{itemize}
    \item 提案位置は低水位運用でも貯水池水理に律速されず，最大揚水流量を通せる．
    \item pumping（吸込）は強く非対称：片側に高流速，反対側に再循環域．構造物が貯水池軸に沿い流れが正面と片側から進入する立地に起因．generating（吐出）は平面的にほぼ対称だが鉛直に不均一で，開口上壁から剥離して下部を通過し，構造物直近の表面流は構造物へ向かって戻る．
    \item 最終形状で再循環域は半減（$572 \to 301$~$\mathrm{m}^2$）．ただし本文では流れは依然非対称に進入し，最大開口 No.~6 が全流量の 25~\% で design 1 と不変，損失水頭の低減もわずか．Abstract の「接近流がより対称・配分が均一化」は本文より強い表現なので，本文側を採る．
  \end{itemize}
  \item \textbf{限界・適用範囲}：site-specific．定常平均場のみで非定常性の議論はない．Froude 相似のみで，CFD と実験は再循環域の広がりが一致せず未決着．
  \item \textbf{本研究との接点}：「同一構造物でもモードで流況が異なる」実施設での傍証．ただし非対称の起源は立地由来で可逆性に内在せず，§3 Ferrara 2018 の $d_\mathrm{in}/d_\mathrm{out}$ とは位置を所与とする点で論点が異なる．2D 浅水で遠方場・3D で近傍場という役割分担（物理模型は rigid-lid 仮定の検証と起動時サージ・渦の確認を担う）は，本研究の 2D 表面 PIV と 3D 構造の関係（Valero 2022, Kantoush 2008a）に呼応．
\end{itemize}
\noindent（確認: p.483--495）


\noindent\litref{M\"uller et al. (2017)}\quad 可逆ポート＋矩形水槽で inflow/outflow 系列を直接扱った実験（M\"uller 2012 学位論文の実験部の論文版．本論文単体からは章番号を確認できないが，学位論文の目次・Ch7 で Ch6 ＝本論文と確定）．
\begin{itemize}
  \item \textbf{問い}：反復的な inflow/outflow 系列は，細粒土砂の沈降過程と系の土砂収支にどう影響するか．
  \item \textbf{手法}：2 槽（main basin $4.0 \times 2.0 \times 1.5$~m，水深 1.15~m ＋ mixing tank）を可逆水回路でつないだ閉鎖系．前壁に可逆ポート（ベルマウス，内径 48~mm，高さ $z_i/B = 0.125/0.25/0.375$），$Q = 0.3$--$1.1$~L/s，20 runs．クルミ殻粉末＋濁度計 2 台のみで評価（Froude 相似，$\lambda_L = 50$--$100$ 相当）．指標は SSR（懸濁土砂比），$\mathrm{INC}_\mathrm{SSR}$（静水沈降曲線基準の正規化増分），ISR/ESR，$\mathrm{SB} = \mathrm{ISR} - \mathrm{ESR}$．時間は $t_m = V/Q$ で正規化．なお IN 系列＝池へ流入＝本研究の inflow（Berm\'udez 2017 とは逆で，こちらは同じ向き）．UVP で流速場も測っているが本論文には非掲載で，流況は M\"uller 2018 / M\"uller 2012 を見ること．
  \item \textbf{主要知見}：
  \begin{itemize}
    \item SSR は inflow 系列で高く outflow 系列で低い．主因は流入ジェットが持ち込む運動エネルギーによる混合と結論（混合槽からの流入では差を説明できず，再懸濁は底面流速が低く寄与小と推定）．
    \item 流量で効き方が質的に変わる．低流量ではサイクルと相関せず静水比 10--40~\% 増，高流量では同期して振動し 50--80~\% 増．系列長 $/t_P < 1$（定常流を作らせず反転）ほど浮遊維持に有利．
    \item 取放水口高さは非単調：$z_i/B = 0.125$ と $0.375$ が $0.25$ より 20~\% 高い（鉛直再循環の増幅）．
    \item 最初の 1 サイクル（$t/t_m < 0.2$）で細粒の約 60~\% が沈降し，サイクル規模・頻度にほぼ依存しない．
    \item 土砂収支はほぼ平衡（$-0.03 < \mathrm{SB} < 0.04$）．往復させるが 2 槽間で再分配しない．規模・頻度は効かず，相対系列長のみが効く（SB が概ね正なのは最初の inflow 系列が最大の投入を生むため）．
  \end{itemize}
  \item \textbf{限界・適用範囲}：浅水槽ではない（$h/B \approx 0.58$，機構は鉛直再循環）で本研究とは別レジーム．濁度は点計測を全体の代表と仮定，再懸濁は未定量．Table~2 の $\mathrm{SSR}_{m,\mathrm{final}}$ は $K_{t_P}$ ごとに測定時刻が異なり生値比較不可．
  \item \textbf{本研究との接点}：「同一の可逆ポートでも inflow/outflow は非対称」を機構付きで示した最も近い実験．問い (1) の傍証．将来の土砂実験に有用：指標の定義，静水参照曲線の取り方，クルミ殻トレーサと濁度校正，$t_m$ 正規化．土砂に着手する際は本論文と Jenzer Althaus 2015 を最初に読み直す（§6）．流況の文脈（§4 表面 PIV）では引けない（UVP 非掲載）．
\end{itemize}
\noindent（確認: p.155--170）

\noindent\litref{M\"uller et al. (2010)}\quad（補助）Grimsel II 下池の現地 ADCP 観測（本論文単体からは対応関係を確認できないが，学位論文の目次・Ch7 で Ch3 ＝ M\"uller 2018 とその会議版＝本論文と確定．ページは Zotero の「1146」に対し本文は 1139--1145）．
\begin{itemize}
  \item \textbf{手法}：放流/取水構造物前面の湖底に ADCP 3 台を固定し，3D 流速プロファイルを運転データと比較（FFT でピーク周波数の一致を確認）．5 分平均をさらに 25 分移動平均．
  \item \textbf{主要知見}：generating では北・南軸で外向きにならず西側の急斜面に曲げられ，10 枚の案内壁による均等配分の設計意図が崩れる．Abstract の「両モードで期待方向の流れ場を観測」は本文より強い表現．主流は放水口軸のやや上・厚さ 5--10~m，底層と中間層は動かず，表層はジェット層と逆向き（風は未分離）．
  \item \textbf{限界・適用範囲}：水平分解能 0.12~m/s 程度＝「流速なし」は「測れず」の意．平均化のためキーワードに反し乱れは未取得．深い湖（最大 100~m），morning glory 型構造で浅水流ではない．
  \item \textbf{本研究との接点}：モードで流況が異なる現地証拠だが，非対称の起源は地形．平均場のみで問い (3) には使えない．核は M\"uller 2012 / 2018．
\end{itemize}
\noindent（確認: p.1139--1145）

\noindent\litref{Zhang et al. (2020)}\quad（補助）鉛直管型 inlet/outlet の形状最適化．Zhu 2023・Bai 2026 と同一の天津大グループ．
\begin{itemize}
  \item \textbf{問い}：鉛直管型流入/放流口の形状パラメータは水理性能（損失水頭・流速分布）をどう決めるか．双方向流の下で最適形状は．
  \item \textbf{手法}：中国 Xilongchi 揚水発電所の鉛直管型 inlet/outlet を baseline とする．CFD（ANSYS Fluent 17.0, RANS, Realizable $k$-$\varepsilon$ が最良）＋応答曲面法 RSM ＋遺伝的アルゴリズム NSGA-II の統合最適化．設計変数 4 つ：オリフィス高さ比 $H^*$，ディフューザ短半軸比 $a^*$，長半軸比 $b^*$，カバープレート半径比 $R_c^*$．OLHS で 80 サンプル点．GCI で格子収束を評価．ADV 実験で検証．なお $C_{V,\mathrm{in}}$ と $C_{V,\mathrm{out}}$ は分母が異なり直接比較できない（前者は鉛直管の平均流速 $v_d$，後者は当該測線の平均流速で正規化）．baseline の $0.136$ 対 $2.508$ は不均一度の比ではない．
  \item \textbf{主要知見}：
  \begin{itemize}
    \item 双方向流の本質的トレードオフ：発電モードとポンプモードの最適水理条件は同時に満たせない（``incompatible contradiction''，原文の表現）．総損失は $\xi_\mathrm{total} = \xi_\mathrm{in} + \xi_\mathrm{out}$ で評価せざるを得ない．
    \item 最適設計（$H^* = 0.422$, $a^* = 1.177$, $b^* = 5.363$, $R_c^* = 2.115$）は基準設計に対し $\xi_\mathrm{total}$ を 4.687~\%，$C_{V,\mathrm{in}}$ を 11.765~\%，$C_{V,\mathrm{out}}$ を 38.596~\% 低減．ただし $\xi_\mathrm{in}$ は 9.551~\% 増，$\xi_\mathrm{out}$ は 11.005~\% 減＝両モードのトレードオフが数値的に明示．
    \item 鉛直方向のモード非対称：発電モード（貯水池へ吐出）はピークが開口の下部（$h^* \approx 0.20$），ポンプモード（吸込）はピークが上部（$h^* = 0.86$）で下部 $h^* = 0$--$0.25$ に逆流域．→吐出時に開口下部へ流れが集中する点は Berm\'udez 2017・Guo 2024 と独立に一致し，3 研究で共通する頑健な知見．
    \item 特筆：ポンプモードで格子断面に生じていた負の流速が完全に消失．ただしディフューザ上部の壁近傍には逆流が残り，構造物全体から消えたわけではない．
    \item 支配パラメータは指標ごとに異なる（$\xi_\mathrm{in}$ は $a^*$，$\xi_\mathrm{out}$ は $b^*$，$C_{V,\mathrm{in}}$ は $R_c^*$，$C_{V,\mathrm{out}}$ は $H^*$ が最大）．単一パラメータが全目的を支配しないことが多目的最適化を要する理由．
    \item Pareto front 上で $\xi_\mathrm{total}$ は 3~\% 幅しか動かないのに $C_{V,\mathrm{out}}$ は 87~\% 幅動く＝損失をほとんど犠牲にせず均一性を大きく改善できる．
  \end{itemize}
  \item \textbf{限界・適用範囲}：構造物（鉛直管・ディフューザ）内部の最適化で，貯水池側の流況ではない（貯水池は境界条件を与えるための領域）．RANS の精度限界を著者自身が指摘．検証が最も弱いのが逆流領域で，まさに看板の主張が住む場所（著者は「逆流領域には差異があったが許容できると判断」と明記）．
  \item \textbf{本研究との接点}：「流入と流出で最適が両立しない」という双方向流固有の困難は，本研究の可逆ポートのモード非対称を工学的に裏づける論点．上記の鉛直非対称（吐出＝下部集中／吸込＝上部集中＋下部逆流）は本研究が表面流を見る際の解釈材料になりうる．ただし本研究は水平矩形断面ポートで形状は異なる．本論文は M\"uller 2010/2018 と Berm\'udez 2017 を引用しており，文献系譜の確認にもなる．
\end{itemize}
\noindent（確認: p.1499--1518）

\noindent\litref{Bai et al. (2026)}\quad（補助）側方 inlet/outlet の多スケール乱流を LES＋POD/DMD で解析．§4（POD/DMD）と §1（揚水発電）の橋渡し（Zhu et al.\ 2023 と同一グループ）．
\begin{itemize}
  \item \textbf{問い}：側方 inlet/outlet の局所損失（全系損失の 6.7--8.3~\%）を生む多スケール乱流機構は何か．
  \item \textbf{手法}：実施設寸法の LES（$\mathrm{Re} = 2.5\times10^7$，35 FTP，自由水面は $\mathrm{Fr} = 0.07$ として rigid lid）＋POD＋DMD．Zhu の LDV 実験で平均場のみ検証（MAE は付着流 4--8~\%，剥離域 11--14~\%）．なおメッシュ数は Abstract「4200 万」と本文「medium ＝ 3200 万を採用」が矛盾しており要確認．
  \item \textbf{主要知見}：
  \begin{itemize}
    \item 幾何拡大が作る中間遷移域（$x/D = 2.0$--$5.0$）で TKE 生成が入口基準の約 12 倍＝損失の主座．
    \item POD 第 1 モード（0.09~Hz）が全 TKE の 25~\%，上位 5 モードで 70~\% 超．剥離バブルの膨張/収縮＝大規模 breathing モードが大域的な剥離トポロジーを支配（自己相関でも $T = 11$~s を確認）．
    \item デュアルパスウェイ：大規模エネルギーの 43~\% が逐次カスケードを迂回して中間スケールへ直接注入され，高周波エネルギーの 65~\% が低周波構造との結合に由来（PSI ＝ 0.78）．近傍場の $-5/3$ は見かけで，2 次構造関数の指数 1.6（理論 2/3）が非局所性を暴く．
    \item 含意：小スケール乱流ではなく低周波の大域不安定（0.09~Hz）を抑えることが最適化の鍵．
  \end{itemize}
  \item \textbf{限界・適用範囲}：放流方向しか計算しておらず，ポンプ方向は未計算＝可逆性・モード非対称は扱っていない．構造物内部が対象で貯水池側の流況ではない（rigid lid のため表面流も出ない）．数値のみで，非局所伝達の定量値には直接的な実験検証がないと著者自認．単一 Re．
  \item \textbf{本研究との接点}：POD/DMD の実装作法の参照（DMD モード数を振幅閾値の感度解析で決定，自己相関で支配周波数を先に押さえる，PSD と空間構造の並置でモード番号＝スケールを確立）．浅水場への適用は Valero 2022．「時間平均では捉えられない低周波大域振動が支配的」は $I_\mathrm{unst}$ の解釈枠組みの候補．ただし 0.09~Hz は拡散部の剥離バブルの breathing であって basin の現象ではなく，類推止まりで機構は同一でない．
\end{itemize}
\noindent（確認: art.\ 141236）

  

\noindent\litref{Guo et al. (2024)}\quad（補助）側方 inlet/outlet の鉛直拡散角 $\alpha$ を LES で最適化．Bai 2026 の姉妹論文（同じ側方 inlet/outlet，同じ WALE-LES，同じ Zhu 実験で検証．差分は Guo ＝幾何パラメータ $\alpha$ の設計最適化，Bai ＝モード分解）．
\begin{itemize}
  \item \textbf{問い}：鉛直拡散角 $\alpha$ は剥離・乱れ強度・水頭損失をどう決めるか．
  \item \textbf{手法}：LES（WALE，3000 万セル，31 FTP）で $\alpha = 0/1.5/3/4.5/6/7.5^\circ$ の 6 水準．outflow 方向のみ（管路→貯水池）で可逆性は未計算．中央ピアで対称と仮定し半断面のみ＝開口間の配分の偏りなど非対称現象は構造的に出ない．検証は Zhu 実験だが寸法が別（$H = 7.2$ vs $9.0$~m，$\alpha = 2.35^\circ$）で同一形状の検証ではない（MAPE 8.21~\%）．
  \item \textbf{主要知見}：
  \begin{itemize}
    \item 逆流係数 $\chi$（式 6）＝統計期間中に流速が負だったサンプルの割合．$\chi = 1$ で常時逆流，0 で常時順流．定義が単純で time-resolved データにそのまま適用できる．本論文で最も独自性のある指標．
    \begin{itemize}
      \item $\chi = 0.5$ の等値線は平均流速 $U = 0$ の等値線とほぼ一致するが，平均場で剥離が消える $\alpha = 3^\circ$ でも $\chi$ には瞬時逆流が残り，完全消失は $\alpha = 0^\circ$（$\chi < 0.01$）＝間欠性は平均場より粘る．
      \item 高乱れは再循環の内部（$\chi \geq 0.5$，むしろ低乱れ）ではなく縁の帯 $0.01 < \chi < 0.5$ でピーク．コア流に直接さらされて運動量交換が激しいため．Q 基準でも再循環内部に渦はほぼない．
    \end{itemize}
    \item $\alpha$ 増で上壁から剥離しコア流が押し下げられ有効断面が減少．損失係数は $0.117 \to 0.138$（$+17.9$~\%）．ただしラック部の局所流速は $\alpha$ 低下に対し「まず減り後に増える」ため $\alpha = 3$--$5^\circ$ を推奨．
    \item トラッシュラック部の卓越周波数 0.09~Hz（Zhu 実験 0.1~Hz，Bai の POD Mode 1 と一致）．固有振動数 15.57~Hz と離れ共振せずとするが，ラック自体は未モデル化で間接推論（著者も FSI が必要と明記）．
    \item ヘアピン渦パケットの成長角が乱れ強度の伝播方向と一致．下部 9.1--11.2$^\circ$（$\alpha$ 非依存，平板境界層の 10.5$^\circ$ に近い）に対し上部は $16.7 \to 11.2^\circ$ と $\alpha$ に強く依存．
  \end{itemize}
  \item \textbf{限界・適用範囲}：構造物内部の管路・拡散部が対象で貯水池側の流況ではない．大規模コヒーレント構造は扱わず今後の課題（そこは Bai 2026 が担当）．$\alpha$ 以外は固定で交互作用は未検討．
  \item \textbf{本研究との接点}：$\chi$ は本研究の secondary metric の候補．$I_\mathrm{unst}$ が $E(t)$ の変動の「大きさ」を測るのに対し，$\chi$ は「符号反転の頻度」を測る別種の量．outflow の非定常性の記述に使える可能性．ただし $\chi$ は方向が定義された流れ（拡散角で射影）が前提なので，basin の 2D 表面流への定義し直しは要検討．「平均場で消えても間欠的逆流は残る」は問い (3) と同型の主張．ただし対象は拡散部で basin ではない．
\end{itemize}
\noindent（確認: art.\ 114291）

\noindent\litref{Constantinescu \& Patel (1998)}\quad（補助）ポンプ取水槽まわりの 3D 流れと渦を数値的に予測できるかを検証．
\begin{itemize}
  \item \textbf{問い}：3D 数値モデルはポンプ取水（pump-intake）まわりの流れと渦を予測できるか（手法の feasibility）．
  \item \textbf{手法}：3D RANS ＋2 層 $k$-$\varepsilon$（Chen \& Patel 1988，$\varepsilon$ の壁境界条件が不要）＋完全陰解法 ADI．矩形水路から円管 1 本が吸い込む pump sump（Melville et al.\ 1994 準拠：$S/D = 2.0$, $X/D = 0.4$）．55 万格子，$\mathrm{Re} = 60{,}000$，自由表面は平坦・対称面．なお寸法は強い自由表面渦を避けるよう選択されている．
  \item \textbf{主要知見}：
  \begin{itemize}
    \item 3 種の渦を同定：V1（角部），V2（後壁と管の間の自由表面渦），V3（管外壁近傍で最強＝外壁境界層に起因）．V2 が V1 並みに弱いのは上記の条件設定の帰結で一般的性質ではない（著者も留保）．
    \item $\mathrm{Re} > 30{,}000$ で自由表面渦への粘性効果は無視できるという通説を，非次元渦度が概ね支持．
    \item 底面の摩擦線（Fig.~8d）＝壁面せん断応力ベクトル場（限界流線）．掃流砂は $\tau_w$ に駆動され，動く向きは上空ではなく摩擦線が決める．3D 流れでは底面流向が表層と食い違うため（河川湾曲部の二次流で砂が内岸に溜まるのと同型），この図から堆積・洗掘の地図が読める：摩擦線の収束線＝底面流体が集まり上昇→堆積，発散線・付着点＝せん断が高く掃流→洗掘．ただし本論文の記述は 1 文のみで展開はなく，向きの情報しか持たない（$|\tau_w|$ と限界掃流力の比較は別途要）．
    \item 著者の留保：理想化対称条件で予測される対称渦は実際にはめったに観測されない．実験検証もない．
  \end{itemize}
  \item \textbf{限界・適用範囲}：吸込側のみで可逆運転ではない．定常 RANS（著者は「非定常特徴を確実に予測する能力は現時点で存在しない」と明記）．rigid lid（ただし圧力場から水面変形 0.56~mm を事後推定し仮定を自己検証）．
  \item \textbf{本研究との接点}：
  \begin{itemize}
    \item 「対称条件を課した数値解析は対称解しか返さない」という留保は，本研究のモード非対称の論点と同型（§3 Dewals, Kantoush 2008a）．Guo 2024 が中央ピアで対称を仮定していることへの警告としても読める．
    \item 表面 PIV への含意：表面の streamline は底面の摩擦線ではない．「低流速域 $\phi_\mathrm{lv}$ に土砂が溜まる」という推論は $h/W$ が小さく流れが 2D 的なほど成り立ち，鉛直循環が効くほど崩れる（M\"uller 2017 の取放水口高さ $z_i/B$ の非単調性がその裏側）．port 近傍は 3D 性が強く $\Omega$ から除外済み．土砂実験に着手する際，表面の指標と底面の輸送をどう繋ぐかを明示的に扱う必要がある（→§6）．
    \item Zhu 2023 との呼応は要確認：本論文は吸込（basin 側から見れば outflow）で，Zhu の「流出」が構造物から貯水池への吐出を指すならモードが逆になる．
  \end{itemize}
\end{itemize}
\noindent（確認: p.123--134）

\noindent\litref{Imberger et al. (1976)}\quad（補助）成層した有限矩形タンクの線シンク流．主題（成層取水層）は本研究に使えず，価値は\textbf{極限の記述}数箇所に限られる．\texttt{discussion.tex} の「流出の吸込みはポテンシャル流に近く渦度供給源をもたない」を支える文献として引く．
\begin{itemize}
  \item \textbf{問い}：成層貯水池で線シンクを始動すると取水層はどう形成されるか．後壁は何を変えるか．
  \item \textbf{手法}：理論（Boussinesq ＋ Laplace 変換）＋particle-in-cell 数値解．2 次元\textbf{鉛直断面}，線シンク．支配パラメータ $R = F\,Gr^{1/4}$ ほか全てが成層強度 $N$ を含む．TVA 実機（$L \sim 16$~km）と照合．知見の本体は，始動→内部せん断フロント列→後壁で反射→定在波→減衰後に取水層．分割流線が後壁に付着せざるを得ないため\textbf{定在波がなくても厳密な定常解は存在しない}．
  \item \textbf{転用できる箇所（実質的な収穫）}：
  \begin{itemize}
    \item p.490／式 (10)：シンク始動直後は\textbf{ほぼ水平なポテンシャル流}が確立し，水平流速は端壁に向かって線形にゼロへ減少する．＝有限矩形タンク内のシンク流の初期応答はポテンシャル流．
    \item §4 p.496：「重力効果は取水層を生むのに十分でなければならず慣性力に完全に支配されてはならない．\textbf{さもなくば単純なポテンシャル流にしかならない}」＝浮力が効かない極限＝ポテンシャル流と著者が明示．本研究は $N = 0$ でまさにこの極限．
    \item p.494：\textbf{分布シンク＋緩やかな始動}の両方（式 13, 14 の $a$ と $b$）が満たされると「せん断フロントも定在波も現れず，運動は常に準定常で初期のポテンシャル流に等しい」．本研究の port は幅 0.11~m の有限開口だが，条件は分布シンク単独では足りない点に注意．
    \item §7 p.506：流入水が中立浮力で\textbf{運動量が小さければ，流入と流出のスケーリングは厳密に対称}（流線関数の符号が反転するだけ．違いは slug の有無）．＝非対称が生じるならその原因は運動量．M\"uller 2017 の「主因は流入ジェットの運動エネルギー投入」と理論側から整合し，主張 (a) を裏返しに支える．ただし著者は「跳水・巻き込みを伴う潜り込み流は層流理論の範囲外」と留保．
  \end{itemize}
  \item \textbf{限界・適用範囲}：\textbf{成層が主題}でパラメータが全て $N$ を含み，均質流体では取水層という現象自体が存在しない．\textbf{断面が鉛直面}で本研究の水平面とは向きが異なり，平面内でシンクが側壁とどう相互作用するかは何も言えない．壁は slip 条件で底面摩擦を扱わない．
  \item \textbf{本研究との接点}：
  \begin{itemize}
    \item \textbf{引き方}：「成層取水層を扱った研究だが，浮力が慣性に対して無視できる極限では，有限矩形タンク内の流れは初期ポテンシャル流のまま準定常にとどまる」と極限の記述として引く．「Imberger らが吸込み流を扱った」と一般化すると成層が主題である事実を隠す．
    \item \textbf{今後}：Imberger 1980 のレビューも同じ成層の枠組みで均質流体には届かない．主張を厚くするなら教科書（Batchelor, Lamb）のポテンシャルシンクの方が素直．残る穴は Fischer et al.\ 1979（未確認）．
  \end{itemize}
\end{itemize}
\noindent（確認: p.489--512，全 24~p 通読）

\noindent\litref{M\"oller et al. (2015)}\quad（周辺）取水口渦による空気連行の定量（VAW/ETH Zurich）．
\begin{itemize}
  \item \textbf{内容}：水平・鋭縁の取水管（$D = 389$~mm）で空気連行率 $\beta = Q_a/Q_w$ を脱気装置により準連続計測（34 runs）．$\beta$ は合成 Froude 数 $\mathsf{F}_{co} = v_D/\sqrt{gh}$ の指数関数（$R^2 \geq 0.9$）．従来は渦の目視で決めていた限界潜没を，$\beta = 10^{-5}$ を inception と定義して逆算し直した（$(h/D)_{cr} = -2.5\,\mathsf{F}_D^{-0.45} + 5.3$，$h$ は管軸から）．
  \item \textbf{本研究に効くのは 1 点＝スケール効果}：DN200/300/400/500 の系列試験で，DN500 基準の連行率比の中央値が DN200 で 0.10，DN300 で 0.33，DN400 で 1.05．\textbf{$D < 400$~mm では Froude 相似が破れ連行が 1 桁小さく出る}．無次元の閾値は $\mathsf{R}_D \geq 6\times10^5$ かつ $\mathsf{W}_D > 3.2\times10^3$（古典値は $\mathsf{R}_{D,cr} = 3.2\times10^4$，$\mathsf{W}_{D,cr} = 121$）．著者自身「プロトタイプへの外挿は 4 点の中央値に基づくだけで厳密には illusory」と留保．
  \item \textbf{本研究との接点}（→ Limitations に 1--2 行）：
  \begin{itemize}
    \item 本研究の $U_p \approx 0.12$~m/s は Kobus 1984 の空気連行開始速度 0.8--1~m/s の 1/7 以下で\textbf{空気連行は起きない}．本論文の主題（$\beta$ の定量，限界潜没）は適用外．空気連行の有無を問われたらこの見積もりで足りる．
    \item スケール効果は，自由水面の変形や渦の巻き込みに関わる現象を実機へ外挿できないことの具体例として引ける（本研究の概算 $\mathsf{R}_D \approx 7\times10^3$，$\mathsf{W}_D \approx 12$ は閾値を 2 桁下回る．要検算）．
    \item \textbf{$h/D$ と $h/a$ は同型ではない}：$h/D$ は空気連行渦を「避ける」ための設計基準，$h/a$ は「調べる」対象．加えて本研究の $h/a = 0.55$--$0.91$ は port が完全水没していない領域で，M\"oller の範囲と重ならない．
  \end{itemize}
\end{itemize}
\noindent（確認: art.\ 04015026）

\subsection*{論文横断整理}

\noindent\textbf{(1) この節の全文献が指す一点：流入と流出は「別の流れ」である}

同じポートを使っていても，水を送り込むときと吸い込むときでは流れの性質が根本的に違う．理由は単純で，

\begin{itemize}
  \item \textbf{流入はジェット}である．運動量を持って池に入り，その運動量で遠くまで進み，まわりの水を巻き込んで大きな循環をつくる．
  \item \textbf{流出は吸い込み口}である．運動量を持ち込まないので，ポート近傍だけが速く，離れると急速に弱まる．
\end{itemize}

\noindent ジェットは「構造をつくる」，吸い込みは「構造をつくらない」．この非対称は幾何形状や規模に依存しない一般的な性質であり，だからこそ本節の文献は，スケールも手法も違うのにそろって同じことを言っている．

\medskip
\noindent\textbf{(2) 数 cm から数百 m まで，同じ非対称が現れる}

\begin{itemize}
  \item \textbf{構造物の中}（Zhu 2023, Zhang 2020）：流出時は流速が激しく変動し逆流域も生じるが，流入時は穏やか．
  \item \textbf{室内水槽}（M\"uller 2012）：流入はジェット＋2 つの循環セル．流出はポート近傍だけが動き，後方は準静止．
  \item \textbf{実貯水池}（M\"uller 2018）：吐出は 150~m 先まで届き大規模循環セルをつくるが，吸込は取水口ごく近傍にとどまる．
  \item \textbf{実施設の設計}（Berm\'udez 2017）：吸込は片側に偏り，吐出は平面的には対称でも鉛直には偏る．
\end{itemize}

\noindent 格子断面という数 cm のスケールから水深 100~m の湖まで，同じ非対称が観察されている．本研究が「流入と流出でモードが変わる」と予想することには，この意味で十分な根拠がある．

\medskip
\noindent\textbf{(3) ただし「流出は局所的」は本研究には当てはまらないかもしれない ― ここが独自性}

M\"uller の室内実験も現地観測も，「流出時は取水口の近くだけが動き，遠方は止まっている」と報告する．ところが本研究の予備観察では，流出時にも池全体に単一の循環が生じている．

食い違いの原因として最も疑わしいのは\textbf{水深}である．M\"uller の水槽は $h/B \approx 0.58$，本研究は $h/W \approx 0.02$--$0.05$ で 1 桁以上違う．深い系では吸い込み流は 3 次元的に広がれるので距離とともに急速に弱まる．しかし極端に浅く側壁で囲まれた系では鉛直方向に逃げ場がなく，流れは 2 次元的に広がるしかない．結果として吸い込みの影響が池全体に及びうる．

つまり「流出は局所的」という既往の描像は，\textbf{深い系に限った話である可能性がある}．浅い系で流出モードを面的に測った研究は見当たらず，ここが本研究の主張しうる新規性である．

\medskip
\noindent\textbf{(4) 文献は「構造物の中を見る群」と「池を見る群」に分かれ，両者は繋がっていない}

\begin{itemize}
  \item \textbf{構造物の中を見る群}（Zhu 2023, Zhang 2020, Bai 2026, Guo 2024, Constantinescu 1998）：ポートやディフューザの内部を CFD・模型で詳細に解き，損失水頭・流速の均一性・渦・振動を扱う．貯水池は境界条件を与えるための領域にすぎず，池の流況は出てこない．
  \item \textbf{池を見る群}（M\"uller 2010/2012/2017/2018, Jenzer Althaus 2015, Berm\'udez 2017）：池の中の流れや土砂を扱う．ただし目的は排砂・土砂管理であって流況そのものの分類ではなく，水深も深い．計測も UVP・ADCP の点／断面計測である．
\end{itemize}

\noindent したがって\textbf{「浅い矩形池の表面流況レジームを，可逆ポートの流入・流出で比較分類する」という位置は空いている}．本研究はこの隙間に入る．なお 3 研究（Zhang 2020, Berm\'udez 2017, Guo 2024）が独立に「吐出時は開口の下部へ流れが集中する」と報告しており，これは頑健な知見として扱ってよい．

\medskip
\noindent\textbf{(5) 本研究の設計・解釈に直接効く教訓}

\begin{itemize}
  \item \textbf{ヒステリシス}（M\"uller 2012）：流出時の流況は直前の流入に影響され，大きな循環セルは散逸するまで残る．→ 各 run を静水から始め，残留がないことを $E(t)$ で確認する手続きが要る．
  \item \textbf{対称条件は対称解しか返さない}（Constantinescu 1998）：理想化した対称条件を課した数値解析は対称解を返すが，現実にはそのような対称流はめったに観測されない．→ モード非対称・双安定を論じる際は，数値の対称仮定が結論を先取りしていないか疑う．Guo 2024 が中央ピアで対称を仮定し半断面だけ計算しているのは，まさにこの落とし穴に当たる．
  \item \textbf{表面の流線は底面の流線ではない}（Constantinescu 1998）：土砂を動かすのは底面のせん断応力で，その向きは表層と食い違いうる．→ 表面 PIV から堆積を推論するときは，流れが 2 次元的であるほど安全，鉛直循環が効くほど危うい．
  \item \textbf{見るべきは小スケールでなく低周波の大規模振動}（Bai 2026, Guo 2024）：エネルギーと損失を支配していたのは低周波の大規模「呼吸」モードだった．また平均場で剥離が消えても瞬時場には逆流が残る（Guo の $\chi$）．→ 時間平均だけを見ると本質を取り逃す．$I_\mathrm{unst}$ のような非定常指標が必要な理由．
\end{itemize}

\medskip
\noindent\textbf{(6) 読み方の注意：この分野では Abstract が本文より強い}

照合作業で繰り返し現れたので記録しておく．Berm\'udez 2017 の「接近流が対称化した」，M\"uller 2010 の「両モードで期待どおりの流れ場を観測」は，いずれも本文を読むと限定がつく．Bai 2026 はメッシュ数が Abstract と本文で食い違う．Zhang 2020 は看板の主張（ポンプモードの逆流消失）が住む逆流領域こそ検証が最も弱い．\textbf{引用の際は必ず本文側の記述を採ること．}

\subsection*{原稿転記メモ}

% ================================================================
\section{浅い池・貯水池の流動パターン}
% ================================================================

\subsection*{主要な問い}

\subsection*{論文メモ}

\noindent\litref{Dufresne et al. (2010)}\quad 矩形浅水池の流況分類と再付着長の支配パラメータを実験で体系化．
\begin{itemize}
  \item \textbf{問い}：矩形浅水池の流況は形状・水理条件を横断してどう分類でき，再付着長の中央値と自然変動を何が支配するか．
  \item \textbf{手法}：Liège 大学 10.40~m 水路（幅 0.985~m，入口区間長 2.0~m，池長 $L$ は可変），染料注入可視化，40 条件．
  \item \textbf{主要知見}：
  \begin{itemize}
    \item 流況 4 種：S0（対称・再付着なし），A1（再付着 1 点），A2（再付着 2 点），A3（再付着 2 点 + 離脱 1 点）．遷移域で非周期変動の不安定 A1/S0．
    \item 遷移基準：形状パラメータ $\mathrm{SP} = L/(\Delta B^{0.60} b^{0.40})$．$\mathrm{SP} < 6.2$ 対称，$6.2$--$6.8$ 不安定，$\mathrm{SP} > 6.8$ 非対称（高 $h/\Delta B$，$\mathrm{F} \approx 0.20$）．
    \item 第 1 再付着長 $R_1 \approx 3.43\,\Delta B^{0.75} b^{0.25}$（長い池，$\mathrm{F} \approx 0.20$），変動小（$\sigma_1/\Delta B \approx 0.29$）．
    \item 第 2 再付着長 $R_2 \approx 15.9\,\Delta B^{1.7} b^{-0.7}$，変動大（$\sigma_2/\Delta B \approx 6.0$）で A2 は本質的に不安定．
    \item $h/\Delta B > 0.8$ で深さ非感受，$\mathrm{F}$ の影響は $R_1$ の平行移動程度．
  \end{itemize}
  \item \textbf{限界・適用範囲}：一方向流（流入のみ）に限定，可逆単一ポートの流向切り替えには非対応．
  \item \textbf{本研究との接点}：流況分類（S0, A1, A2）と SP 遷移基準を本研究の遷移指標と比較する基本枠組み．本研究の池形状がどの SP 領域に属するかは別途確認を要す．
\end{itemize}
\noindent（確認: p.197--203）

\noindent\litref{Dufresne et al. (2011)}\quad 2D 深さ平均浅水モデル（Wolf 2D）の対称・非対称流況再現能力を検証した数値研究．
\begin{itemize}
  \item \textbf{問い}：2D 深さ平均浅水モデル（Wolf 2D）は対称・非対称流況とその遷移をどの程度再現できるか．
  \item \textbf{手法}：FVM Wolf 2D，2 スケール $k$-$\varepsilon$，Colebrook 摩擦（較正不要，ガラス壁 $k_s = 0$），格子 10~mm（検証用 25~mm），GCI 誤差評価，2 プロトコル（P1：非対称 IC+BC →無擾乱化，P2：二分法で臨界長を追跡）．
  \item \textbf{主要知見}：
  \begin{itemize}
    \item 誤差は GCI で $R_1$ が 5--10~\%，$R_2$ が 2~\% 未満．
    \item 遷移基準は概ね再現．P1 は過大評価（$\Delta B/b = 4.38$ で 7.5--7.7 vs 実験 6.2--6.8，小 $\Delta B/b$ ほど顕著），P2 は臨界 $\approx 7.0$ と実験上限に近い．
    \item $R_1$ 再現精度は $\Delta B/b = 4.38$ で $< 5\,\%$，$0.52$ で最大 30~\% 誤差．
    \item 予測（$\Delta B/b = 5$--$10$）でも臨界形状パラメータ $\approx 7.0$，$R_1/\Delta B$ 漸近値 $\approx 2$．
    \item 小 $\Delta B/b$ の不一致は，再付着長変動の再現不可と流路曲率による底面摩擦の過小評価に起因．
  \end{itemize}
  \item \textbf{限界・適用範囲}：可逆ポートの流向切り替えを含まず，深さ平均仮定が成立しない近ポート 3D 流には非対応．
  \item \textbf{本研究との接点}：将来 CFD を用いる際の方法論的先例．較正不要の Colebrook + $k$-$\varepsilon$ が浅水池で機能．初期条件依存の議論は PIV 計測の初期流況依存評価の根拠にもなる．
\end{itemize}
\noindent（確認: p.247--256）

\noindent\litref{Peltier et al. (2014)}\quad 浅水矩形池の蛇行噴流の発生条件と特性を LSPIV+POD で定量化．
\begin{itemize}
  \item \textbf{問い}：浅水矩形池の蛇行噴流はどの条件で発生し，周波数・波長・横方向広がりを何が支配するか．
  \item \textbf{手法}：Liège 同一施設（$B = 0.985$~m，$b = 0.08$~m（一部 $0.06$~m），池長 0.7--1.6~m），表面 LSPIV（$1 \times 1$~m，25~Hz，直径 2~mm トレーサ），POD．80 条件中 50 条件で蛇行（うち 39 を POD 解析）．
  \item \textbf{主要知見}：
  \begin{itemize}
    \item 蛇行発生は $\mathrm{F} > 0.21$ かつ $\mathrm{SF} < 6.2$（$\mathrm{SF} = L/(\Delta B^{0.6} b^{0.4})$）．$\mathrm{SF}$・$\mathrm{F}$ は有無を判別するが蛇行特性とは無相関．
    \item 蛇行特性は摩擦数 $S = \lambda \Delta B/8H$ が支配：$\Lambda_x/H = 145\,S + 2.4$（$R^2 = 0.85$），$\Lambda_y/H = 20\,S^{0.45}$（$R^2 = 0.62$），$\mathrm{St} = 0.004\,S^{-0.776}$（$R^2 = 0.66$）．
    \item 摩擦数増（水深浅）で縦拘束により蛇行弱化．POD 第 1 対モードが蛇行 coherent structure を代表．
    \item LSPIV の妥当性は Kantoush et al.\ の先行研究に依拠．
  \end{itemize}
  \item \textbf{限界・適用範囲}：可逆単一ポートには非対応．池幅 $B \approx 1$~m は本研究の $B = 2.0$~m と異なりスケーリング外挿に注意．
  \item \textbf{本研究との接点}：LSPIV+POD の非定常解析手法が適用可．「対称流」は一様でなく，安定対称流（S0）と対称的非定常蛇行流は別物で，この区別を欠くと短池の非定常な混合・堆積拡散を見落とす．本研究の流入流況が $\mathrm{F} > 0.21$ の蛇行域に該当するかは摩擦数 $S$ の試算で確認を要す．
\end{itemize}
\noindent（確認: p.699--710）

\noindent\litref{Chagdali et al. (2025)}\quad 入口条件（開水路 vs 加圧円形管噴流）と噴流鉛直位置が流況に与える影響を実験で精査．
\begin{itemize}
  \item \textbf{問い}：入口条件（自由水面水路 vs 加圧円形管噴流）と噴流の鉛直位置 $Z_\mathrm{jet}/H$ は浅水矩形池の流況にどう影響するか．（背景：実施設は管路流入が多いのに研究が乏しい）
  \item \textbf{手法}：Liège 同一施設，2 池長・3 噴流高さ・3 流量の 22 条件，表面 LSPIV ＋断面 ADVP（Nortek Vectrino）．
  \item \textbf{主要知見}：
  \begin{itemize}
    \item 開水路入口は既往どおり $\mathrm{SF}$・$\mathrm{Fr}$ で S0・蛇行 O・A1 を概ね説明．
    \item 円形管噴流入口では入口の鉛直位置 $Z_\mathrm{jet}/H$ が効く．
    \item 底面近傍噴流は短池・長池とも不安定流 U を誘発しやすい．
    \item 中深・水面近傍噴流は多くの場合定常再付着流．ただし長池・低 $\mathrm{Fr}$ では噴流高さによらず不安定流．
    \item 加圧噴流条件で近入口域の 3D 構造が顕著，表面速度（LSPIV）は底面付近より低くなりうる．
    \item 短池で再付着点が下流側の A1\* は Dufresne (2010) になかった新型．
  \end{itemize}
  \item \textbf{限界・適用範囲}：出口は常に自由水面水路で，可逆ポートの流向切り替えは扱わない．
  \item \textbf{本研究との接点}：開水路入口で構築した既往分類を管路流入の実施設へ直接適用するのは危険．本研究のポートは水平矩形断面だが「底面付近の加圧噴流」の性格は共通し，流入時 U 型不安定の参照例．LSPIV+ADVP の併用は本研究に最も近い．
\end{itemize}
\noindent（確認: p.209--221）

\noindent\litref{Kantoush et al. (2008b)}\quad（補助）Kantoush 2008a（§4）と同一施設・枠組みの補完．池長依存の双安定の追加証拠．
\begin{itemize}
  \item \textbf{手法}：EPFL $6 \times 4$~m 基準池，入口/出口 0.25~m，$Q=7$~L/s，$h=0.2$~m．LSPIV/UVP/WOLF 2D．池長 6/5/4/3~m（AR=1.5/1.25/1.0/0.75）で流況を比較．
  \item \textbf{主要知見}：
  \begin{itemize}
    \item 対称設定でも長池（$L=6$~m, AR=1.5）で非対称流（Coanda 偏向，左右どちらにも倒れる双安定）．池長短縮で対称化（4 渦→2 渦）．
    \item WOLF 2D は完全対称入力→完全対称解（数値アーティファクトなし）．微小擾乱（$q_\mathrm{in}=q_0+q_1\cdot 2y/b$, $q_1/q_0=1$--$5\,\%$）で非対称解を再現＝対称解の不安定性を示す．擾乱強度 1--5~\% に結果は鈍感．
  \end{itemize}
  \item \textbf{本研究との接点}：Dufresne/Camnasio/Dewals と同系列の「池長依存の双安定」追加証拠．新規性は低く補助扱い（核は §2 Dufresne・§3 Dewals・§4 Kantoush 2008a）．
\end{itemize}
\noindent（確認: 会議論文・全体）

\subsection*{論文横断整理}

\subsection*{原稿転記メモ}

% ================================================================
\section{可逆な流入・流出とモード非対称性}
% ================================================================

\subsection*{主要な問い}

\subsection*{論文メモ}

\noindent\litref{Dewals et al. (2008)}\quad 対称設定でも非対称流（双安定）が生じる機構を実験・2D 数値・線形安定解析で解明．
\begin{itemize}
  \item \textbf{問い}：幾何・流入・流出が対称でもなぜ非対称流が生じ，いつ対称流が不安定化するか．
  \item \textbf{手法}：実験（$6 \times 4$~m 基準，8 形状，入口・出口幅 0.25~m，$Q = 7$~L/s，$h_0 = 0.2$~m，$\mathrm{Fr} = 0.10$）＋2D 数値（WOLF 2D，深さ平均）＋線形安定解析．
  \item \textbf{主要知見}：
  \begin{itemize}
    \item 対称設定でも test 1, 5, 6, 7 が非対称（2, 3, 4, 8 は対称）．微小擾乱の成長による鏡像 2 安定解＝双安定．
    \item 偏向は Coanda 効果（片側流速増→低圧→偏向増幅，遠心力と釣合いで定常非対称）で説明．
    \item 数値：完全対称入力→対称解．微小擾乱（$q_1/q_0 = 2\,\%$）→非対称解で実験一致（再付着 2.53~m vs 2.65~m，$< 5\,\%$）．対称解の外乱高感受性＝不安定を示す．
    \item 非対称度を中心線まわりの軸方向速度モーメント $m$，長手平均 $M = \mathrm{mean}|m|$ で定量化（$M$ は擾乱強度 1--5~\% に非依存）．
    \item 分岐は $R = \mathrm{ER} \cdot \mathrm{AR}^{8/9}$（$\mathrm{ER} = B/b$，$\mathrm{AR} = L/B$），臨界 $R_c$ で対称/非対称が分かれる（限界安定曲線 $\mathrm{ER} \approx (0.035\,\mathrm{AR})^{-8/9}$）．
    \item 水深を浅くすると底面摩擦で偏向弱化も試験範囲（$h_0 = 0.075$--$0.200$~m）で非対称残存．線形安定解析では test 8 の擾乱減衰が test 1 の 3 倍以上強く，最大増幅率約 1/20 で，はるかに安定．
  \end{itemize}
  \item \textbf{限界・適用範囲}：一方向流，入口・出口は対称配置で可逆単一ポートではない．
  \item \textbf{本研究との接点}：「対称設定でもモード非対称（双安定）が生じる」核心メカニズムの基盤．非対称度モーメント $m$ は本研究の $I_\mathrm{asym}(t)$ と対応づけ可能な候補指標．Coanda・外乱感受性・分岐の枠組みは可逆ポートのモード非対称議論の土台．
\end{itemize}
\noindent（確認: p.31--54）

\noindent\litref{Camnasio et al. (2011)}\quad UVP で浅水池の全面速度場を測り Dufresne 分類を独立施設で精緻化．
\begin{itemize}
  \item \textbf{問い}：$L/B$ と $B/b$ は流況型をどう決めるか．UVP の全面速度場で分類はどう精緻化されるか．
  \item \textbf{手法}：EPFL 矩形浅水池（最大 $6 \times 4$~m，$h = 0.20$~m 固定，$Q = 7$~L/s，$F_{in} = 0.10$，$R_{in} = 112{,}000$，計測高さ $z = 0.08$~m $= 0.4h$ を平均代表），UVP，41 試験．
  \item \textbf{主要知見}：
  \begin{itemize}
    \item 5 主分類：CH-L（ほぼ 1 次元水路流），S0（各側 1 渦），S1（各側 2 渦，上流側は stagnation zone），A1，A2．加えて遷移型 A1/S1．
    \item 遷移は形状パラメータ $T = L/[(B-b)^{0.6} b^{0.4}]$ で整理，$4.09 < T < 4.48$ で双安定的・不安定的挙動．
    \item 固定幅では池長短縮で $\mathrm{A1} \to \mathrm{S1} \to \mathrm{S0}$ と対称化．
    \item 固定長では側壁近接でまず対称流が不安定化し A1 的非対称，さらに極端に狭めると大規模循環が抑制され A2 的下流チャネル化を経て CH-L へ近づく．
    \item 非対称流（A1/A2）で主ジェットが片側壁へ偏向・付着する機構として Coanda 効果．
  \end{itemize}
  \item \textbf{限界・適用範囲}：一方向流・対称な入口/出口配置で可逆単一ポートではない．$T$ は Dufresne（$\Delta B = (B-b)/2$）と定義が異なり数値閾値を直接比較しない．
  \item \textbf{本研究との接点}：Dufresne (2010) の分類を独立施設の全面速度場（UVP）で裏付け．不安定遷移帯の双安定挙動はモード非対称の核．
\end{itemize}
\noindent（確認: p.352--358）

\noindent\litref{Miozzi \& Romano (2020)}\quad 自由水面浅水噴流の蛇行（横方向不安定）発生機構を高解像 PTV で精査．
\begin{itemize}
  \item \textbf{問い}：自由水面の浅水噴流で横方向の局所不安定はどう発生・伝播し，全体の蛇行運動を生むか．Fr の影響は．
  \item \textbf{手法}：大型水槽（$2 \times 2$~m）への自由噴流（$W = 4.2$~cm，噴流/水槽 $\approx 1/50$），時間分解 PTV（Optical Flow），$\mathrm{Re} = 5000$ 固定，$\mathrm{Fr} = 0.15$--$0.6$，$h/W = 0.5$--$1.1$．
  \item \textbf{主要知見}：
  \begin{itemize}
    \item 局所横方向不安定が増幅され全体蛇行（軸直交振動）へ発展．時間平均は壁噴流似・自己相似で両側に対称な再循環域（平均は左右釣合い）．
    \item 蛇行は自己維持機構：両側で反回転渦がジグザグに並び，建設的干渉で横変動を増幅．軸・横変動が強結合し，擾乱は $\approx 0.8\,U_0$ で伝播．
    \item ジグザグ波長 $\lambda_x \approx 3.5$--$4$（$W$ 単位，Fr 非依存）．蛇行振幅は下流で増加・Fr で減少，周波数（St）は逆（Fr で増加・下流で減少）．
    \item Fr 増で噴流配置が下流へ移動．渦同定は Okubo--Weiss．自己維持機構はアスペクト比 $\approx 1$（低 Fr）で活発．
  \end{itemize}
  \item \textbf{限界・適用範囲}：閉じ込め貯水池ではなく大型水槽への自由噴流（再付着なし）．可逆ポートなし．Coanda 壁付着でなく自由噴流の蛇行が主題．
  \item \textbf{本研究との接点}：「対称的だが時間的に蛇行する非定常流」の発生機構を最も精緻に示す．安定対称流と蛇行流を区別すべきという論点（Peltier）の物理的裏付け．高解像 PTV + Okubo--Weiss + 伝播速度解析は非定常解析の手法参照．Peltier 2014 と相補的．
\end{itemize}
\noindent（確認: p.1--18）

\noindent\litref{Ferrara et al. (2018)}\quad 入口・出口チャネルの位置を系統的に振り，再付着と双安定への影響を数値的に解明（本研究の可逆ポートに最も近い既往研究）．
\begin{itemize}
  \item \textbf{問い}：入口・出口を中心線からずらした位置（$d_\mathrm{in}$, $d_\mathrm{out}$）は，矩形浅水池の流況（再付着の有無・側）と双安定にどう影響するか．
  \item \textbf{手法}：検証済み 2D 数値 WOLF 2D（2 スケール $k$-$\varepsilon$）．$B = 4$~m 固定，$L = 4.5/5.5/6$~m，$b = 0.25$~m，$h = 0.2$~m，$\mathrm{F} \approx 0.1$，$Q = 7$~L/s（本文 §2.3 の $0.7$~L/s 表記は誤記の可能性．$b$, $h$, $\mathrm{F} \approx 0.1$ から逆算すると $7$~L/s で，Camnasio 2011 と同条件）．入口位置 11 × 出口位置 → 約 280 計算．静水 IC と再付着流 IC の 2 段階手順で双安定を判定．
  \item \textbf{主要知見}：
  \begin{itemize}
    \item 3 種の流況：片側壁へ再付着／再付着なし（入口→出口へ直行）／反対側へ再付着．
    \item 入口・出口が同じ側 → 両者がほぼ同等に効き，閾値は楕円 $(\alpha\,d_\mathrm{in}/B)^2 + (\beta\,d_\mathrm{out}/B)^2 = 1$（$\alpha$, $\beta$ は $L$ ごとに別値）．
    \item 入口・出口が反対側 → 主に入口位置が支配し，$|d_\mathrm{in}| \gtrsim B/4$（長池）〜$3B/8$（短池）で入口に近い側壁への再付着が発生（それ未満では再付着なし）．
    \item 双安定（流れの履歴で異なる定常解＝ヒステリシス）が出現し，$L/B$ 増で双安定域が拡大．
    \item $d_\mathrm{in} \lesssim -3B/8$ なら再付着が確実＝設計の頑健側．Camnasio 2013 の L-R 土砂遷移（双安定で再付着なし側へ）とも整合．
  \end{itemize}
  \item \textbf{限界・適用範囲}：数値のみ（実験は Camnasio 2013 の C-C/L-L/L-R/C-R を再現）．単一水理条件．入口と出口は別々で，可逆な単一ポートではない．深さ平均で 3D 二次流は陽に扱わない．
  \item \textbf{本研究との接点}：Camnasio 2013 が非中心配置を 3 形状で扱ったのに対し，本論文が初めて入口・出口位置を系統的に振った数値研究で，本研究のポート配置・可逆性の検討に最も近い既往研究．双安定・ヒステリシスはモード非対称そのもの．可逆運転での流向反転は入口と出口の入れ替えに相当し，$(d_\mathrm{in}, d_\mathrm{out})$ 平面が参照枠となる．「双安定配置は設計で避ける」が実務指針．
\end{itemize}
\noindent（確認: p.689--695）

\noindent\litref{Westhoff et al. (2018)}\quad 対称から非対称再付着流への遷移を，せん断層の最大エネルギー散逸原理で説明できるかを検討．
\begin{itemize}
  \item \textbf{問い}：矩形浅水池で対称流から非対称再付着流へ切り替わる理由を，せん断層の最大エネルギー散逸原理で説明できるか．
  \item \textbf{手法}：既往の実験・2D 浅水流計算で観測された対称/非対称流況遷移を対象に，新たな浅水流数値計算ではなく簡略エネルギー散逸モデルで遷移を説明しようとした．主ジェットと両側再循環域の間のせん断層に着目し，対称/非対称の流れ場を三角形分割で近似．$c_s/c_b$ を変えて各流況でせん断層散逸が最大となる状態を求め，最大散逸が大きい流況を実現流況とみなす．遷移位置を Dufresne et al.\ の経験的流況分類と比較．
  \item \textbf{主要知見}：
  \begin{itemize}
    \item 短い池では対称流の方が，長い池では非対称流の方が，最適化後のせん断層散逸が大きくなり，観測される流況遷移と定性的に一致．
    \item せん断層摩擦は，主ジェットから再循環域への運動量・エネルギー輸送の代理量として扱われる．
    \item 最大散逸は，著者らの解釈では最大パワー/最大エントロピー生成と対応するが，なぜ系がその状態を選ぶかの決定的説明はまだない．
    \item 再付着長や小循環中心はモデル内で完全には決まらず，観測値を与えたときに最もよく遷移を再現する．
    \item 小スケール乱流を詳細に解かずに，マクロな流況遷移を予測できる可能性を示す．
  \end{itemize}
  \item \textbf{限界・適用範囲}：対称流と最初の非対称流況のみ．より長い池の第 2 再付着・側壁跳び・meandering jet は対象外．ジェット幅・速度・水深一定，側壁摩擦無視，ジェット/再循環域の明確分離など強い簡略化を含む．
  \item \textbf{本研究との接点}：対称/非対称モード選択を，線形安定解析や Coanda 効果とは別のエネルギー散逸原理から説明しようとする候補理論．ただし可逆ポート問題へは，直接の予測式ではなくモード選択を考える概念枠組みとして使うのが安全．
\end{itemize}
\noindent（確認: p.221--230）

\noindent\litref{Valero et al. (2022)}\quad 非定常 URANS（3D）で浅水矩形池の蛇行を再現し，3D 効果と 2D の限界を POD で評価．
\begin{itemize}
  \item \textbf{問い}：非定常 URANS（3D）は浅水矩形池の蛇行噴流をどこまで再現できるか．3D 効果は何か．
  \item \textbf{手法}：FLOW-3D（URANS, VOF），3 乱流モデル（RNG $k$-$\varepsilon$, $k$-$\varepsilon$, $k$-$\omega$）．Peltier 2014b/2015 の 4 ケース（F/FT/NFT/NF，摩擦数 $S = 0.18 \to 0.03$）．$L=1$~m, $B=0.98$~m, $b=0.08$~m．自由表面の最エネルギー 4 モードを POD 抽出し，実験(LSPIV)・2D(WOLF)と比較．渦は Q-criterion で同定．メッシュ感度（鉛直 9 段）．
  \item \textbf{主要知見}：
  \begin{itemize}
    \item 3D URANS はモード周波数を良好に再現（実験の 88.7~\%），振幅(70.7~\%)・エネルギー(47.7~\%)はやや劣る．摩擦支配ケースでモード減衰，非摩擦ケースで増幅．3 乱流モデルの差は小さい．
    \item 蛇行は 2 種の不安定（sinuous＝モード 1,2 / varicose＝モード 3,4）の相互作用．
    \item 3D は 2D が捉えない局所鉛直構造（変曲点が深さで異なる）を捕捉．自由表面の 2 次元 footprint はより複雑な 3D 現象の投影にすぎず断続的．入口でジェットコアに渦が注入され蛇行を駆動，出口付近で渦が伸張・強化．
  \end{itemize}
  \item \textbf{限界・適用範囲}：工学的視点（広く使う URANS+等方渦粘性の評価）で LES は対象外．可逆ポートなし．
  \item \textbf{本研究との接点}：
  \begin{itemize}
    \item 表面 PIV で見る蛇行は 3D 構造の footprint にすぎないという留意点（Kantoush 2008a の壁近傍 2 次流，Dracos の中間場と整合）．POD/Q-criterion は非定常解析の手法参照．
    \item 一方で，主対象が定常的な A1/A2 型の再付着モードなら，本論文は中心文献というより，非定常 meandering 領域まで視野を広げる補助文献．
    \item 摩擦数 $S$ で蛇行モードのエネルギーが変わる＝§5 スケーリングと接続．Peltier 2015（2D 版）と対．
  \end{itemize}
\end{itemize}
\noindent（確認: p.12--20）

\noindent\litref{Peltier et al. (2015)}\quad 2D 深さ平均浅水モデルが蛇行流をどこまで再現できるかを POD で検証．
\begin{itemize}
  \item \textbf{問い}：2D 浅水方程式（深さ平均）は蛇行流を再現できるか．浅さ（底面摩擦）の影響は．
  \item \textbf{手法}：WOLF 2D（2 スケール $k$-$\varepsilon$：水平 $k,\varepsilon$ ＋鉛直 Elder）．Peltier 2014b の 4 ケース（F/FT/NFT/NF）．3 粗度高さ（$k_s=0/0.1/1$~mm）．POD で実験(LSPIV)と数値のモードを比較（モードエネルギー・時間/空間モード）．
  \item \textbf{主要知見}：
  \begin{itemize}
    \item 平滑床でも第 1・2 モード（蛇行を担う sinuous 対モード）はエネルギー・時空間変動とも実験とよく一致（浅さによらず）．
    \item 第 3 モード以降（特に空間モード形状）は再現不十分．粗度を上げると改善＝底面生成乱流（鉛直長さスケール）が大規模構造の横拡大を抑制．
    \item 非常に浅い場合は粘性効果が残り一致が劣化．$k$-$\varepsilon$ の二長さスケールで較正不要．
  \end{itemize}
  \item \textbf{限界・適用範囲}：一方向流・可逆ポートなし．蛇行のみ（再付着/双安定は別）．
  \item \textbf{本研究との接点}：「2D 深さ平均で蛇行の主要モードは再現可能」という方法論的根拠（将来 CFD）．POD によるモード比較手法．底面摩擦（鉛直長さスケール）が大規模構造を制御＝Babarutsi/Dracos の深さスケーリングに連なる．Valero 2022（3D 版）と対をなす．
\end{itemize}
\noindent（確認: 04015008）

\subsection*{論文横断整理}

\noindent 安定性解析・分岐論の観点からの整理は §7 に集約した（本節の Dewals 2008・Mullin 2003・Westhoff 2018 を含む）．

\noindent\litref{Mullin et al. (2003)}\quad 対称急拡＋下流収縮水路の対称性破壊分岐を分岐解析＋実験で精査（安定性・分岐理論の古典）．
\begin{itemize}
  \item \textbf{問い}：対称 2D 水路の急拡（expansion ratio $E=3$）＋下流収縮（outlet ratio $\gamma$）で，対称性破壊分岐はアスペクト比 $\alpha$・出口比 $\gamma$ にどう依存するか．
  \item \textbf{手法}：有限要素分岐解析コード ENTWIFE（極限点・対称性破壊点・Hopf 点・余次元 $>1$ 特異点の経路追跡）＋水路実験（$\mathrm{Re} \approx 16$--$95$，層流，粒子追跡）．$E=3$ 固定，$\gamma=1.5$--$2$ 等．
  \item \textbf{主要知見}：
  \begin{itemize}
    \item ピッチフォーク分岐で対称流→非対称流（Fearn 1990 系）．下流収縮が対称流を再安定化（補臨界ピッチフォーク→ある $\mathrm{Re}$ 域で対称復活→さらに上で再び非対称）．
    \item $(\alpha, \gamma)$ 平面に余次元 2 点 $D$（normal form $x^3 - \lambda^3 x + \alpha x + \beta\lambda x = 0$）を含む多重分岐構造．$\gamma>1$（拡大優位）と $\gamma<1$（収縮優位）で挙動が質的に変化．
    \item 微小幾何欠陥がピッチフォークを大きく分断（Fearn 1990）．$\mathrm{Re} \gtrsim 60$ で 3D 効果が顕在化．
  \end{itemize}
  \item \textbf{限界・適用範囲}：閉じ込め水路（自由表面でない），低 $\mathrm{Re}$ 層流〜遷移．浅水池の高 $\mathrm{Re}$ 乱流レジームとは直接対応しない．
  \item \textbf{本研究との接点}：「対称急拡における対称性破壊ピッチフォーク」「下流収縮（＝出口）による再安定化」「欠陥への高感受性」はモード非対称・双安定の古典的・理論的基盤．出口の有無/位置が分岐構造を変える＝可逆ポートで流向反転時のモード選択を ODE/分岐論で論じる際の参照（Dewals 2008 の線形安定解析と相補）．
\end{itemize}
\noindent（確認: p.433--452）

\subsection*{原稿転記メモ}

% ================================================================
\section{表面 PIV と物理模型}
% ================================================================

\subsection*{主要な問い}

\subsection*{論文メモ}

\noindent\litref{Kantoush et al. (2008a)}\quad 同型の浅水貯水池で表面 LSPIV・全層 UVP・2D 数値を相互比較し，表面 PIV の妥当性を実証．
\begin{itemize}
  \item \textbf{問い}：浅水矩形貯水池の流れ場を表面 LSPIV・全層 UVP・2D 数値（CCHE2D）でどこまで一致して捉えられるか．表面速度は流れ場を代表できるか．
  \item \textbf{手法}：EPFL LCH の矩形浅水貯水池（$6.0 \times 4.0$~m，水深 0.20~m，入口/出口 0.25~m，$Q=7.0$~L/s，$n=0.01$）．表面 LSPIV（白色プラスチック粒子 3.4~mm，CMOS 33~fps），全層 UVP（Metflow 2~MHz，可動枠の 3 本組プローブで断面 3D，トレーサはクルミ殻 $d_{50}=50~\mu$m），CCHE2D（2D 深さ平均，放物型渦粘性）を相互比較．
  \item \textbf{主要知見}：
  \begin{itemize}
    \item UVP で鉛直速度は水平速度より十分小さく浅水 2D 性を確認．ただし水平軸まわりの渦（壁近傍の鉛直循環＝2 次流）が存在し完全 2D ではない．
    \item LSPIV と UVP は渦パターン・渦中心位置がよく一致．CCHE2D も両者とよく一致．
    \item 流れは入口から平面噴流として進入し片側へ偏向，大規模主渦＋角部の三角渦＋小渦を形成（非対称）．初期条件を少し変えると鏡像パターン＝双安定．
    \item LSPIV は低速・浅水でも有効で現地計測への適用可能性．CCHE2D は小さな BC/IC 変化で流向が切り替わる（高感受性）．
  \end{itemize}
  \item \textbf{限界・適用範囲}：一方向流，対称な入口/出口で可逆単一ポートではない．「表面≒深さ平均」は浅水かつ低鉛直速度条件下での近似．
  \item \textbf{本研究との接点}：本研究と同型の浅水貯水池構成で「表面 LSPIV $\approx$ 全層 UVP $\approx$ 2D 数値」を実証した，表面 PIV 採用の直接的な方法論的根拠．壁近傍 2 次流という留意点も同時に与える．シーディング・LSPIV 処理は Methods の参照．
\end{itemize}
\noindent（確認: p.139--144）

\noindent\litref{Adrian (1991)}\quad PIV の基礎総説．粒子画像流速計測の分類・設計基準・理論的性能限界を体系化（成果論文ではなく手法の体系化）．
\begin{itemize}
  \item \textbf{問い}：粒子画像による流速計測の原理・分類・理論的性能限界は何か．
  \item \textbf{手法}：総説．pulsed-light velocimetry を上位概念として手法を体系化し，性能限界を無次元数で定式化．
  \item \textbf{主要知見}：
  \begin{itemize}
    \item \textbf{分類は 2 段階・2 つの無次元数}．数える体積が違う別量なので混同しないこと．
    \begin{itemize}
      \item $N_s$（源密度）＝解像セル（粒子像 1 個に対応する流体体積）内の粒子数．$N_s > 1$ で像が重なりスペックル＝LSV．実務的な濃度ではまず起きず，「LSV とされた実験の多くは実は PIV」と指摘．
      \item $N_I$（像密度）＝検査領域（相関窓）内の粒子像数．$N_I \gg 1$ で高像密度 PIV（粒子群の平均変位を相関で求める），$N_I \ll 1$ で PTV（個別追跡）．本研究＝高像密度 PIV．
    \end{itemize}
    \item \textbf{設計基準}（Keane \& Adrian 1990，p.293）＝実務上いちばん使える部分：$N_I > 10$--$20$，面内変位 $< d_I/4$（検査窓の 1/4），面外変位 $< \Delta z_o/4$，窓内の速度変動 $|\Delta u|/|u| < 0.2$．これで検査領域の 90--95~\% で有効ベクトルが得られる．「1/4 ルール」の出典．
    \item \textbf{不確定性原理}（式 22）：$\sigma_u \, \Delta x_\mathrm{max} = c \, u_\mathrm{max} \, d_t$．速度精度と空間分解能は同時に良くできず，測定レンジ $u_\mathrm{max}$ を広げるほど積が悪化する．両者を同時に改善する唯一の道は粒子像 $d_t$ を小さくすること．
    \item \textbf{微分量には別基準}：データ密度 $N_\lambda = C \lambda_T^3$（$\lambda_T$ ＝ Taylor マイクロスケール＝速度変動／速度勾配で定義される「勾配の長さスケール」）．$N_\lambda \gg 1$ でなければ渦度・ひずみ速度は不正確．「補間で滑らかにした場は微分できるが結果は不正確」と明記．一方で結論部では「渦度とひずみ速度の測定は PIV が提供する最も歓迎すべき能力の一つ」とも述べる．
    \item \textbf{速度バイアス}：速い粒子ほど光シートを抜け・重なりやすく，検出が低速側に偏る．PTV では深刻，高像密度 PIV では小さい．相関ピークが返すのは窓内の重み付き平均速度（式 33）で，重みは低速側に偏る＝上記 $|\Delta u|/|u| < 0.2$ はその帰結．
    \item 散乱光は小粒子域（1--10~$\mu$m）で像の明るさが $d_p^3$ に比例（散乱エネルギーは $d_p^3$ で増えるが，像の大きさは回折で決まり粒径に依存しないため）．スリップ速度は $d_p^2$ に比例するので，最も明るく写る粒子が最も流れに追従していない＝均一粒径が望ましい．ただし幾何散乱域（100~$\mu$m 以上）では像の明るさは粒径に依存しなくなる（式 16）．
    \item \textbf{装置を $m$ 倍にすると像の明るさは $m^{-5}$ で低下}（式 15．$d_o$, $d_i$, $\Delta y_o$ が各 $m$ 倍のため）．大きな装置ほど PIV は原理的に不利．→本研究の basin $3 \times 2$~m という大視野で mm 級の浮遊トレーサを使う必然性の根拠になる（実験室 PIV の視野は $25 \times 35$~mm 程度）．
    \item 緑 $>$ 赤（明るさ $\lambda^{-3}$，フィルム感度，光シート最小厚 $\sim \lambda$ の 3 効果）．ただし本研究はレーザー・光シートを使わないため，この論点は適用外．
  \end{itemize}
  \item \textbf{限界・適用範囲}：1991 年時点．記録媒体（フィルム vs ビデオ），Young's fringe，光学的相関の各節は現在ほぼ歴史的意義のみ．相互相関は「更なる研究が必要」の段階で multi-pass は登場しない．time-resolved PIV は結論部の展望のみ．\textbf{光シートで内部断面を照らす前提の理論で，水面を上から撮る表面 PIV（LSPIV）は範囲外}．→面外ロス $\Delta z_o$ の基準はそのまま移せない．
  \item \textbf{本研究との接点}：
  \begin{itemize}
    \item Methods の PIV 設定（現在プレースホルダ）の選定根拠として直接引ける：$N_I > 10$--$20$，面内変位 $<$ 窓の 1/4．PIVlab の推奨値もこの系譜．
    \item \textbf{dual-$\Delta t$ の理論的理由}：速度のダイナミックレンジが広いと，長い $\Delta t$ は高速域で 1/4 ルールを破り，短い $\Delta t$ は低速域で変位がピクセル誤差に埋もれる．不確定性原理への直接の対処．
    \item $I_\mathrm{circ}$（渦度）の妥当性：$N_\lambda$ の議論＝ベクトル間隔が微細スケールに対し十分かを Methods で示す根拠．ただし $\lambda_T$ は等方性乱流を念頭にした量で，本研究の準 2 次元浅水流（$h/W \approx 0.02$--$0.05$）では水深 $h$ やジェット幅の方が自然な物差しになりうる→要検討．
    \item $\phi_\mathrm{lv}$（低流速域面積率）を主要指標とする以上，速度バイアスは認識しておくべき論点．逆に高像密度＋相関法を採る正当化材料にもなる．
    \item 表面 PIV の直接の先行例は Fujita 1998（LSPIV），実装は Thielicke \& Stamhuis 2014（PIVlab）．本論文はその下敷きとなる一般理論という位置づけ．
  \end{itemize}
\end{itemize}
\noindent（確認: p.261--304，全 45~p 通読）

\noindent\litref{Fujita et al. (1998)}\quad LSPIV の定礎論文＝標準的な引用先（手法自体の初出は Fujita \& Komura 1994 / Fujita et al.\ 1997 で，本論文は 3 つの適用例による体系化）．
\begin{itemize}
  \item \textbf{問い}：PIV を実験室スケール（数 cm$^2$）から水理工学の大スケール（4--45{,}000~m$^2$）へどう拡張するか．
  \item \textbf{手法}：LSPIV の体系化．従来 PIV との相違は 3 点＝(i) 自然光＋通常ビデオ（30~Hz 以下）で足りる．\textbf{理由は視野の大きさではなく「水理応用では流速が相対的に低い」から}．ただし遠距離ゆえ照度が落ちるため画像強調が必要（不適切な強調はむしろ劣化させると警告）．(ii) 射影変換（8 係数，既知座標マーカー最低 4 点，最小二乗）で歪み補正．PIV の\textbf{前}に行う．(iii) トレーサ．相関・サブピクセル（放物線フィット）・誤ベクトル除去は従来 PIV と共通．適用例：跳水模型（$1.15 \times 0.95$~m）→合流部氷輸送模型（$5.65 \times 4.14$~m）→淀川洪水（$140 \times 320$~m）．
  \item \textbf{主要知見}：
  \begin{itemize}
    \item \textbf{トレーサの思想が Adrian 1991 と異なる}．個々の粒子像ではなく「流れに乗る灰色レベルの一貫したパターン」を追うため\textbf{粒子サイズは本質的に重要でない}．自然のパターン（漂流物・泡・さざ波）が使え，不均一な場合にのみ追加散布する．大視野の本研究に適用されるのはこちらの立場．
    \item \textbf{精度は平均流速で 3--5~\%}（結論部の明記）．個別には ADV と 10~\% 以内，流量が水位流量曲線と 3~\% 以内．
    \item \textbf{乱れは測れない}．跳水の例で LSPIV の変動速度 RMS が ADV の局所乱れ強度より\textbf{1 桁大きい}．トレーサ（2--3~mm 超の気泡・ビーズ）が乱流の微細構造に追従しないため．提供できるのは「平均場の記述」＝解像されるのは大スケール構造のみ（著者は「水理学的記述」と表現）．
    \item 解像限界の指針：検査領域は解像したい最小長さスケールの 1/5 以下（Stevens \& Coates 1994）．
    \item 検査領域サイズへの鈍感性：物理座標で $10 \times 10$~m 程度までは流量計算がこのパラメータに依存しない．
    \item 探索領域は連続画像から最大変位を目視同定して決める（計算速度のため平均流速場の関数として選ぶ）．
  \end{itemize}
  \item \textbf{限界・適用範囲}：表面速度のみ（深さ方向は 1/7 乗則などの仮定に依存）．トレーサ分布・照明条件に品質が左右される．試行錯誤が不可避で「技術と流れ特性への完全な習熟が前提」と著者が明記．岸近傍は精度低下．3 例はいずれも単方向流で可逆流・閉鎖 basin ではなく，time-resolved な使い方の議論もない（$\Delta t = 1$~s，時間平均が主目的）．1998 年の機材だが\textbf{アルゴリズムと補正理論は現在も有効}．
  \item \textbf{本研究との接点}：
  \begin{itemize}
    \item 本研究の表面 PIV ＋射影変換という手法構成の直接の出典．Methods は Adrian 1991 が理論的基準（$N_I > 10$--$20$，面内変位 $<$ 窓の 1/4），本論文が大視野での実装指針，という二段構えで引ける．\texttt{doc/piv\_preprocessing\_protocol.md} の射影変換手順の引用根拠．
    \item \textbf{$I_\mathrm{unst}$ の解釈に直結}：LSPIV の速度変動は真の乱れより 1 桁大きく出るため，$I_\mathrm{unst}(B_k)$ を「乱流強度」と解釈してはならない．捉えているのは basin スケールの大規模非定常性（ジェット振動・循環セルの変動）．この区別を Methods か Discussion で明示すべき．
    \item $I_\mathrm{circ}$（渦度）の分解能の議論は，Adrian の $N_\lambda$（等方性乱流前提の Taylor スケール）より，本論文の「検査領域 $<$ 解像したいスケールの 1/5」の方が浅水流には実用的な物差しになりうる．
    \item 精度 3--5~\% は本研究の計測精度を述べる際の標準的な引用値．検査領域サイズへの鈍感性は Appendix の感度解析を設計する際の先行例．
    \item Kantoush 2008a が浅水貯水池で LSPIV の妥当性を示した，その手法自体の出典（Kantoush 側は未確認）．
  \end{itemize}
\end{itemize}
\noindent（確認: p.397--414）

\noindent\litref{Thielicke \& Stamhuis (2014)}\quad（補助）使用ツール PIVlab の出典（software metapaper）．
\begin{itemize}
  \item \textbf{内容}：MATLAB の GUI ベース DPIV ツール PIVlab の紹介．前処理（CLAHE / intensity capping / highpass）→相関（DCC or DFT，多重パス，窓変形）→ピーク検出（ガウス $2 \times 3$ 点 or 9 点）→検証（速度閾値・正規化中央値）→内挿（boundary value solver）→平滑化，の各段でアルゴリズムを比較し最良を実装．Fig.~2 のワークフロー図が設定項目のチェックリストになる．
  \item \textbf{主要知見}：
  \begin{itemize}
    \item 最良は\textbf{窓変形つき多重パス DFT}．DFT は A と B が同サイズゆえ変位とともに背景ノイズが増えるため\textbf{変位を検査領域の 1/4 に抑える}のが望ましい（Keane \& Adrian 由来）．DCC は変位 32~px までノイズが増えないが 1--2 桁遅い．
    \item 最適条件下でバイアス誤差 $< 0.005$~px，ランダム誤差 $< 0.02$~px．ただし最適条件＝粒子像直径 $\approx 3$~px，密度 5--15 個/窓，ノイズ・粒子対ロス・モーションブラー・せん断なし＝実験では達成し得ない理想値．
    \item 前処理は有効ベクトル検出確率を CLAHE で $4.7 \pm 3.2$~\%，capping で $5.2 \pm 2.5$~\% 改善．ただし\textbf{intensity highpass は低周波の変位情報も抑制する}と明記．
    \item 平滑化はどの手法でも品質を上げるが，効くのは主に外れ値（最大絶対誤差 $47 \to 17$~\%，平均絶対誤差は $1.26 \to 0.82$~\% と差は小さい）．
  \end{itemize}
  \item \textbf{限界・適用範囲}：\textbf{精度評価は合成画像のみ}で，実画像による検証は前処理の有効ベクトル率改善に限られる（商用ソフトとの比較は本論文にはない）．精度の詳細は本論文になく Thielicke 2014 博士論文が一次出典．レーザーシート＋中性浮力粒子が前提で表面 PIV は想定外（アルゴリズム自体は画像があれば動く）．2014 年版（MATLAB R2010a ベース）で現行版とは異なる．
  \item \textbf{本研究との接点}：
  \begin{itemize}
    \item 本研究が PIVlab を用いる場合の直接の引用先．理論的背景は Adrian 1991，水理大スケール化は Fujita 1998，本論文は実装を担う三段構え．使用バージョンを Methods に明記すること．
    \item Methods の PIV 設定は Fig.~2 の順に記述すれば漏れがない（AGENTS.md が \texttt{pivlab\_setting*.mat} を authoritative な記録とする方針は，この情報量からして妥当）．
    \item \textbf{精度値をそのまま引かないこと}：$0.005$/$0.02$~px は理想条件の合成画像の値．大視野・mm 級トレーサの LSPIV である本研究の精度の目安は Fujita の 3--5~\% 側．
    \item dual-$\Delta t$ の実装上の理由：DFT の「変位 $<$ 窓の 1/4」を単一 $\Delta t$ では低速域と高速域の両方で満たせないため．
    \item \textbf{highpass は慎重に}：basin スケールの緩やかな流れを見る本研究では低周波成分を落とす副作用が効く．CLAHE と capping は比較的安全．
  \end{itemize}
\end{itemize}
\noindent（確認: art.\ e30）

\noindent\litref{Berkooz et al. (1993)}\quad（補助）POD の基礎総説（全 37~p 通読）．本研究は現時点で POD 未使用だが，着手する際はまず本論文を参照する．
\begin{itemize}
  \item \textbf{問い}：乱流解析における POD とは何か．なぜ最適なのか（Lumley 1967 の導入者自身による総括）．
  \item \textbf{手法}：総説．2 点相関 $R$ の固有値問題としての定式化と性質の体系化．証明は Berkooz 1991b 学位論文．
  \item \textbf{主要知見}：
  \begin{itemize}
    \item \textbf{最適性（核心）}：任意の線形分解の中で「与えられたモード数で平均的に最大の運動エネルギーを捉える」意味で最適（Prop.~2.3）．係数 $a_i(t)$ は互いに無相関．固有値 $\lambda_i$ が第 $i$ モードの平均エネルギー．＝「上位 $n$ モードで全 TKE の $X$~\%」に意味があるのはこの最適性による．
    \item \textbf{モード数の実用的基準}（Sirovich 1989）：捕捉エネルギー 90~\% 以上，かつ無視するどのモードも平均して第 1 モードの 1~\% 未満．「非ゼロ固有値の数」を次元と呼ぶのは誤り（Hausdorff 次元と無関係）．
    \item \textbf{スナップショット法}（Sirovich 1987）：$N \times N$ を $M \times M$ に縮約．PIV データへの適用にはこれが前提．
    \item \textbf{対称性}：2 点相関が並進不変なら固有関数は Fourier モードに一致する（Prop.~2.4，必要十分）＝一様方向では新情報を与えず，非均質な方向でこそ有用．basin は全方向に非均質で POD が活きる設定．
    \item 対称性が破れた実現（正方 Rayleigh--B\'enard セルで単一ロールの回転方向が発生時にランダムに選ばれ以後変わらない）は非エルゴード性を示唆し，固有関数の対称性検査で検出できる（Prop.~2.6）．この場合\textbf{対称性を使ってアンサンブルを人為的に増やすと真の性質を覆い隠す}と警告．
  \end{itemize}
  \item \textbf{限界・適用範囲}：\textbf{POD は線形手続き}で，著者自ら「線形性はその限界の源泉」と明記．エネルギーが唯一の基準ゆえ，小さいが力学的に重要なモードは切り捨てられる．1993 年であり DMD は登場しない．
  \item \textbf{本研究との接点}：
  \begin{itemize}
    \item 対称性の破れ＝非エルゴード性の論点は本研究の双安定（左右どちらに倒れるかが初期条件で決まる）と同型で，4 回反復の設計はこの問題への対処と読める．ただし $I_\mathrm{asym}$ は瞬時場の左右エネルギー比であって固有関数の対称性検査ではなく，理論的裏づけと呼ぶには飛躍がある．
    \item \textbf{注意}：本論文の文脈は乱流のコヒーレント構造．本研究の LSPIV は乱流を解像しない（Fujita 1998）ため，POD を使う目的は「basin スケールの大規模流動パターンの分解」であり乱流構造の抽出ではない．
    \item 役割分担：理論は本論文＋Taira 2017，浅水への適用は Valero 2022，実装例は Bai 2026．
  \end{itemize}
\end{itemize}
\noindent（確認: p.539--575）

\noindent\litref{Sirovich (1987, I--III)}\quad スナップショット POD の原典．3 部作を 1 メモに統合（図は一枚もない）．Part I: Coherent structures (p.561--571) / Part II: Symmetries (p.573--582) / Part III: Dynamics (p.583--590)．
\begin{itemize}
  \item \textbf{問い}：大規模な乱流データからコヒーレント構造を実用的に抽出するにはどうすればよいか．
  \item \textbf{主要知見}：
  \begin{itemize}
    \item \textbf{Part I／スナップショット法＝最重要}：直接法は格子点数 $N$ の固有値問題になり計算不能（$3 \times 10^3$ 格子ですら当時の容量超）．相関時間以上に離れた $M$ 個の瞬間場から $K$ を構成すると $M \times M$ に縮約される（$CA = \lambda A$）．サンプリング点数は計算に本質的に入らないため，空間解像度が高い PIV データでも行列サイズは $M$ で決まる．実装は「データ行列を SVD」．非圧縮なら固有関数も自動的に非圧縮かつ境界条件を満たす．エネルギー 99~\% 捕捉を実用基準として提示（Berkooz が引く「90~\% かつ第 1 モードの 1~\% 未満」は Sirovich 1989 の別基準）．
    \item \textbf{Part II}：幾何の対称群（並進・鏡映・回転）で既存データを水増しできる．正方断面なら 8 倍，正方 B\'enard なら 16 倍．副次効果として，ゼロであるべき量が群平均で正確にゼロになる．
    \item \textbf{Part III}：固有関数を基底に NS を Galerkin 射影して ODE 系 $da_k/dt = F_k(a)$ を得る．展開は全流れ量でなく\textbf{平均からの変動}について行うべきと明言（打ち切った固有関数集合では平均流が十分適合されず，固有値の順序だけで基底を選べなくなるため）．
    \item \textbf{留保}：実験で言うコヒーレント構造（例：ヘアピン渦）は本定義の固有関数と一致しない可能性．固有関数はむしろ基本モードで，それらが時間的に相互作用して観察される構造を作る，と推測．
  \end{itemize}
  \item \textbf{限界・適用範囲}：3 部作とも具体的な数値例を含まない（方法論の提示）．Part II/III は均質方向・周期境界を前提とする議論が中心で，閉じた basin には直接適用できない．
  \item \textbf{本研究との接点}：
  \begin{itemize}
    \item スナップショット POD の原典．本研究で POD を使う場合の実装上の前提．
    \item Part III の「変動を展開すべき」は，POD を何に適用するかの選択の根拠として Methods に書ける．
    \item \textbf{Part II の対称性拡張は本研究では使わない方がよい}．basin は中心線に鏡映対称だが，本研究の主題は左右非対称（$I_\mathrm{asym}$）で偏向方向が実現ごとに選ばれうる．水増しすると測りたい非対称を人為的に消す．Berkooz 1993 の非エルゴード性の警告がそのまま当てはまる場面．
  \end{itemize}
\end{itemize}
\noindent（確認: p.561--571, p.573--582, p.583--590）

\noindent\litref{Taira et al. (2017)}\quad（補助）モード解析の現代的総説（9 名共著）．現時点で本研究は POD/DMD 未使用のため補助扱い．使う段になれば手法選択の地図として最初に見る．
\begin{itemize}
  \item \textbf{読解範囲}：II--V 章（固有値/特異値分解・POD・balanced POD・DMD）を精読．Koopman・大域安定性・resolvent の各章は未読．
  \item \textbf{問い}：流体解析で使えるモード分解手法には何があり，それぞれ何を与え何が弱点か．
  \item \textbf{手法}：総説．データベース手法（POD, balanced POD, DMD ＝流れ場データが入力）と作用素ベース手法（Koopman, 大域安定性, resolvent ＝支配方程式が必要）を線形代数の枠組みで整理し Table~1 に一覧．＝\textbf{実験データしか無い場合に使えるのは前者のみ}という切り分けが明快．
  \item \textbf{主要知見}：
  \begin{itemize}
    \item \textbf{POD}：エネルギー（$L^2$）で最適．長所は直交・低次元・\textbf{非コヒーレントなノイズが高次モードに現れるため高次を捨てるだけで除去できる}（ノイズが信号より小さい限り）．短所 4 つ：2 次相関のみ／\textbf{時間係数が複数周波数の混合になる}／エネルギー順であり力学的重要度順ではない／保持モード数の決め方に決定打がない．
    \item \textbf{DMD}：各モードが単一周波数と成長/減衰率をもち，POD の周波数混合の弱点を補う．「POD と離散 Fourier 変換の好ましい面を組み合わせたもの」．短所：モードが直交せず客観的順位がない／時間解像データが必要／非線形性が重要な場合に信頼できない結果を与えうる．
    \item \textbf{Spectral POD}：Welch 法でブロック分割し各周波数の相互スペクトル密度行列を固有値分解．離散周波数ごとに直交モードを与え，かつエネルギーで順位づけられる＝周波数混合への POD 側の解．
    \item \textbf{進行波はモード対で表れる}（式 34）：$\sin(\xi - ct) = \cos(ct)\sin(\xi) - \sin(ct)\cos(\xi)$．実数値 POD では進行構造を単一モードで表現できず，移流方向に位相 $\pi/2$ ずれた 2 モードの組で現れる．Fig.~6 の例ではモード 1・2 のエネルギーが 5.8~\% ずつと完全に等しい．
    \item 実装：スナップショット法が最も広く使われるが，\textbf{SVD ベースの方が丸め誤差に頑健}．MATLAB は \texttt{svd(A,'econ')} が reduced SVD（Table~2）．balanced POD は随伴計算が要るため実験では不可．
  \end{itemize}
  \item \textbf{限界・適用範囲}：総説で新規データなし．例はすべて航空宇宙の高速流（翼・キャビティ）で水理・浅水流は扱わない．各手法の記述は簡潔で実装詳細は原論文に譲る（「一つの午後で入門を得る」ための文書）．
  \item \textbf{本研究との接点}：
  \begin{itemize}
    \item 手法選択の判断基準：エネルギーで分けるなら POD，周波数で分けるなら DMD，両方なら spectral POD．outflow の非定常性に周波数情報が要るなら DMD か spectral POD が候補．
    \item \textbf{進行構造は POD モード対で表れる}は Brevis 2011 で観察された現象の一般的定式化であり，本研究で inflow のジェット振動を POD にかけた際の解釈の理論的根拠（Peltier 2014/2015 や Valero 2022 が蛇行を第 1・2 モード対で捉えたとの伝聞は未確認）．
    \item \textbf{POD の時間係数は周波数が混ざる}ため，POD 係数の PSD が単一ピークになる保証はない．Brevis で綺麗に分離したのは周期的渦列という特殊例．$I_\mathrm{unst}$ の補完に使う際の注意．
    \item 浅水流での機能は本論文からは分からない（Brevis 2011 / Valero 2022 で補う）．
  \end{itemize}
\end{itemize}
\noindent（確認: p.4013--4041，II--V 章）

\noindent\litref{Brevis \& Garc\'ia-Villalba (2011)}\quad POD を浅水流に適用した先行例．
\begin{itemize}
  \item \textbf{問い}：PIV/LDA を使わず，染料可視化の画像列に POD をかけるだけで浅水コヒーレント構造の空間特性と発生周波数を定量化できるか．
  \item \textbf{手法}：浅水円柱後流（渦列）の染料可視化．$D = 125$~mm，$h = 18$~mm，$R_D = 16{,}200$，浅さ数 $S = 0.06$．民生ビデオ 25~Hz を\textbf{感度解析により 5~Hz にリサンプル}し 600 スナップショット（低いと係数の時間発展を解像できず，高いとモード収束に過大な枚数を要する）．平均を引いた変動場にスナップショット POD．
  \item \textbf{主要知見}：
  \begin{itemize}
    \item \textbf{モードは対で現れる}：$\varphi_1/\varphi_2$ と $\varphi_3/\varphi_4$ が酷似し固有値もほぼ等しく，係数は時間シフトしているだけ＝同一構造の空間シフト版．構造が一定速度で移流されるときの標準的状況．
    \item 第 1・2 モード＝流下方向の移流，$f = 0.2$~Hz で LDA と厳密に一致．第 3・4 モード＝横方向運動，$f = 0.4$~Hz ＝その 2 倍．＝縦移流と横振動が別のモード対に分離される．上位 4 モードで約 61~\%．
    \item \textbf{POD をフィルタとして使える}：4 モード再構成で大規模 2DCS は再現され，底面摩擦による小スケールの筋状パターンは落ちる．
    \item \textbf{未解決}：速度場で報告されている第 1 のモード対が本研究には現れず，著者は理由を示せない．候補は「速度場とスカラー場の POD モードは一致する必要がない」．
  \end{itemize}
  \item \textbf{限界・適用範囲}：濃度場の POD であり速度場ではない（上記の未解決点）．著者自ら「PIV/LDA の代替ではなく補完．予備段階で着目すべき周波数・領域を絞るのに役立つ」と位置づける．円柱後流であり貯水池・可逆ポートではない．
  \item \textbf{本研究との接点}：
  \begin{itemize}
    \item \textbf{POD $\times$ 浅水流}の直球の先行例．水深も $h = 0.018$--$0.04$~m と本研究（$0.03$--$0.10$~m）と同オーダー．ただし流れの型は周期的渦列で basin 内循環ではない．
    \item \textbf{モード対の現れ方}は本研究の POD 結果を解釈する雛形．inflow のジェット振動が固有値のほぼ等しいモード対として出る可能性が高く，事前に知らないと「なぜ Mode 1 と 2 が同形か」で戸惑う．
    \item サンプリング周波数を感度解析で決める作法は本研究にも移せる．POD 係数の PSD は $I_\mathrm{unst}$ を補完しうる（$I_\mathrm{unst}$ は変動の大きさ，PSD は「どの空間構造がどの周波数で振動するか」まで分かる）．
  \end{itemize}
\end{itemize}
\noindent（確認: p.586--594）


\noindent\litref{Sciacchitano (2019)}\quad（補助）PIV 不確かさ定量化の現代的総説（Topical Review，Open Access，全 31~p 通読）．
\begin{itemize}
  \item \textbf{問い}：PIV の測定不確かさをどう客観的に定量化するか．誤差源は何で，どの手法が使えるか．
  \item \textbf{手法}：総説．ISO-GUM に従い誤差（random/systematic）と不確かさを定義し，計測連鎖の各段の誤差源を Fig.~4 に体系化（主に系統的／主にランダムを色分け）．UQ 手法を a-priori（合成画像・モンテカルロ・実験的評価で手法一般の目安を出す）と a-posteriori（実データから局所・瞬時の不確かさを出す）に大別し，後者 5 手法を Table~4 で比較．§5 で導出量・統計量への不確かさ伝播を扱う．
  \item \textbf{主要知見}：
  \begin{itemize}
    \item \textbf{「PIV の不確かさは 0.1~px」は単純化しすぎ}と冒頭と結論で明言．実験ごと・場所ごと・時刻ごとに大きく変わり，「予備的で粗い推定」としてのみ扱うべき．→単一の代表値で済ませない．
    \item \textbf{結像系}：誤差は粒子像径 $d_\tau \approx 2$~px で最小（Fig.~5）．1~px 未満だと peak locking（変位が整数画素へ引き寄せられる）が支配し，これだけで 0.1~px オーダー．大きすぎるとランダム誤差が径に比例して増加．$32 \times 32$ 窓での最小誤差は 0.05--0.1~px．モンテカルロでの典型値は 0.01--0.2~px．
    \item \textbf{トレーサ}：スリップ速度は粒子径の 2 乗に比例．液体中は密度差が小さく 10--50~$\mu$m で十分（Melling 1997）．→本研究の mm 級トレーサはこの基準から外れるが，Fujita 1998 の LSPIV は「個々の粒子でなくパターンを追う」立場なので前提が異なる．
    \item \textbf{a-posteriori 5 手法}：uncertainty surface／correlation SNR metrics（PPR, PRMSR, PCE 等）／particle disparity／correlation statistics／error sampling．比較評価では\textbf{correlation statistics が最も信頼できる}（推定精度が多くの場合 85~\% 超＝実誤差 RMS との食い違いが 15~\% 未満．coverage とは別量で，Fig.~13(g) の coverage は 65~\% 程度）．SNR metrics は相関面の情報だけで済み追加計算コストがほぼゼロ．粒子視差法と相関統計法は窓内粒子数が有限のため推定値自体に $(2N_I)^{-1/2}$ ＝ 5--25~\% の内在ばらつきがある．
    \item \textbf{error sampling}（§3.2.2.4）：装置まで含めた系統誤差を捉えられる唯一の手法．実験をできるだけ多くの側面を変えながら繰り返す．Fig.~11 は同一装置で 1 ヶ月毎日繰り返した例で，\textbf{速度の変動は標準的な不確かさ伝播が予測する値より遥かに大きかった}．ただし瞬時ベクトルのランダム不確かさは出ず，「実務では往々にして不可能」．
    \item \textbf{渦度の不確かさ（式 16）}：$U_\omega = (U/h)\sqrt{1 - \rho(2h)}$．$U$ は速度の局所標準不確かさ，$h$ は格子間隔，$\rho(2h)$ は距離 $2h$ の 2 点における\textbf{計測誤差の相互相関係数}．誤差が空間相関しているほど渦度の不確かさは下がる．報告値は $32 \times 32$ 窓で\textbf{オーバーラップ 75~\% なら $\rho = 0.45$，50~\% 以下では $\rho = 0$}．＝オーバーラップの選択が微分量の不確かさに効く．関連して Luff 1999：渦度の不確かさが最小になるのは標準的な 3 点中心差分＋局所多項式回帰・ガウシアン平滑化の併用．Fouras \& Soria 1998：空間サンプリング間隔を広げるとバイアスが増え\textbf{ピーク渦度を過小評価する}．
    \item \textbf{統計量}（§5.1）：平均値の不確かさには有効サンプル数 $N_\mathrm{eff}$ を使うべき．time-resolved PIV では取得周波数が高くサンプル間に時間相関があり $N_\mathrm{eff} \ll N$ になる．バイアス不確かさの推定は困難で「現時点では uncertainty surface 法のみが直接アクセスを提供する」．
    \item \textbf{推奨}：有限窓による空間変調の系統不確かさを定量するには，異なる検査窓サイズで処理して結果のばらつきを報告せよ（§5.1 末尾）．
  \end{itemize}
  \item \textbf{限界・適用範囲}：平面 PIV が主で\textbf{レーザーシート＋中性浮力粒子の実験室 PIV が前提．LSPIV・表面 PIV は想定外}．射影変換の誤差や自由表面の擾乱は扱わない．a-posteriori 手法は\textbf{画像に符号化された不確かさのみ}を定量でき，非一様流入・実験中の Re 変動・粒子スリップ・較正・同期は別途評価が必要と結論で明記．系統不確かさの定量はランダムに比べ研究が遅れており，較正誤差・未解像スケールによる打ち切り誤差・peak locking が主因で\textbf{非定常・乱流の研究でとくに重要}．
  \item \textbf{本研究との接点}：
  \begin{itemize}
    \item \textbf{静水ノイズフロアの位置づけ}：Table~2 の ``Flow measurement with null or uniform displacement''（Willert 1996, Forliti et al.\ 2000）＝a-priori の実験的評価に該当する．error sampling ではない．
    \item \textbf{4 回反復が error sampling に対応}．反復間のばらつきを Appendix で報告することは，単なる再現性確認ではなく本研究で実行可能な唯一の UQ．Fig.~11 の教訓（理論的伝播より実際のばらつきは大きい）を添えられる．
    \item \textbf{式 16 は $I_\mathrm{circ}$ に直接使える}．$U$ と格子間隔が分かれば計算できる．オーバーラップを埋める際に $\rho(2h)$ の議論を持っておく．ピーク渦度の過小評価も系統誤差として言及しうる．
    \item \textbf{Appendix の QC 設計}：検査窓サイズを変えた感度解析（§5.1）．Fujita 1998 の「窓サイズへの鈍感性」の確認と組み合わせると QC 節が構成できる．
    \item \textbf{outflow の低流速域に効く}：表面流速が小さく変位がピクセル誤差に埋もれる領域で，peak locking と $d_\tau$ の議論が最も効く．
    \item 役割分担：理論は Adrian 1991，大視野の実装は Fujita 1998，ツールは Thielicke 2014，不確かさの現代的枠組みが本論文．ただし\textbf{LSPIV 特有の誤差（射影変換・水面擾乱・トレーサの表面追従）は本論文の射程外}で，そこは Fujita の 3--5~\% 側で押さえる．
  \end{itemize}
\end{itemize}
\noindent（確認: art.\ 092001，全 31~p 通読）

\subsection*{論文横断整理}

\subsection*{原稿転記メモ}

% ================================================================
\section{スケーリングと相似則}
% ================================================================

\subsection*{主要な問い}

\subsection*{論文メモ}

\noindent\litref{Dracos et al. (1992)}\quad 浅水平面噴流を水深 H でスケールし，近傍/中間/遠方場と蛇行の起源を解明した基盤研究．
\begin{itemize}
  \item \textbf{問い}：有界（浅い）流体層中の平面乱流噴流は，水深 H を長さスケールとしてどのような場構造をもち，いつ蛇行を始めるか．
  \item \textbf{手法}：水深 H の有界層への平面噴流（スロット幅 $B=1$~cm，$H=20$--$360$~mm，$H/B=2$--$36$，$\mathrm{Re}_0 \approx 10^4$）．LDA 計測＋染料・紙片可視化．受水域 $5 \times 6$~m，深さ 1~m．
  \item \textbf{主要知見}：
  \begin{itemize}
    \item 通常の平面噴流は理想的には奥行き方向に一様な 2 次元乱流噴流だが，浅水噴流は水深方向の拘束のため，途中から 2 次流（上下境界による乱流非等方性ゆえの第二種二次流れ的なヘリカル循環）や準 2 次元渦が支配的になる．
    \item 水深 H をスケールに $\xi'=x/H$ で 3 領域に分割：近傍場 $\xi' \lesssim 2$（古典 2D 噴流的，B でスケール），中間場 $\xi' \approx 2$--$10$（2 次流が全層に及び中央面計測は代表性を失う），遠方場 $\xi' \gtrsim 10$（2 次流が消え噴流が中央面まわりに蛇行）．
    \item 場構造は $x/B$ でなく $x/H$ が支配＝水深が本質的スケール．
    \item 遠方場の蛇行（$\xi' \approx 10$ 〜）は両側の大規模な反回転準 2D 渦列（ペアリングで成長）を伴い，準 2D 乱流（エンストロフィー・カスケード，スペクトル $-3$ 域）．蛇行に伴う準 2D 乱流が巻き込み流体の混合を有意に増す．
    \item 平均速度は自己相似（$U/U_m=\exp(-A\eta^2)$，$A=0.693$），渦通過 $\mathrm{St}_v \approx 0.07$．
  \end{itemize}
  \item \textbf{限界・適用範囲}：大型受水域への自由平面噴流で，再付着を伴う閉じ込め貯水池や可逆ポートではない．理想化形状だが，深さスケーリングと蛇行起源は一般性が高い．
  \item \textbf{本研究との接点}：浅水噴流で「水深 H を長さスケールにとる」根拠の原典．Peltier の $S=\lambda\Delta B/8H$ や蛇行波長の H 正規化，Chagdali の $H/\Delta B$ はこの深さスケーリングに連なる．蛇行の起源（準 2D 反回転渦列）として §3（Miozzi, Peltier）と接続．
\end{itemize}
\noindent（確認: p.587--614）

\noindent\litref{Babarutsi et al. (1989)}\quad 浅水再循環流の底面摩擦数 $S=c_f d/(2h)$ と二長さスケール枠組みを確立した原典．
\begin{itemize}
  \item \textbf{問い}：浅水開水路の再循環流（急拡部）で，底面摩擦は再付着長・再循環流量・乱流をどう変えるか．支配する無次元数は．
  \item \textbf{手法}：開水路（幅 0.610~m，水深可変，長さ 7~m）の片側急拡（拡幅 0.305~m）．熱膜流速計．5 試験（うち 2 試験は床にワイヤメッシュで摩擦増）．$S=c_f d/(2h)=0.0098$--$0.2529$．
  \item \textbf{主要知見}：
  \begin{itemize}
    \item 再循環流は 2 つの長さスケール（拡幅 $d$ と底面摩擦長 $h/c_f$）で特徴づけられ，その比が底面摩擦数 $S=c_f d/(2h)$．
    \item 再付着長 $L$：深水（$S<0.05$）で $L \approx 8d$，$S$ 増で短縮，浅水極限で $L \approx 1.2\,h/c_f$（水平長 $d$ に非依存）．再循環流量も同様（$q \approx 0.08\,U_o d \to 0.018\,U_o h/c_f$）．
    \item 底面摩擦の二面性：小スケール乱流を生成しつつ，横方向せん断（大スケール乱流）を安定化・抑制．浅水ほど乱流・再循環域が縮小．
    \item $S$ による分類：深水 $S<0.05$／遷移 $0.05$--$0.1$／浅水 $S>0.1$．島後流（現地）とも整合．
  \end{itemize}
  \item \textbf{限界・適用範囲}：片側急拡（後ろ向きステップ型）の一方向流で可逆ポートではない．熱膜は流向反転を直接検出できず鏡像再構成．
  \item \textbf{本研究との接点}：摩擦数 $S$ と「$d$ と $h/c_f$ の二長さスケール」枠組みの原典で，Peltier の $S=\lambda\Delta B/8H$・Chu 2004・Chagdali の源流．とくに本研究のような「出口なし・壁が多い・流入口のみの水槽」では再循環が主要構造となり，その長さ・強さ・混合効率が底面摩擦に強く制御されうるため，関連は大きい．再付着長の摩擦依存は §2 とも接続．
\end{itemize}
\noindent（確認: p.906--924）

\noindent\litref{Babarutsi \& Chu (1991)}\quad 浅水再循環流の濃度場を計測し，摩擦数に加え閉じ込め $R=d/h$ と準 2D 乱流の視点を与える．
\begin{itemize}
  \item \textbf{問い}：浅水再循環流（急拡部）で，濃度（混合）場は底面摩擦・閉じ込めにどう支配されるか．
  \item \textbf{手法}：1989 と同一水路（拡幅 30.5→60.1~cm）で色素濃度場を吸光プローブ計測（速度場の 1989 を補完）．4 試験＋予備 5．同じ $S$ 枠組み．
  \item \textbf{主要知見}：
  \begin{itemize}
    \item 混合は底面摩擦が支配：小スケール乱流を生成しつつ大スケール横方向混合を安定化・抑制．横方向混合は約 1 摩擦長 $h/c_f$ を超えると抑えられる．壁沿い濃度は断面平均の 10--15 倍（浅水）〜4--5 倍（深水）．
    \item 2 無次元数を導入：底面摩擦数 $S=c_f d/(2h)$ と 閉じ込め $R=d/h$（幅深比）．閉じ込め効果は $R=d/h \approx 9$--$12$ で効き始める．
    \item 遷移域（$S=0.05$--$0.1$）で混合が単調変化しない異常は，摩擦の安定化と閉じ込めの不安定化が相殺するため．
    \item 浅水では大スケール運動が自由表面と床に挟まれ準 2D 化 → 逆カスケード（Kraichnan 1967）．濃度場はモデル較正に有用（速度場は乱流モデルに鈍感だが濃度場は識別力が高い）．
  \end{itemize}
  \item \textbf{限界・適用範囲}：片側急拡の一方向流で可逆ポートではない．色素（濃度）研究．
  \item \textbf{本研究との接点}：摩擦数に加え「閉じ込め $R=d/h$」と準 2D 乱流（逆カスケード）の視点．とくに閉じた／出口なし水槽で染料・土砂・トレーサ分布を見る場合，速度場が似ていても濃度場はかなり異なりうる：底面摩擦が強いと混合が抑えられる一方，水深が浅いことによる準 2D 化は大規模渦を強め混合を促進する場合がある．混合・分散の物理は §6（土砂・堆積）に直結．濃度場によるモデル較正は手法面で参考．
\end{itemize}
\noindent（確認: p.643--659）

\noindent\litref{Camnasio et al. (2014)}\quad（補助）形状パラメータ $S$ を運動エネルギー含有量へ拡張し，深さ平均 $k$-$\varepsilon$ の予測能を検証．
\begin{itemize}
  \item \textbf{問い}：流況（S0/S1/A1/A2）の運動エネルギー含有量は形状パラメータ $S$ でどう整理できるか．深さ平均 $k$-$\varepsilon$ はそれを予測できるか．
  \item \textbf{手法}：EPFL $6 \times 4$~m 池，11 形状（$L/\Delta B=1.6$--$34.3$, $\Delta B/b=0.7$--$7.5$, $S=L/\Delta B^{0.6} b^{0.4}=3.6$--$29.7$），UVP で全面速度．WOLF 2D（深さ平均 2 スケール $k$-$\varepsilon$ vs 代数 Elder モデル）と比較．
  \item \textbf{主要知見}：
  \begin{itemize}
    \item 非次元比エネルギー $e_\mathrm{nd}$ は $\ln(S)$ に線形：$e_\mathrm{nd} = 1 - 0.5\ln(S/2S_\mathrm{cr})$（$S_\mathrm{cr}=6.5$）．$S$ が幾何（$L,B,\Delta B,b$）を 1 つの数に集約．$S$ 増（高トラップ効率側）で $e_\mathrm{nd}$（理想的な直進ジェットに比べてどれくらい運動エネルギーを保っているか）単調減少．
    \item 深さ平均 $k$-$\varepsilon$ は平均速度・比運動エネルギー・主ジェットの乱流エネルギーを良好に再現．代数モデルはジェット拡散を過大評価し A2 を plug flow と誤予測．再循環域の乱流は両者とも過大評価．
    \item 鉛直 2 次流（再循環域）は 2D で未再現＝3D 課題（Haun et al.\ 等）．
  \end{itemize}
  \item \textbf{限界・適用範囲}：一方向流・可逆ポートなし．
  \item \textbf{本研究との接点}：一番使いやすい結論は，$k$-$\varepsilon$ 付き 2D 浅水モデルは平均流と平均運動エネルギーはかなり再現できるが，再循環域の TKE には限界がある，という点．本研究に引きつけると，平均流況の比較には 2D や表面 PIV でもかなり議論できる一方，堆積・再懸濁まで踏み込むなら再循環域の乱流エネルギー評価には注意が必要，という読みになる．
\end{itemize}
\noindent（確認: p.586--597）

\noindent\litref{Jirka (2001)}\quad 浅水流の基礎概念（2DCS）を定義した総説．Jirka 自身と共同研究者（Chen, Dracos, Giger）の一連の研究の総括．
\begin{itemize}
  \item \textbf{問い}：浅水流に現れる大規模流れ構造（2DCS）はどう定義され，どの機構で生成され，混合にどう効くか．
  \item \textbf{主要知見}：
  \begin{itemize}
    \item \textbf{浅水流の定義}：水平 2 方向の広がりが水深を大きく上回る流れ．乱流条件は深さ Reynolds 数 $\mathrm{Re}_h = UH/\nu \gg 10^3$．「実流れでは容易に満たされるが\textbf{実験室では困難を生じうる}」と明記．
    \item \textbf{2DCS}：水深全体に一様に広がり位相相関した渦度をもつ大規模乱流塊．渦度ベクトルは鉛直に揃い，水平スケール $\lambda/H \gg 1$．
    \item \textbf{生成機構 3 分類}（強い順）：Type A 地形強制（島・岬・突堤による剥離）／Type B 内部横方向せん断不安定（浅水ジェット・後流・混合層）／Type C 基本流の 2 次不安定（証拠は限定的）．本研究の inflow ジェットは Type B に該当（論文自身が shallow jet を Type B の例とする）．
    \item \textbf{発達には距離が要る}：浅水ジェットは 3 領域に分かれ，遠方場は $x/H > 10$．ここで初めて渦が水深を超えて成長し 2 次元的性格を持つ．近傍場（$x/H < 1$）は平均的には 2D でも小スケールは高度に 3D．
    \item \textbf{減衰は底面摩擦}：渦が 1 回転で回転エネルギーを失う条件から最大渦サイズ $\lambda_\mathrm{max} \sim 2H/c_f$．$c_f = 0.005$（現地）--$0.01$（実験室）で $\lambda_\mathrm{max}/H = O(100)$．
    \item 遠方場では 3D 乱流（$-5/3$）から 2D 乱流（$-3$，エンストロフィーカスケード＋逆カスケード）へ遷移し，渦の合体により低周波側のピークが下流ほど増大．ジェットの $\mathrm{St} = fb/U_m \sim 0.08$--$0.10$．
    \item \textbf{施設要求}：2DCS への境界影響を避けるには basin は少なくとも幅 4~m・長さ 10~m 必要．
  \end{itemize}
  \item \textbf{限界・適用範囲}：総説で新規データなし．\textbf{単方向流が前提}で可逆流・閉鎖 basin は対象外．曲率効果なし．2001 年時点で LES は未実装（$k$-$\varepsilon$ は Type A にはまずまず，Type B には「かなり疑わしい」）．
  \item \textbf{本研究との接点}：
  \begin{itemize}
    \item 「2DCS」の概念と Type 分類の出典．本研究の inflow ジェットを Type B と位置づけられる．
    \item \textbf{$x/H > 10$} は $h = 0.03$--$0.10$~m で $x = 0.3$--$1.0$~m．port 近傍は浅水流として振る舞わず，$\Omega$ から除外する判断の物理的根拠になる．また水深掃引は「2 次元化がどこから始まるか」を連続的に変えている．
    \item \textbf{$\lambda_\mathrm{max} \sim 2H/c_f$} を本研究に当てると $h = 0.055$~m, $c_f = 0.01$ で約 11~m ＝ basin（$3 \times 2$~m）より大きい．つまり\textbf{渦サイズを決めるのは底面摩擦ではなく壁}．Jirka の非拘束浅水流との本質的な違い．
    \item \textbf{施設サイズ $4 \times 10$~m 基準を本研究（$3 \times 2$~m）は下回る}．壁の影響が本質的に含まれることを認識すべき（ただし本研究の目的は閉鎖 basin のレジームなので限界というより設定の違い）．
    \item \textbf{$\mathrm{Re}_h$ の確認が要る}：本研究の basin 内代表流速次第では $\mathrm{Re}_h \gg 10^3$ を満たさない可能性がある．
    \item \textbf{outflow は Jirka の枠組みで語れない}：吸込は運動量を注入しないため Type A/B/C のいずれでもない．inflow のみが浅水流の理論に乗る，という非対称がここにもある．
  \end{itemize}
\end{itemize}
\noindent（確認: p.567--573）

\noindent\litref{Chen \& Jirka (1997)}\quad 浅水後流の絶対不安定／対流不安定．Chen \& Jirka 1995 の実験分類を理論側から説明した対の論文．将来の安定性解析に着手する際の入口．
\begin{itemize}
  \item \textbf{問い}：浅水後流の 3 パターン（渦列／非定常バブル／定常バブル）が浅さ数 $S$ で決まるのはなぜか．
  \item \textbf{手法}：深さ平均浅水方程式に底面摩擦項を加えて線形化＝\textbf{底面摩擦つき修正 Orr--Sommerfeld 方程式}．基本流は平行流と仮定し，速度分布を $U(y) = \bar{U}(1 - R + 2R\,\mathrm{sech}^2(y/l))$ でモデル化．$R$ ＝（中心線速度と周囲流速の比）で，$R = 0$ 一様流，$R = -1$ 中心線速度ゼロ，$R < -1$ で逆流．分散関係の鞍点（群速度 $d\beta/d\alpha = 0$）で絶対／対流不安定を判定．Chebyshev 擬スペクトル法．
  \item \textbf{主要知見}：
  \begin{itemize}
    \item \textbf{絶対不安定 vs 対流不安定}：絶対＝擾乱が上流にも広がりその場に留まって成長し，外部入力なしに自励振動する（発振器）．対流＝擾乱は成長しつつ下流へ流され，その場では消える（増幅器）．渦列は絶対不安定が非線形飽和した最終産物．
    \item 底面摩擦は絶対不安定を減衰させる．$R = -1$ で臨界 $S_{bca} = 0.331$（非粘性極限）．
    \item $R$--$S_b$ 平面が\textbf{絶対不安定／対流不安定／安定}の 3 領域に分かれる安定性図（Fig.~6）．左下（速度欠損大・摩擦小）ほど不安定という直感どおりの構図．
    \item \textbf{$S$ を増やすことは $\mathrm{Re}$ を減らすことと等価}．浅水後流の $S_C > S_A > S_U > S_V$ が非拘束低 $\mathrm{Re}$ 後流の $\mathrm{Re}_K > \mathrm{Re}_{IK} > \mathrm{Re}_A > \mathrm{Re}_C$ に完全対応（Table~1）＝底面摩擦が粘性と同じ役割を果たす．
    \item \textbf{局所絶対不安定は必要条件だが十分条件ではない}（Chomaz et al.\ 1988）．絶対不安定領域が有限の臨界サイズに達して初めて大域モードが自励する．→理論の閾値（$S_A = 0.79$）と実際に渦列が見える閾値（$S_V = 0.2$）がずれる理由．
    \item 粘性は効かない（$\mathrm{Re}_b > 1000$ で流れはほぼ不変）＝摩擦パラメータ $S$ が唯一の支配因子．周波数も理論 $\mathrm{St}_d = 0.20$ が実験 $0.21$ とよく一致．
  \end{itemize}
  \item \textbf{限界・適用範囲}：\textbf{平行流の仮定}を著者自ら「結果はせいぜい近似的な現実」と明記．線形理論なので振幅は予測できない．後流が対象で噴流は Chen \& Jirka 1998，貯水池・可逆ポートは対象外．単方向流．局所解析であり大域不安定（非平行効果）は扱わない．
  \item \textbf{本研究との接点}：
  \begin{itemize}
    \item \textbf{将来の安定性解析の中核候補}：底面摩擦を含む深さ平均モデルの線形安定性解析という枠組み．
    \item 「絶対＝自励，対流＝外乱駆動」という区別は，本研究の outflow の非定常性が自励的なのか port の流量変動に応答しているだけなのかを論じる際の軸になりうる．ただし判定には基本流の速度分布と平行流の仮定が要るため，本研究のデータから直接は判定できない．概念としての参照に留める．
    \item \textbf{$S$ はそのまま使えない}：$S = c_f D/h$ の $D$ は障害物スケールで，本研究には障害物がない．port 幅や basin 幅を代入することは形式的には可能だが，物理的意味づけは別途必要．
  \end{itemize}
\end{itemize}
\noindent（確認: p.157--172）

\noindent\litref{Chen \& Jirka (1999)}\quad 浅水層に拘束された平面噴流の LIF 計測．水深 2.5--8~cm で\textbf{本研究の 3--10~cm とほぼ一致する唯一の先行実験}．
\begin{itemize}
  \item \textbf{問い}：浅水層中の平面噴流の濃度場・混合はどうなるか．浅さは何を変えるか．
  \item \textbf{手法}：$7.8 \times 6$~m 水槽，スロット幅 $B = 1$~cm，$\mathrm{Re}_B \sim 10^4$，水深 $h = 2.5/4/6/8$~cm の 4 条件．LIF（Rhodamine B，5~W Ar レーザー，CCD 30~fps）で中央水深の濃度場を計測．弱い共流（1--2~cm/s）で下流からの再循環を防止＝\textbf{閉鎖 basin とは正反対の設定}．大スケールゆえ染料によるレーザー吸収が無視できず，光線に沿って画素ごとに逐次補正（既知濃度の容器で検証，係数が濃度によらず一定）．
  \item \textbf{主要知見}：
  \begin{itemize}
    \item \textbf{$x/h \sim 10$ で蛇行開始}．根拠は単純な幾何で，ジェットの広がりが $db/dx \sim 0.1$ なので $b = h$ となるのが $x = 10h$．それより手前ではジェットは「浅い」ことを知らず，通常の 3D 平面噴流として振る舞う．
    \item \textbf{全ての濃度特性はスロット幅 $B$ ではなく水深 $h$ でスケールする}（速度場での Giger 1991 と一致）．4 つの水深のデータが $x/h$ で並べると重なる（$x/B$ では近傍場で重ならない）．
    \item \textbf{中心線の濃度変動 RMS が V 字}＝$x/h \sim 16$ で最小，その後 $X_\mathrm{max}$（16--48）で最大．3D 乱流による変動の減衰曲線と，2D 蛇行による変動の増加曲線が交差した形．谷の位置が遷移の指標になる．
    \item \textbf{平均は自己相似だが変動は非相似}．平均濃度は $x/h = 5$--$40$ でガウス分布に重なる（＝通常の噴流と同じ）のに，RMS は重ならず下流ほど増大．平均場だけ見ると 3D→2D 遷移を見落とす．
    \item RMS の横断分布は\textbf{二峰性}で，中心でも外縁でもなく半値幅の位置が最大＝ジェット内外を往復する間欠性が最大の場所（Guo 2024 の $\chi$ が再循環の「縁」で乱れ最大なのと同じ構図）．
    \item \textbf{濃度は速度より 1.7 倍広く広がる}（$b/B$ の傾き 0.17 対 速度場 0.10）．蛇行時，速度はベクトルゆえ左右の振れが平均で相殺するが，濃度はスカラーで相殺しないため．通常の平面噴流の 1.35 倍より大．
    \item \textbf{瞬時濃度は同位置の平均の 4 倍}．2D 乱流は渦伸長がなく渦が砕けないため濃い塊が保存される（「鉛直方向の制約からより弱い希釈しか起こらない」）．平均濃度で汚染を評価すると過小評価になる．
    \item スペクトルは $x/h = 33.6$ の $-5/3$ から $92.8$ の $-3$ へ遷移．自己相関は下流ほど周期が長い（渦の合体）．
  \end{itemize}
  \item \textbf{限界・適用範囲}：濃度場のみで速度場ではない（速度は Dracos 1992 / Giger 1991 の LDV）．単方向噴流＋共流で閉鎖 basin ではない．近傍場のスペクトルは CCD 30~fps の制約で測れないと著者が明示．高 Schmidt 数（$\sim 1000$）で微細スケールは追随しない．
  \item \textbf{本研究との接点}：
  \begin{itemize}
    \item 「浅水層に拘束された平面噴流」＝本研究の inflow ジェットに最も近い基礎設定で，水深範囲も一致．
    \item \textbf{$x/h \sim 10$} は $h = 0.055$~m で $x \sim 0.55$~m．port から半メートルは「浅水」として振る舞わない．$\Omega$ から port 近傍を除外する判断，および水深掃引が「2 次元化の開始位置」を連続的に変えていることの根拠．
    \item \textbf{平均は相似・変動は非相似}は $I_\mathrm{unst}$ を primary metric に置く直接の根拠．平均ベクトル場と streamline だけでレジーム判定すると遷移を見落としうる．
    \item \textbf{変動は 2 状態の境界で最大}は本研究の空間分布（$\phi_\mathrm{lv}$ の境界付近など）を解釈する視点．
    \item \textbf{2D 化で混合が悪化し濃い塊が保存される}は将来の土砂実験に効く．「低流速域に土砂が溜まる」という平均場ベースの推論は，瞬時の濃い塊による輸送を見落とす可能性がある（→§6）．
    \item ただし彼らは共流で再循環を防いでおり，壁に当たって戻る本研究の basin では遠方場の性質が変わるはず．スカラー場と速度場の統計が一致する必要もない（Brevis 2011 の未解決点と同じ論点）．
  \end{itemize}
\end{itemize}
\noindent（確認: p.817--826）

\noindent\litref{Giger et al. (1991)}\quad（補助）浅水平面噴流の速度場（LDA）．Dracos 1992 の姉妹論文で，Chen \& Jirka 1999（濃度場）と Jirka 2001 総説の一次データ（全 29~p 通読）．
\begin{itemize}
  \item \textbf{問い}：浅水中の平面噴流はどう周囲流体を巻き込み混合するか．水深はそれをどう変えるか．
  \item \textbf{手法}：water table（$1.80 \times 2.80$~m）に $B = 10$~mm のスロットから $U_0 = 1.10$~m/s，$\mathrm{Re} \sim 10^4$．水深を $H/B = 4$--$36$ に変化．LDA 2 成分．側方給水で無限水域を近似（Taylor 1958 の解と照合）．なお大半の実験は\textbf{上面もガラス板}（両面 no-slip，$n = 2$）．自由水面は $n = 1$．
  \item \textbf{主要知見}：
  \begin{itemize}
    \item \textbf{3 領域分割}：近傍場 $x/H < 2$（通常の平面噴流）／中間場 $2 < x/H < 10$（二次流が支配し平均速度が 3 次元になる）／遠方場 $x/H > 10$（平均場が再び 2 次元）．分割の根拠は\textbf{二次流の有無}であって，Jirka 2001 の「2 次元乱流への遷移」とは着目点が異なる．同じ $x/H = 10$ でも意味が違う．
    \item 中間場では 4 象限に螺旋セルが立ち，中心線速度が最大 16~\% 低下（$x/H \sim 6$ で最強，$> 10$ で消滅）．ただしこの運動量フラックス減少は\textbf{見かけのもの}で，境界面近くへの再配分にすぎない．
    \item \textbf{底面摩擦は広がりに影響せず，中心線速度だけを下げる}．広がり係数 $C_b = 0.10$ は摩擦に依らず古典値と同等．従来モデル（Lee \& Jirka 1980 ら，摩擦が巻込み係数を変えないと仮定）への反論．
    \item \textbf{巻込みは下流で減少し有限距離でゼロになる}（式 41, 42）．摩擦なしなら $\alpha = 0.054$ で一定だが，摩擦があると単調減少．→それ以上希釈されない．最大希釈率は\textbf{アスペクト比 $H/B$ と摩擦係数 $f$ だけで決まる}（式 43, Fig.~16）．$H/B$ が小さいほど希釈は頭打ちになる．
    \item \textbf{遠方場では混合効率が落ちる}．噴流内の乱流部分は近傍場 95~\%，中間場 99~\%（二次流が促進），$x/H \sim 40$ で 80~\%．取り込んだ流体が混ざらないまま運ばれる．
    \item 遠方場の乱れは非等方（$u$ 成分は 0.2 で一定，$v$ 成分は増加し続ける）＝蛇行が横変動を生み続ける．
    \item 運動量フラックスが保存しない問題は\textbf{未解決}．Appendix~A の 27 研究で $M_\infty/M_0 = 0.52$--$1.77$ とばらつく．圧力勾配は実測で原因から排除（Fig.~11）．物理的理由は不明と著者が結論．
  \end{itemize}
  \item \textbf{限界・適用範囲}：単方向・非拘束の自由噴流で，側方給水により無限水域を近似＝閉じ込め貯水池ではない．摩擦係数 $f = 0.043$ は Fig.~13 の最良フィットで決めた値であり独立測定ではない．
  \item \textbf{本研究との接点}：
  \begin{itemize}
    \item \textbf{$x/H \sim 10$ の一次的な出典}．Fig.~5a（$x/B$ で散らばる）と Fig.~5b（$x/H$ で揃う）の対比が「水深がスケーリング変数」の根拠．本研究が $h$ を軸に整理する方針を支える．ただしこの系譜は「開口寸法は忘れて水深で整理せよ」という立場なので，\textbf{$a$（開口高さ）が効くかは本研究が示すべきこと}．
    \item \textbf{条件の違いが 3 つあり，そのまま適用できない}：(i) 上面ガラス板 $n = 2$ に対し本研究は自由水面 $n = 1$ で摩擦の効きが半分．(ii) 側方給水で無限水域を近似しており，\textbf{壁で囲まれた本研究とは正反対}．(iii) $H/B = 4$--$36$ に対し本研究の port は幅 0.11~m・水深 0.03--0.10~m で相当量は 0.3--0.9 と彼らの下限を大きく下回る．→Fig.~16 に従えば本研究の条件は\textbf{希釈が最も頭打ちになる側}．
    \item \textbf{巻込みが止まる距離}（式 42）は本研究でも計算可能（$n = 1$，$f$ を仮定）．ジェットがどこまで周囲を取り込み続けるかは basin 内循環の駆動力に関わるので，検討の価値がある．
    \item \textbf{遠方場で混合効率が落ちる}は Chen \& Jirka 1999 の「瞬時濃度が平均の 4 倍」と同一現象の速度側からの記述．将来の土砂実験で，basin 遠方部では取り込まれた流体が十分混ざらない可能性を示唆．
    \item 本研究では $x/H \sim 10$（$h = 0.055$~m で $x \sim 0.55$~m）に達した後，ジェットは必ず壁に当たって反射する．遠方場の性質は彼らの非拘束条件とは異なるはず．
    \item Dufresne 2010 が引用しており浅水貯水池の文脈で既に参照されている基礎（Dufresne 側は未確認）．
  \end{itemize}
\end{itemize}
\noindent（確認: p.615--642）

\noindent\litref{Heller (2011)}\quad 物理模型のスケール効果の標準的な総説（Forum paper，全 14~p 通読）．\texttt{methods.tex} の「Reynolds 数を含む全ての無次元量を実機と同時に一致させることは意図していない」および Limitation の外挿不可の主張は，本プロジェクトの中心的な方法論的立場でありながら現在\textbf{無引用}．その出典．
\begin{itemize}
  \item \textbf{問い}：物理水理模型のスケール効果はどう生じ，どの条件下で無視でき，どう回避・補償・補正するか．
  \item \textbf{手法}：総説．相似の 3 階層（幾何・運動学・力学）と 5 つの力比（$\mathsf{F}, \mathsf{R}, \mathsf{W}, \mathsf{C}, \mathsf{E}$）を整理し，相似を得る 4 手法（inspectional analysis／dimensional analysis／calibration／scale series）を比較．地すべり誘起段波を例に全手法を適用し，回避・補償・補正の実務を示す．
  \item \textbf{主要知見}：
  \begin{itemize}
    \item \textbf{$\lambda \neq 1$ なら必ずスケール効果が生じる}．同一流体（水）を使う限り力比のうち 1 つしか一致させられない（式 13）．「最も関連する力比を選び他は無視できる程度に抑える」しかない＝Froude 限定相似の正当化そのもの．
    \item \textbf{$\lambda$ 単独は判定基準にならない}．$\lambda = 2$ でも空気連行に有意な差が出る例（Fig.~2）がある一方，大きい $\lambda$ でも効かない例もある．「$\lambda$ だけを限界基準にするのは不十分」と明記．→本研究が「1:X 模型」と書かず prototype-informed と位置づける方針の裏づけ．
    \item \textbf{判定はパラメータごとに行う}：同じ模型でも流量には効かないが空気濃度には効く（Fig.~1）といったことが起きる．「各パラメータについて個別の判断が要る」．
    \item \textbf{スケール効果は damping 側に働く}：模型では粘性・表面張力が相対的に強いので，相対波高・相対流量・輸送土砂量などは\textbf{実機より小さく出る}のが通例．予測が過大か過小かの判断はつく．
    \item \textbf{Table~1 が実務の核心}：24 行の現象別「significant scale effects を避ける限界基準」．本研究に近いのは，Air-entraining free vortex at horizontal intake（Anwar et al.\ 1978）の $Q_i/(z_i\nu) > 30{,}000$ かつ $\rho Q_i^2 z_i/(A_i^2\sigma) > 10{,}000$，および Water wave, theoretical effect of surface tension（Hughes 1993）の $T > 0.35$~s かつ $h > 0.02$~m で表面張力の寄与が 1~\% 未満．→本研究の下限水深 0.03~m はこれを満たす．
    \item Table~2 は現象別の典型 scale（intake は 1:50--1:100）．ただし「経済性との妥協でありスケール効果が無視できることを意味しない」と明記．
    \item 移動床と distensible structure はスケール効果なしの模型化がとくに難しい，が主要結論の一つ．粒径をそのまま $\lambda$ で縮小すると $d_g < 0.22$~mm で付着力が効き，$d_{50} > 0.5$~mm が要るなどの制約．
  \end{itemize}
  \item \textbf{限界・適用範囲}：総説で新規データなし．例は段波・堰・跳水・洗掘など古典的構造物が中心で，浅水貯水池の平面流動レジームは Table~1 に項目がない．可逆流・閉鎖 basin も扱わない．Table~1 の基準は「その現象・その prototype feature に対して」定義されたもので，他の量への流用を著者が繰り返し警告．
  \item \textbf{本研究との接点}：
  \begin{itemize}
    \item \textbf{Froude 限定相似の立場の出典}：\texttt{methods.tex} に直接引ける．同時に「何が犠牲になっているか（$\mathsf{R}$, $\mathsf{W}$ の不一致）」も本論文の枠組みで書ける．
    \item \textbf{primary metrics は界面現象を含まない}：$E$, $\phi_\mathrm{lv}$, $I_\mathrm{asym}$, $I_\mathrm{circ}$, $I_\mathrm{unst}$ はいずれも速度場から作る量で，空気連行や砕波のような $\mathsf{W}$ が効く現象ではない．上記「パラメータごとに判断」に従えば，これらは Froude 相似で扱える側だと主張できる．逆に将来の土砂実験では本論文の警告がそのまま効く（→§6）．
    \item \textbf{下限水深 0.03~m}：Table~1 の $h > 0.02$~m は満たすが，これは\textbf{波長に対する表面張力の寄与}の基準で，本研究の定常流とは現象が違う．援用するなら「桁の目安として」と限定すること．なお表面膜による表面層の切り離しは scale effect ではなく model/measurement effect で本論文の範囲外．
    \item \textbf{damping 側という指摘}は Limitations の書き方を決める：観察される非定常性や混合は，実機ではこれ以上に強く出る可能性がある（過小評価側）と述べられる．
    \item Table~1 の取水口渦の基準は M\"oller 2015 と接続．あちらが具体，こちらが一般論．
  \end{itemize}
\end{itemize}
\noindent（確認: p.293--306，全 14~p 通読）

\subsection*{論文横断整理}

\subsection*{原稿転記メモ}

% ================================================================
\section{土砂移動・CFD 連携}
% ================================================================

\subsection*{主要な問い}

\subsection*{論文メモ}

\noindent\litref{Camnasio et al. (2013)}\quad 入口・出口配置が流況と土砂堆積に与える影響，および堆積から流況へのフィードバックを実験＋2D 数値で解明（§6 の核）．
\begin{itemize}
  \item \textbf{問い}：入口・出口チャネルの位置（対称/非対称配置）は流況と土砂堆積パターンにどう影響するか．堆積は流況にフィードバックするか．
  \item \textbf{手法}：EPFL LCH 矩形浅水貯水池（$L=4.5$~m, $B=4$~m, $b=0.25$~m, $h=0.2$~m, $Q=7$~L/s, $F_\mathrm{in}=0.1$, $R_\mathrm{in}=112000$）．入口/出口の 4 配置 C-C/L-L/L-R/C-R．浮遊砂はクルミ殻（$d_{50}=89~\mu$m, $C_\mathrm{in}\approx 2$~g/L, 4 時間）．UVP で速度場，レーザで堆積厚を測定．WOLF 2D（深さ平均, 2 スケール $k$-$\varepsilon$）に，実測堆積厚を時間変化地形として与え，流況変化を再現．幾何縮尺 1:50 で Rhône 沿い実貯水池を念頭に設定．
  \item \textbf{主要知見}：
  \begin{itemize}
    \item 堆積は主ジェット経路に沿って最大（供給量が輸送能力を上回るため），再循環中心では $<5$~mm．Stovin/Dufresne 型の「供給<輸送能で渦核に堆積」とは異なる条件．
    \item L-R 配置では，清水時は左側壁への再付着流だが，土砂供給後に再付着なし流況へ完全に変化．堆積厚が 3--4~mm に達する約 30 分後に遷移し，流れと土砂の双方向カップリングを示す．
    \item WOLF 2D では，L-R 配置で初期条件により再付着あり/なしの 2 つの定常解が得られ，双安定性が示された．
    \item 堆積地形を時間変化底面として与えると L-R の流況遷移を再現できる．一方，底面・壁面粗度のみの変更では流況はほぼ不変．$S=c_f(B-b)/(4h)\approx 0.02$--$0.04$ で非摩擦的条件にあるため，粗度効果は小さい．
    \item WOLF 2D は速度プロファイルを良好に再現し，堆積地形による流況切替も再現したが，遷移時刻は実験より遅く評価した．
  \end{itemize}
  \item \textbf{限界・適用範囲}：一方向流（流入→流出）であり，単一可逆ポートではない．供給量>輸送能力の特定条件であり，低供給量・自然砂条件とは堆積位置が異なりうる．深さ平均 2D モデルであり，3D 二次流や鉛直構造は明示的には扱わない．
  \item \textbf{本研究との接点}：流れと土砂の双方向カップリングと「堆積が流況モードを切り替える」現象を示す中核文献．入口/出口の非対称配置（C-C/L-L/L-R/C-R）は，本研究のポート配置・可逆性を考える近い参照枠．WOLF 2D＋時間変化地形の手法は §6 CFD の参照になる．
\end{itemize}
\noindent（確認: p.535--547）

\noindent\litref{Camnasio et al. (2012)}\quad（補助）Camnasio 2013 JHR の会議版コンパニオン（ICSE6, Paris 2012．同一実験・同一 4 配置）．
\begin{itemize}
  \item 問い・手法・主要知見は Camnasio (2013) とほぼ同じ．
  \item 強調点：堆積から流況へのフィードバック機構のうち「乱流減衰（turbulence damping）」を主因と推定．浮遊砂が乱流を減衰させ（Nino \& Garcia 1998 の粒子レイノルズ閾値 $<2$），L-R 配置の流況遷移を説明．
  \item \textbf{本研究との接点}：乱流減衰フィードバックの根拠として Camnasio 2013 と併せて引用．
\end{itemize}
\noindent（確認: 会議論文・全体）

\noindent\litref{Esmaeili et al. (2014)}\quad 3D 数値モデル（SSIIM）で浅水貯水池の流況と 3D 二次流を再現し，排砂運用への適用を検討．
\begin{itemize}
  \item \textbf{問い}：3D 数値モデルは浅水貯水池の流況（対称/非対称）と 3D 二次流をどこまで再現できるか．土砂・排砂運用にどう活かせるか．
  \item \textbf{手法}：完全 3D RANS モデル SSIIM（FVM, $k$-$\varepsilon$）．EPFL 施設の 4 ケース T8/T9/T11/T13（$L,B$ を変え ER=$B/b$, AR=$L/B$, SF を振る）．固定床＋排砂後の平衡変形床．Fr=0.1, Re=$1.75\times10^4$．LSPIV（表面）＋UVP（全層 3D）で検証．土砂：浮遊（移流拡散）＋掃流（Van Rijn）．
  \item \textbf{主要知見}：
  \begin{itemize}
    \item 3D モデルは対称・非対称流況を再現（非対称は入口 2.5~\% 擾乱で誘発）．高 SF（T11,T13）は擾乱に鈍感で対称，低 SF（T8,T9）は非対称．Dufresne の SF 基準と整合．
    \item 2D 深さ平均では扱えない鉛直二次流・鉛直速度を捕捉．土砂堆積帯の予測に重要．
    \item LSPIV 表面＋UVP 全層プロファイルとよく一致．
    \item 排砂後の平衡変形床上の流況も計算（drawdown flushing 運用，Dashidaira/黒部川を想定）．
  \end{itemize}
  \item \textbf{限界・適用範囲}：一方向流で可逆ポートではない．SSIIM 特有のモデル仮定．
  \item \textbf{本研究との接点}：3D CFD（SSIIM）で 3D 二次流を捉える＝Kantoush 2008a が指摘した壁近傍 2 次流を陽に扱う道．排砂・flushing 運用の文脈．表面 LSPIV と全層の対応づけ．
\end{itemize}
\noindent（確認: ページ未確認）

\noindent\litref{Dufresne et al. (2010, IJSR)}\quad 流況（S0/A1/A2）が土砂堆積分布とトラップ効率をどう決めるかを実験で体系化（§2 の JHR「流況分類」と同一実験キャンペーンの土砂側）．
\begin{itemize}
  \item \textbf{問い}：流況（S0/A1/A2）は土砂堆積分布とトラップ効率にどう影響するか．
  \item \textbf{手法}：Liège 10.4~m 水路（$\Delta B = 0.35$~m，$b = 0.285$~m，3 池長 $L = 1.8/2.2/7.0$~m）．流れは染料可視化で再付着長（中央値＋自然変動）．土砂は粒状プラスチック（Styrolux 656C，$\rho = 1020$~kg/m$^3$，$d = 2.4$~mm，沈降 25~mm/s，ストームウォーターの粗粒分 $d > d_{90}$ を模擬），$C_\mathrm{in} = 0.5$~g/L，10 分注入，出口ネットで 3 時間帯のトラップ効率 $\eta$ を計測．
  \item \textbf{主要知見}：
  \begin{itemize}
    \item 流況は S0/A1/A2 と，遷移域の A1-S0 不安定流に整理される．土砂試験は主に安定流況 S0/A1/A2 について実施．
    \item 再付着長は時間変動をもつ．特に第 2 再付着長 $R_2$ は変動が大きく，A2 は瞬時的には A1 的状態との間で揺らぐ．
    \item 堆積は流況に強く依存：
    \begin{itemize}
      \item S0 ＝ 入口両隅にほぼ対称に堆積し，捕捉効率は低い．
      \item A1 ＝ 両入口隅と大循環核に堆積し，循環核に粒子が捕捉されるため局所堆積が大きい．
      \item A2 ＝ 一方の入口隅と下流域に広く堆積し，A1 より分散的．
    \end{itemize}
    \item トラップ効率は S0 ＝ 0--20~\%（平均約 6~\%），A1 ＝ 10--40~\%（平均約 28~\%），A2 ＝ 30--60~\%（平均約 48~\%）．特に S0 から A1 への遷移で効率が急増する．
    \item A1 は池長あたりの捕捉ポテンシャルが最大．そのため，長池を A1 流況が成立する長さの複数池に分割すると，総捕捉効率を高められる可能性がある．
  \end{itemize}
  \item \textbf{限界・適用範囲}：一方向流・可逆ポートなし．プラスチック粒子，少数の土砂試験．
  \item \textbf{本研究との接点}：流況モードが堆積パターン・トラップ効率を決める核心．§2 の流況分類（同一 2010 キャンペーンの JHR 版）と対をなす土砂側．可逆ポートで流入／流出モードが変われば堆積分布も変わる基盤．
\end{itemize}
\noindent（確認: p.258--270）

\noindent\litref{Dufresne et al. (2012)}\quad（補助）Dufresne 2010 IJSR と同一実験の実務向け再整理．形状パラメータ SP と設計指針を前面に出す．
\begin{itemize}
  \item \textbf{問い}：矩形浅水池の流況パターンは，堆積分布・捕捉効率・設計指針にどう結びつくか．
  \item \textbf{主要知見}：
  \begin{itemize}
    \item 形状パラメータ $\mathrm{SP} = L/(\Delta B^{0.6} b^{0.4})$．$\mathrm{SP} < 6.2$ 対称（堆積抑制），$\mathrm{SP} > 6.8$ 非対称（堆積促進）．
    \item 堆積分布・捕捉効率の流況依存は Dufresne 2010 IJSR と同じ（S0 ＝ 0--20~\%，A1 ＝ 10--40~\%，A2 ＝ 30--60~\%）．池長あたりの捕捉ポテンシャルは A1 が最大．
    \item 設計指針：堆積最大化＝池を A1 流況が成立する長さに分割，最小化＝S0 になる短池化．
  \end{itemize}
  \item \textbf{本研究との接点}：Dufresne 2010 IJSR と併せて引用．SP による堆積制御の実務指針．
\end{itemize}
\noindent（確認: p.504--510）

\noindent\litref{Dewals et al. (2024)}\quad（補助）流れ場の診断指標として「体積平均 water age 分布」を提案し，浅水池の混合・短絡構造を低次元に表現．
\begin{itemize}
  \item \textbf{問い}：体積平均 water-age 分布関数 $\mu(\tau)$ は，浅水池の流れ場について何を明らかにするか．
  \item \textbf{手法}：Camnasio et al.\ 2014 が 2D 浅水方程式＋深さ平均 $k$-$\varepsilon$ で計算した 10 形状（$L = 3$--$6$~m，$B = 0.6$--$4$~m）の流れ場と，Dewals et al.\ 2020 の位置依存 water-age 分布 $c(x,y,\tau)$ を用い，体積平均年齢分布 $\mu(\tau)$ を解析．完全混合槽・plug flow と比較し，特性時間 $T_m = V/Q$，$T_j = L/U$ で整理．
  \item \textbf{主要知見}：
  \begin{itemize}
    \item $\mu(\tau)$ は，小 $\tau$ で一定（速い経路＝主ジェットの plug-flow 的通過），その後ステップ列（再循環域に捕捉された水粒子の遅延経路），最終的に指数減衰（古い水粒子については完全混合槽的挙動）を示す．
    \item 平坦部の上限 $\tau^*$ は $T_j$ と同程度．直進ジェットの短池（$L/B < 1.5$）では $\tau^* \approx T_j$．再付着ジェットや第 2 再付着を伴うケースでは，ジェットの曲がり・減速・拡散により $\tau^*/T_j$ が大きくなる．
    \item 体積平均年齢分布は，位置依存 water-age 分布や速度場に比べて次元を大幅に削減しつつ，速い経路・遅い経路・長期滞留域という混合・短絡情報を保持する．
    \item 主ジェット経路と，再循環域に $n$ 回捕捉される遅延経路の重ね合わせとして，$\mu(\tau)$ のステップ状分布を説明する簡易概念モデルを提示．
  \end{itemize}
  \item \textbf{限界・適用範囲}：新規実験ではなく既往の 2D 数値流れ場と water-age 計算の再解析．定常・一方向流で可逆ポートや非定常運転は扱わない．2D 深さ平均で成層・鉛直循環・3D 構造は明示的に扱わない．土砂年齢や粒子沈降は扱わず water age のみ．
  \item \textbf{本研究との接点}：流況モードが短絡・滞留・混合効率をどう変えるかを低次元に診断する枠組み．トラップ効率そのものではないが，それを左右する流れ場・滞留構造の定量指標候補．可逆ポートで流入／流出モードが変わる場合，速度場比較に加えて water-age 分布や粒子 age 分布でモード差を評価できる可能性がある．
\end{itemize}
\noindent（確認: p.75--93）

\subsection*{論文横断整理}

\subsection*{原稿転記メモ}

% ================================================================
\section{流れの安定性と分岐（手法横断メモ）}
% ================================================================

\noindent 本節は\textbf{索引}である．線形安定解析・固有値解析・分岐論の枠組みで対称流の不安定化やモード選択を論じた文献を，将来の安定性解析（ODE ベース）に向けてここに集約する．
各文献の本体メモは，現象の文脈が失われないよう\textbf{それぞれの節に置いたまま}とし，ここには手法・枠組みの観点からの要約と参照先のみを記す．
この節を主戦場とする文献（Chomaz, Huerre \& Monkewitz, Fearn/Cherdron など）を今後入手した場合は，本体メモを下記「論文メモ」に置く．

\subsection*{論文メモ}

\subsection*{論文横断整理}

\noindent\textbf{(1) 三つの異なる枠組みが同じ問いに答えている}

いずれも「対称な設定でなぜ非対称流が生じるか」「いつ流れは自励的に振動するか」を扱うが，出発点が異なる．

\begin{itemize}
  \item \textbf{固有値問題として解く}（Dewals 2008, Chen \& Jirka 1997）— 基本流を仮定し，微小擾乱の成長率を分散関係から求める．深さ平均モデルに底面摩擦を加えた修正 Orr--Sommerfeld 型が浅水流での標準形．
  \item \textbf{分岐構造として解く}（Mullin 2003）— 制御パラメータに対する解の分岐図を追跡する．対称性破壊ピッチフォーク，余次元 2 点，欠陥による分断．
  \item \textbf{変分原理から選ぶ}（Westhoff 2018）— 安定性を解かず，複数の候補状態のうち散逸が最大のものが実現すると仮定する．上の 2 つとは独立の第三の説明原理．
\end{itemize}

\medskip
\noindent\textbf{(2) 文献一覧（本体メモの所在つき）}

\begin{itemize}
  \item \textbf{Dewals et al.\ (2008)}［本体は §3］— 底面摩擦を含む修正 Orr--Sommerfeld 型の固有値問題で対称流の不安定化を評価．test 8 は擾乱減衰が test 1 の 3 倍以上強く，最大増幅率が約 1/20 で，はるかに安定．分岐パラメータ $R = \mathrm{ER} \cdot \mathrm{AR}^{8/9}$．\textbf{浅水矩形池そのものを扱った唯一の安定性解析}．
  \item \textbf{Chen \& Jirka (1997)}［本体は §5］— 浅水後流の\textbf{絶対不安定／対流不安定}．「絶対＝外部入力なしに自励振動する発振器，対流＝擾乱が下流へ流される増幅器」という区別．$R$--$S_b$ 平面の 3 領域安定性図．底面摩擦は絶対不安定を減衰させ，$S$ を増やすことは $\mathrm{Re}$ を減らすことと等価．\textbf{局所絶対不安定は必要条件だが十分条件ではない}（Chomaz et al.\ 1988）．
  \item \textbf{Mullin et al.\ (2003)}［本体は §3］— 対称急拡＋下流収縮水路のピッチフォーク分岐を有限要素分岐解析（ENTWIFE）で精査．\textbf{下流収縮（出口）による対称流の再安定化}，$(\alpha, \gamma)$ 平面の余次元 2 点，欠陥への高感受性．閉じ込め層流系だが ODE/分岐論の古典的枠組み．
  \item \textbf{Westhoff et al.\ (2018)}［本体は §3］— せん断層の最大エネルギー散逸によるモード選択．短池は対称流，長池は非対称流の方が散逸が大きく，観測される遷移と定性的に一致．ただし\textbf{なぜ系がその状態を選ぶかの決定的説明はまだない}と著者も認める．
\end{itemize}

\medskip
\noindent\textbf{(3) 本研究への適用可能性と，そのままでは使えない点}

\begin{itemize}
  \item \textbf{$S$ はそのまま持ち込めない}（Chen \& Jirka 1997）：$S = c_f D/h$ の $D$ は障害物スケールだが，本研究には障害物がない．port 幅や basin 幅を代入することは形式的には可能だが，物理的意味づけは別途必要．
  \item \textbf{平行流の仮定}が浅水安定性解析の標準的な前提だが（著者自ら「結果はせいぜい近似的な現実」と明記），閉鎖 basin の循環流はこの仮定から遠い．
  \item \textbf{自励か外乱駆動か}という軸は，本研究の outflow の非定常性を論じる際に有効．ただし判定には基本流の速度分布が要るため，現状のデータから直接は判定できない．概念としての参照に留める．
  \item \textbf{対称条件を課した数値解析は対称解しか返さない}（Constantinescu 1998［本体は §1］，Kantoush 2008a［§4］，Dewals 2008）．安定性を数値で調べる際に結論を先取りしていないかを常に疑う．
\end{itemize}

\medskip
\noindent\textbf{(4) 不足している系譜（将来の収集先）}

現状は浅水流側（Dewals, Chen \& Jirka）と閉じ込め層流側（Mullin）の 2 点しかなく，両者を繋ぐ理論的基盤が欠けている．具体的には：

\begin{itemize}
  \item 絶対／対流不安定と大域モードの理論（Chomaz et al.\ 1988; Huerre \& Monkewitz 1990）— Chen \& Jirka が依拠する枠組みの原典で，未入手．
  \item 突然拡大の対称性破壊分岐の系譜（Cherdron et al.\ 1978; Fearn et al.\ 1990; Durst; Mizushima）— Mullin が拠って立つ系譜で，未入手．
  \item 不完全分岐における欠陥による枝選択とヒステリシス — 「フェアなコイン投げではなく系統的バイアスをもつ」という主張の理論的裏づけ．
\end{itemize}

\subsection*{原稿転記メモ}

% ================================================================
\section{CFD 手法（2 本目・3 本目の材料）}
% ================================================================

\noindent 本節は\textbf{2 本目（3D VOF）・3 本目（2D/3D 比較）のための手法文献}を集める．
1 本目（実験）には CFD の数値を入れない方針（\texttt{paper/README.md}）なので，本節の文献は 1 本目では引かない．
収集の優先順位は「CFD の計画判断に効くか」で決めており，執筆時に引くだけの文献（V\&V の作法，分岐論の古典）は後回しとする．

\subsection*{論文メモ}

\noindent\litref{Komen et al. (2017)}\quad 本節の最重要文献．OpenFOAM を名指しした数値散逸の定量法（全 31~p 通読）．
\begin{itemize}
  \item \textbf{なぜ読むか}：\texttt{cfd/campaign.md} の診断「系が分岐の臨界より下にいる」は現状\textbf{定性的な推論}．laminar（陽な SGS なし）で回している以上，実効的な散逸は数値散逸であり，その大きさが不明．本論文の手法で $\nu_\mathrm{num}$ を測れれば，実験と CFD を Chen \& Jirka の同じ $S$ 軸上に並べられ，「どれだけ下にいるか」「あとどれだけ散逸を減らせば臨界を越えるか」が数値で言える．本研究の no-model 運転は本論文の ``no-model LES (UDNS)'' と同一設定であり，対応が直接つく．
  \item \textbf{問い}：産業用 LES（非スタガード有限体積・2 次精度・implicit top-hat filter．OpenFOAM, STAR-CCM+, FLUENT, Code\_Saturne を名指し）の数値散逸を物理空間でどう定量するか．SGS の寄与との比は．
  \item \textbf{手法}：\textbf{乱流運動エネルギー $k$ の収支残差から測る}．統計的定常では $k$ の収支（生成＋粘性散逸＋乱流輸送＋圧力拡散＋分子拡散）の和がゼロになるはずだが，数値誤差で非ゼロ残差が残る．これを $\epsilon_\mathrm{num}$ と定義（式 14, 15）．数値粘性は $\nu_\mathrm{num} = \nu\,\epsilon_\mathrm{num}/\epsilon$（式 16）．Schranner 2015 との違いは 2 点：(i) あちらは\textbf{全運動エネルギー}，本論文は\textbf{乱流運動エネルギー}で小スケールの寄与を強調でき，収支項の分析から傾向の\textbf{機構}も説明できる．(ii) Schranner はセル $m$ の制御体積と時間刻み $\Delta t$ で数値積分するため\textbf{瞬時・局所評価が可能}，本論文は統計平均を要する．検証は $\mathrm{Re}_\tau = 180$--$950$ のチャネル流．格子は $\Delta x^+$ で命名（90/60/30-Grid，q-DNS 用 9-Grid ＝ 265 万セル）．SGS は Smagorinsky・動的 Smagorinsky・no-model の 3 通り．統計は 60--80 FTT（TKE 収支は 180 FTT）．
  \item \textbf{主要知見}：
  \begin{itemize}
    \item \textbf{q-DNS でも数値散逸は残る}：9-Grid で総散逸が総生成より約 8~\% 小さく，平均 $\nu_\mathrm{num}/\nu \approx 0.08$．参照の擬スペクトル DNS は収支和が実質ゼロ（$-0.1$--$1.0$~\%）＝この 8~\% は有限体積法に固有．
    \item \textbf{数値散逸は SGS 散逸より 2.4--3.5 倍大きい}：no-model LES（$\mathrm{Re}_\tau = 180$）の平均で $\nu_\mathrm{num}/\nu = 0.35$（30-Grid）／$0.60$（60-Grid）に対し $\nu_\mathrm{SGS}/\nu = 0.10$／$0.25$．原文は ``substantially larger'' とし，``This explains why the addition of an explicit SGS model is not effective'' と結論．実際 LES モデルありと no-model の差は\textbf{数 \% 未満}．Fig.~14(b) の局所値は 60-Grid のチャネル中央で $\nu_\mathrm{num}/\nu \approx 1.05$ ＝数値粘性が分子粘性を超える．
    \item \textbf{誤差の署名（4 点セット）}：格子を粗くすると，平均速度が $y^+ > 10$ で過大／$u_\mathrm{rms}$（主流方向）のピークが過大／$v_\mathrm{rms}$・$w_\mathrm{rms}$（横・鉛直）がほぼ全域で過小／$k$ が過大．＝この組み合わせが「数値散逸が効いている」ことの指紋になる．
    \item \textbf{機構（本論文の独自性）}：チャネル流では $\overline{u'u'}$ のみが生成を持ち，pressure strain rate 項が $\overline{u'u'} \to \overline{v'v'}, \overline{w'w'}$ へ再分配する．格子を粗くすると\textbf{3 つの pressure strain 項すべてが過小評価}され，再分配が不足→$\overline{u'u'}$ の正味生成が増え（主流方向 RMS 過大），$\overline{v'v'}$・$\overline{w'w'}$ が受け取れない（横方向 RMS 過小）．また散逸率の正確な計算は生成率より高い解像度を要する（変動の空間勾配のため）．
    \item \textbf{数値散逸が増えると $k$ は「増える」}：より散逸的な Van Leer（TVD 混合＝中心差分に風上を混ぜたもの）を使うと，中心差分より平均速度・$u_\mathrm{rms}$ の過大評価が拡大し $k$ はさらに高くなる（Fig.~12）．＝粗格子で $k$ が人工的に増えるのは中心差分の数値振動のせいではなく，数値散逸そのものが原因．安定化のため風上を混ぜると事態は悪化する．
    \item \textbf{CFL が効く}：最大 CFL を $0.4 \to 0.2$ にすると最大数値散逸率が約 1.4 分の 1（Fig.~16）．
    \item \textbf{Rhie--Chow 補正も数値散逸源ではないか}と著者が推測．非スタガード格子の圧力振動抑制に使うが，本論文では分離できなかったと明記．＝OpenFOAM 固有の未解明因子．
    \item \textbf{アスペクト比は主因ではない}：$\Delta x^+/\Delta z^+$ を 2 倍・半分に振っても影響小．(90,15) 格子が改善するのは\textbf{スパン方向解像度}が上がり近傍壁ストリークを解像できるため（Fig.~17）．
    \item 粗い格子は同じ $\mathrm{Re}_\tau$ を得るのに大きい流量を要する（Table~4：$\mathrm{Re}_\mathrm{bulk}$ が DNS 5650 に対し 30-Grid 5643，60-Grid 6205，90-Grid 6723）＝平均速度過大評価の直接の理由．
  \end{itemize}
  \item \textbf{限界・適用範囲}：\textbf{完全発達乱流チャネル流・統計的定常・周期境界}が前提．式 (15) が使えるのは $R_k = C_k = 0$ のときのみで，チャネル流は平行流ゆえ $k$ の平均対流がゼロ．統計収束に 180 FTT を要する．単相（VOF ではない），自由表面も剥離もない．$\mathrm{Re}_\tau = 180$--$950$．OpenFOAM 2.2.0（2013 年）．Rhie--Chow の寄与は未分離．
  \item \textbf{本研究との接点}：
  \begin{itemize}
    \item \textbf{適用の可否が最大の論点}：本研究は $C_k \neq 0$（ジェットが $k$ を輸送する）ため式 (15) をそのまま使えない．式 (14) に戻り\textbf{平均対流項 $C_k$ を明示的に計算する}必要がある．手間は増えるが原理的に可能．非定常性については，各バンド内で準定常とみなせる区間を切り出す近似が要検討．
    \item \textbf{より軽い代替}：Schranner 2015 はセル単位・時間刻み単位で局所・瞬時に評価でき，統計平均を要しない＝非定常・非周期に強い．本研究にはこちらが実際的かもしれない．
    \item \textbf{参照すべき桁は 0.08 ではなく 0.35--0.60}：0.08 は q-DNS 格子（265 万セル）の値．本研究の no-model 運転に対応するのは no-model LES 格子の値で，実効 Re の低下は 26--38~\%，局所的には 60-Grid 中央で $\nu_\mathrm{num} \approx \nu$（5 割超）．しかも本研究のメッシュ（水深方向 5~mm で 6--20 セル）は Komen の LES 格子より遥かに粗く，$\nu_\mathrm{num}/\nu$ が 1 を超える可能性は十分ある．測る価値が高い．
    \item \textbf{symmetric twin cell への含意}（本論文の主張ではなく推論）：数値散逸は等方的に効かず，pressure strain の過小評価を通じて\textbf{横方向変動を選択的に潰す}．蛇行や非対称モードは横方向の振動なので，主流ジェットを保ったまま非対称モードだけが消える，という機構がありうる．P09 でメッシュ依存性がないことは示せても，「どのメッシュでも数値散逸が十分大きく対称化している」可能性は排除できない．$\nu_\mathrm{num}$ を実測すれば区別がつく．
    \item \textbf{CFD と実験が別の枝にある件の診断}：上記「誤差の署名 4 点セット」と照合する．合致するなら数値散逸が第一容疑者，合致しないなら境界条件・幾何・物理モデルを探すべき，という切り分けができる．
    \item \textbf{設定への含意}：(i) 対流項に風上系（upwind, limitedLinear, vanLeer 等）を混ぜているならそれが $k$ を人工的に増やしている．(ii) CFL を下げると数値散逸が減る．(iii) P10 の $z$ 精密化は Fig.~17 の「アスペクト比より横断方向解像度」と整合し，方向性が支持される．(iv) SGS モデルの選択よりスキームと CFL に時間をかけるべき．
  \end{itemize}
\end{itemize}
\noindent（確認: p.565--595，全 31~p 通読）

\noindent\litref{Schranner et al. (2015)}\quad Komen と対をなす数値散逸定量法．\textbf{本研究で実際に使えるのはこちら}（全 14~p 通読）．
\begin{itemize}
  \item \textbf{問い}：任意のグリッドベース Navier--Stokes ソルバの実効数値散逸率・数値粘性を，ソルバ非依存の後処理として物理空間で定量するには．
  \item \textbf{手法}：\textbf{全運動エネルギーの輸送方程式の残差から測る}．エネルギー減衰率を 2 通りに計算し差を取る：(i) シミュレーションのエネルギーデータを直接時間微分，(ii) NS の空間項から評価．離散化誤差がなければ一致するはずの差が数値散逸 $\epsilon_n$（式 24--26）：
  \[ -\epsilon_n = \frac{dE_\mathrm{kin}}{dt} + F_\mathrm{ekin} + F_\mathrm{ac} - F_v - \Pi + \epsilon_v \]
  （$\Pi$ は非圧縮でゼロ）
  \begin{itemize}
    \item 数値粘性は $\nu_n = \epsilon_n/\mathcal{E}$（式 38--40）．$\mathcal{E}$ は\textbf{散逸関数}（式 14）であって $\epsilon_v$ ではない．$\epsilon_v = \nu\mathcal{E}$ なので $\nu\,\epsilon_n/\epsilon_v$ と等価＝Komen 式 (16) と同形．$\epsilon_n/\epsilon_v$ とすると無次元比になり粘性の次元にならない．\textbf{実装時の要注意点}．
    \item \textbf{必要なのは速度場 3 時刻分}．時間微分は $t \pm \Delta t$ の 2 次中心差分（式 59）だが空間項は $t_n$ で評価するため $t_{n-1}, t_n, t_{n+1}$ を保存する．セル面流束は 6 次中心で再構成（式 60）．
    \item \textbf{評価法は 2 通り}：(a) 各セルで式 (26) を評価し部分領域で積分，(b) 部分領域の\textbf{境界面でのみ}流束を評価．(b) は局所評価が不要で計算負荷が小さい．数学的・数値的には同一．
    \item \textbf{自己完結}：ソルバの出力（速度・圧力・密度）だけでよく追加情報が不要．周期境界も統計的定常も不要．検証は 3 次元 Taylor--Green 渦（層流→遷移→乱流を通過），$64^3$ セル，$\mathrm{Re} = 100$--$\infty$，ALDM モデル．
  \end{itemize}
  \item \textbf{主要知見}：
  \begin{itemize}
    \item 物理空間法とスペクトル空間法の一致は\textbf{相対誤差 $\pm 0.5$~\%}．小スケールが急速に生成される $t \approx 4$ 付近でわずかに増える程度．非周期・非一様・非構造格子に適用できる点がスペクトル法との決定的な差．
    \item \textbf{時間刻みへの依存は弱い}：$\Delta t$ を 20 倍振っても $\epsilon_n$ はほとんど変わらない（Fig.~3）．原文は「CFL $<1$ を満たし時間積分誤差が空間離散化誤差以下なら影響は無視でき，解析のためにとくに小さい $\Delta t$ を考慮する必要はない」＝既存の計算結果をそのまま後処理できる．
    \item \textbf{セル面再構成の次数も効かない}：2 次/4 次/6 次の差は最大 2.8~\%（それも $t \in [8, 9.5]$ の短期間のみ，大半の時間は同一）．理由は\textbf{手法自体が保存的だから，セル間流束の評価方法は問題にならない}．
    \item 任意の非周期部分領域で機能（CSD1--3, NCSD1--2 で検証）．ただし\textbf{小さい部分領域かつ低 Re では食い違いが目立つ}．同じ体積 $6^3$ セルでも $\epsilon_n$ が大きい CSD2 の方が一致が良く，数値散逸が小さい領域ほど相対誤差が大きい．
    \item \textbf{局所的には数値散逸が負にもなる}（Fig.~6 が $\mathrm{sgn}(\nu)$ で色分けした 3 次元可視化）．残差を数値散逸と解釈できるのは\textbf{十分大きい制御体積で散逸的寄与が分散的寄与を上回る}ときのみ（Komen も同じ論理を引用）．
    \item 完全周期領域では流束項が厳密に相殺し，全数値散逸は全運動エネルギーの変化率と全粘性散逸だけで決まる（式 61）．
  \end{itemize}
  \item \textbf{限界・適用範囲}：検証は\textbf{Taylor--Green 渦のみ}＝滑らかで周期的な単相流で，壁境界層も自由水面も剥離も含まない．ALDM という特定の ILES スキームでの検証．粘性散逸項自体も離散化誤差を含み，Appendix~A で議論されるが完全には除去できない（著者は「関心はコードが既に生む粘性散逸を超える分なので $\mathcal{E}$ の誤差は他項より重要でない」と整理）．\textbf{総量しか出ず，どの収支項が壊れているかは分からない}＝そこは Komen の TKE ベース手法の役割．VOF・二相流は想定外．
  \item \textbf{本研究との接点}：
  \begin{itemize}
    \item \textbf{本研究に適用できるのはこちら}．統計的定常も周期境界も要求せず瞬時場から局所評価できるため，充填過程という非定常問題にそのまま使える．Komen が要求する 180 FTT の統計平均は本研究では不可能．ソルバ改造も不要で，$\Delta t$ を小さくし直す必要も高次スキームで回し直す必要もない（上記 2 点より）．
    \item \textbf{部分領域ごとに測れる点が問いに直結}：ジェットのコア／再循環セル内部／低流速域／port 近傍で別々に $\nu_n$ を評価すれば，「数値散逸が非対称モードの住む場所に集中しているか」に直接答えうる．$\Omega$ マスクの妥当性の裏づけにもなる．ただし低 Re ＋小領域は精度が落ちるので，領域は大きめに取り複数サイズで測ってばらつきを見ること．
    \item \textbf{閉鎖 basin では流束項が大幅に簡略化される}：no-slip 壁では $u = 0$ なので $F_\mathrm{ekin}$ も $F_v$ も消える．残るのは\textbf{port 面の流束}のみ（＋自由表面が slip 条件なら粘性流束）．ただし port 近傍は流速最大で勾配も急なので，ここの評価精度が全体を左右する．
    \item \textbf{VOF での適用は未検証}．interFoam の界面圧縮項が運動エネルギー収支にどう入るかは本論文の枠組みに含まれない．界面近傍を除外して評価するなどの工夫が要るかもしれない．実装前に要検討．
    \item Komen との使い分け：Komen ＝乱流運動エネルギー基準・統計的定常向け・\textbf{機構の説明に強い}，Schranner ＝全運動エネルギー基準・非定常/非周期向け・\textbf{実装が軽い}．実務的には「Schranner で $\nu_n$ を測り，Komen の誤差の署名（平均速度過大／主流方向 RMS 過大／横方向 RMS 過小／$k$ 過大）で解釈する」という組み合わせになる．
  \end{itemize}
\end{itemize}
\noindent（確認: p.84--97，全 14~p 通読）

\noindent\litref{Deshpande et al. (2012)}\quad interFoam の標準的な評価論文（Open Access，全 36~p 通読）．
\begin{itemize}
  \item \textbf{なぜ読むか}：数値散逸を測る（Komen/Schranner）にせよ，その解釈には interFoam が運動量場に何をするかの理解が要る．また「laminar + interFoam」という組合せ固有の癖を知らないと，測った $\nu_\mathrm{num}$ をどこに帰属させるか判断できない．
  \item \textbf{問い}：interFoam は界面移流・慣性支配流・表面張力支配流のそれぞれでどの程度の性能をもつか．
  \item \textbf{手法}：\textbf{コードそのものを解析した}アルゴリズム記述（§2）＋検証テスト（§3）．テストは (i) 純移流，(ii) 慣性支配（定在重力波，液滴衝突），(iii) 表面張力支配（曲率精度，spurious current），(iv) 毛管流．OpenFOAM 1.5 で検証，2.1.1 まで有効と著者が明記．
  \item \textbf{主要知見}：
  \begin{itemize}
    \item \textbf{界面圧縮の仕組み}：MULES limiter $\lambda_M$ は界面近傍で 1，それ以外で 0（Fig.~2）．＝界面から離れた領域では $\gamma$ の移流は\textbf{単純な風上差分}（式 13）．界面近傍のみ圧縮流束と高次スキームが加わる（式 14）．$\gamma$ の面値は vanLeer，運動量は PISO 3 反復．
    \item \textbf{$C_\gamma$（cAlpha）は 1 を推奨}：「全計算で $C_\gamma = 1$ を用いた．増やしても減らしても界面曲率の誤差と界面のにじみを悪化させることが分かったため」＝既定値の根拠を引ける．
    \item \textbf{質量保存は極めて良好}：質量誤差 $1.3\times10^{-7}$~\%（level set 32--80~\%，CLSVOF 0.4~\% と桁違い，Table~8）．
    \item \textbf{移流精度}：Zalesak 円板では CLSVOF/ELVIRA と同程度（Table~3）．2D 渦場では THINC-SW より誤差が\textbf{3--5 倍}大きく収束率も低い（Table~6）．3D 渦場は「やや大きいが同等」（Table~7）＝2D の方が変形が激しいため．幾何再構成法には劣る．
    \item \textbf{非構造格子では界面に波紋}（Fig.~4）．構造格子では滑らか．本研究は六面体構造格子なので回避済み．
    \item 補間スキーム（linear/cubic）の違いは移流精度にほとんど影響しない（Table~4, 5）．
    \item \textbf{打ち切り誤差は格子細分で収束しない}：$\tilde{e}_\tau \sim 1/\Delta x$（Table~9，収束率が一貫して $-1.00$）．圧縮流束に置き換えても同じ．一方グローバル誤差指標は収束する．著者は「今後の研究の余地」と結論．
    \item \textbf{spurious current も格子細分で悪化する}：$C_{\Delta x} \sim \sqrt{\sigma/(\rho\Delta x)}$（式 70）．「格子細分化は緩和せず，むしろ悪化させる」と明記．密度比が大きいほど気相側で増幅される（式 68 の $1/\rho$）．
    \item \textbf{定在重力波（§3.2.1）が本研究に最も近い検証}：$H = L = 1$~m，密度比 1000，閉じた容器，slip 壁．$\lambda/\Delta x = 50$ という粗さでも解析解（Lamb）と周波数がよく一致．ただし Fig.~8 では 8 周期ほどで位相ずれが見え始める＝\textbf{長時間計算では位相誤差が蓄積する}．大半のケースを $\sigma = 0$ で実施．
    \item 曲率は格子を細かくしても 10~\% ずれたまま収束する（$\Delta P_\mathrm{exact} = 4$ に対し 3.6，§3.3.2）．
    \item \textbf{時間刻みの安定条件}：$\Delta t \leq \max(10\tau_\mu,\ 0.1\tau_\rho)$，$\tau_\mu = \mu\Delta x/\sigma$，$\tau_\rho = \sqrt{\rho\Delta x^3/\sigma}$（式 84）．$C_1 = 0.01$, $C_2 = 10$ は約 80 本の計算から決めた interFoam 固有の値．これは\textbf{表面張力の陽的処理に対する数値安定条件}であって，表面張力が物理的に効くかとは無関係．$\sigma \neq 0$ なら We が大きくても適用される．ただし著者は「移動界面では spurious current の生成は二次的なので保守的な見積もり」と留保．
  \end{itemize}
  \item \textbf{限界・適用範囲}：検証は理想化ベンチマークが中心．\textbf{乱流との相互作用は扱わない}．OpenFOAM 1.5 ベース．本研究のような\textbf{浅い自由表面流・充填過程}の検証例はない．打ち切り誤差の非収束は未解決のまま．
  \item \textbf{本研究との接点}：
  \begin{itemize}
    \item \textbf{Methods の interFoam 記述の出典}：アルゴリズム（MULES，界面圧縮，PISO，$C_\gamma = 1$）は本論文を引けばよい．使用バージョンを明記すること．
    \item \textbf{質量保存 100~\% の裏づけ}：\texttt{cfd/campaign.md} の「質量 100~\%」は VOF の構造上の性質で本論文の $1.3\times10^{-7}$~\% と整合．ただしこれは\textbf{$\gamma$ 移流の質量保存}であって運動量の保存精度ではない．
    \item \textbf{界面から遠い領域は風上差分}という事実は数値散逸の帰属を考える上で重要．basin 中央部の散逸は $\gamma$ 方程式ではなく\textbf{運動量方程式の対流項スキーム}に帰属すべき．もし \texttt{div(phi,U)} に \texttt{linearUpwind} 等の風上混合を使っているなら，Komen 2017 の Fig.~12（Van Leer が中心差分より顕著に悪化させ，$k$ を人工的に増やす）がそのまま当てはまる．第一に確認すべき設定．
    \item \textbf{時間刻みの再確認}：本研究の $\Delta x = 5$~mm で式 (84) を概算すると $\Delta t \leq$ 約 4.2~ms．CFL $=0.5$ 基準の約 21~ms より 5 倍厳しい．\texttt{adjustTimeStep} を CFL だけで回しているなら要確認（概算，要検算）．
    \item \textbf{$\sigma = 0$ という選択肢}：$\mathrm{Bo} = \rho g h^2/\sigma \approx 410$ で表面張力は物理的に無視できる．$\sigma = 0$ にすれば (i) 上記の時間刻み制約が消え，(ii) spurious current が原理的に消え，(iii) 曲率の 10~\% 誤差が入らず，(iv) Schranner の運動エネルギー収支に表面張力の仕事項を足す必要がなくなる．本論文の定在重力波検証自体が $\sigma = 0$ で行われており validated な構成．
    \item \textbf{細分化の効果は一方向でない}：§3.1.4（打ち切り誤差 $\sim 1/\Delta x$）と §3.3.3（spurious current が悪化）は細分化で悪化する側，Komen の数値散逸は細分化で改善する側．P09/P10 で「細かくしたら変わった」を単純に「収束した」と読めない．Limitations に書くべき論点．
    \item \textbf{Schranner 適用時の障害}：interFoam は修正圧力 $p_d = p - \rho\mathbf{g}\cdot\mathbf{x}$（式 28）で離散化するため，Schranner の圧力仕事項をそのまま取り出せない．加えて (i) 密度が $\gamma$ で変わる（式 31），(ii) 重力の仕事項 $\int \rho\mathbf{g}\cdot\mathbf{U}\,dV$ の追加が要る．表面張力項は $\sigma = 0$ なら不要．
  \end{itemize}
\end{itemize}
\noindent（確認: art.\ 014016，全 36~p 通読）

\noindent\litref{Pope (2004)}\quad LES の概念的基礎を問い直す古典（Open Access，全 24~p 通読）．
\begin{itemize}
  \item \textbf{なぜ読むか}：「底面摩擦／鉛直解像度」が \texttt{cfd/campaign.md} で\textbf{唯一まだ触っていない物理軸}とされている．5~mm 鉛直セルで水深 30--100~mm ＝ 6--20 セルが十分かを判断する基準が要る．加えて Komen/Schranner が「数値誤差をどう測るか」なのに対し，本論文は\textbf{なぜ測るのか，測った後どう論じるか}を与える．
  \item \textbf{問い}：LES の概念的基礎に関する 10 の問い．とくに (i) LES は物理モデルか数値手続きか，(ii) LES をどう「完全」にできるか，(iii) 統計量はフィルタ幅 $\Delta$ にどう依存するか，(iv) LES モデルをどう評価すべきか．問題提起が目的で「答えは決定的でも網羅的でもない」と明言．
  \item \textbf{主要知見}：
  \begin{itemize}
    \item \textbf{LES の 3 分類}：pure physical LES（陽なモデル＋数値誤差が無視できる）／physical LES（数値誤差あり）／numerical LES（$W$ の記述と発展が数値手法と本質的に結びつく．MILES 等）．numerical LES では \textbf{$\Delta = h$}（格子幅）と定義する．
    \item \textbf{「no model」という呼び方は不適切}．陽なモデルがなくても計算場はメッシュと数値手法の両方に依存する．MILES は「不適切な空間解像度の領域で安定に応答するよう数値手法を特別に構成」している．
    \item \textbf{最適な計算は非無視な数値誤差を含む}：§4「$\Delta/h$ の最適値は非無視な数値誤差に対応する可能性が高い」，§7「pure physical LES が有利とは考えにくい」，結論「最適なモデルが非無視な数値誤差を含む可能性は極めて高い」．理由は §9 の費用対効果：精度だけでなく\textbf{許容精度を達成する計算コスト}で比較すべきだから．
    \item \textbf{$M$ による解像度基準（式 1, 2）}：$M(x,t) = k_r/(K + k_r)$．$k_r$ は残差運動の TKE，$K$ は解像運動のそれ．$M = 0$ が DNS，$M = 1$ が RANS．$\epsilon_M = 0.2$ が「運動エネルギーの 80~\% を解像」に対応し，$M$ が $\epsilon_M$ を超える領域は細分，大きく下回る領域は粗くする（adaptive LES）．ただし\textbf{$k_r$ の推定法が要る}と著者も認める．
    \item \textbf{$\Delta$ 依存性の評価は LES 手法の必須要素}：「$\Delta(x)$ への依存を評価することが LES の一般的な作法になっていないのは遺憾」．\textbf{単一の $\Delta$ での比較は極めて誤解を招きうる}（Fig.~5 で 2 モデルの曲線が $\Delta^*$ で交差し，そこだけ見ると結論が逆になる）．
    \item \textbf{収束先は真値とは限らない}：$\Delta$ が慣性小領域にあるとき，収束先は真値 $Q$ ではなく\textbf{中間漸近 $Q_I^m$}（§7 (iii)(a)，Fig.~2）．ただし §7 (iv)：統計量 $Q$ とそれに影響する過程が\textbf{もっぱらエネルギー保有領域に属するなら単一漸近（$Q_I^m = Q$）}となり $Q^m$ は真値に収束する（Fig.~4）．さらに §7 (i)：中程度の Re では中間漸近が識別できないかもしれない，という留保つき．
    \item \textbf{$W$ と $U$ の関係は統計的でしかない}：解像速度場 $W$ はフィルタ済み $\bar{U}$ と等しくなり得ない（$\bar{U}$ はランダム場で将来発展が現在の状態で決まらない）．→a priori テストは疑わしい．
    \item \textbf{壁近傍は LES の弱点}：「粘性壁近傍領域には大きな渦は存在しない」（Bradshaw）．高 Re では壁近傍運動は解像できずモデル化が必須．LES が最も成功するのは\textbf{高 Re の自由せん断流}で，律速過程が解像された大スケールにある場合．
    \item 動的モデルの成功理由を「スケール相似性」ではなく\textbf{予測が $\Delta$ に依存しないようにする原理}から導き直し，標準的な $c_s$ の式（式 30）が本質的にそのまま得られることを示す（§10）．
  \end{itemize}
  \item \textbf{限界・適用範囲}：概念的な問題提起が目的で具体的な手順・数値は最小限．図はすべて概念スケッチで定量的な曲線ではない．例は均質乱流・チャネル流・自由せん断流など古典的な流れ．\textbf{VOF・自由表面・浅水流は扱わない}．高 Re・慣性小領域の存在を前提とする議論が中心．
  \item \textbf{本研究との接点}：
  \begin{itemize}
    \item \textbf{Methods の書き方}：laminar 設定は Pope の枠組みでは「モデルなし」ではなく\textbf{numerical LES}．「陽な SGS モデルを用いず，数値スキームの散逸が残差運動の効果を担う」と書くのが正確で，Komen の「数値散逸が SGS 散逸より 2.4--3.5 倍大きい」という測定結果とも整合する．
    \item \textbf{Limitations の枠組み}：「最適な計算は非無視な数値誤差を含む」は，数値誤差の存在自体が失格の証拠ではないという論拠になる．問題は誤差があることではなく\textbf{その大きさを特徴づけていないこと}．Schranner で $\nu_\mathrm{num}$ を測ることの正当化に直結する．
    \item \textbf{本研究の metrics は有利な側}：$E$, $\phi_\mathrm{lv}$, $I_\mathrm{asym}$, $I_\mathrm{circ}$ はいずれも basin スケールのエネルギー保有運動から決まる量で，§7 (iv) の「単一漸近＝真値に収束する」場合にあたる．燃焼や高 Re 壁流のように律速過程が解像スケール以下にある状況ではない．CFD の妥当性の論拠．
    \item \textbf{ただし中間漸近の理論はそのまま持ち込めない}：本研究は $\mathrm{Re}_h$ が $10^3$ オーダーで慣性小領域が明瞭でなく，準 2 次元で逆カスケードを持つため Kolmogorov 的でない．Figs.~2--5 の漸近の絵は引かず，\textbf{$\Delta$ 依存性を特徴づけよという方法論的主張（§7, §8）に限って引く}のが安全．
    \item \textbf{P09/P10 の位置づけ}：メッシュを振ること自体が Pope の求める作法．ただし「単一 $\Delta$ での比較は誤解を招く」（Fig.~5 の交差）ので，振った範囲の広さが結論の強さを決める．
    \item \textbf{壁近傍の論点は理由が違う}：Pope の議論は高 Re 乱流境界層のストリーク解像が対象で，本研究の底面境界層（$\mathrm{Re}_h \approx 10^3$ で遷移域かそれ以下）にはそのまま移せない．ただし解像度の懸念は正当：6--20 セルでは底面近傍の速度勾配＝$c_f$ が実質 1--2 セルで決まる．Chen \& Jirka の $S$ を通じて分岐の位置を左右するので要確認だが，引用は Pope ではなく境界層解像の議論として立てるべき．
    \item \textbf{$M$ を laminar 設定で見積もる代替案}：$k_r$ の推定には本来 SGS モデルが要るが，複数の $\Delta$ で走らせて\textbf{解像運動エネルギー $K$ の変化}を見れば，$\Delta$ と $\Delta/2$ の間のスケールが持つエネルギーを近似できる．$K$ がまだ増え続けていれば解像不足．SGS モデルに依存しない実際的な代替．
  \end{itemize}
\end{itemize}
\noindent（確認: art.\ 35，全 24~p 通読）

\subsection*{論文横断整理}

\subsection*{原稿転記メモ}

% ================================================================
\section{入手メモ（未入手・保留中の文献）}
% ================================================================

\noindent 読みたいが入手できていない文献を、再調査を避けるためここに記録する（入手でき次第、該当セクションへ移す）．

\begin{itemize}
  \item \textbf{Chu, Liu \& Altai (2004)} ``Friction and confinement effects on a shallow recirculating flow'', \textit{J. Environ. Eng. Sci.} 3(5), 463--475．
  \begin{itemize}
    \item 状況：大学経由で DL 不可（アクセス権なし）．2026-06-16 時点．
    \item 想定セクション：§5（摩擦数 $S = \lambda\Delta B/8H$ と閉じ込め効果の出典．Peltier 2014・Chagdali 2025 が引用）．
    \item 入手の代替：相互貸借（ILL）／著者（V.H.\ Chu, McGill）／内容が重なる入手容易な代替＝Babarutsi \& Chu (1991), Babarutsi, Ganoulis \& Chu (1989)．
    \item 二次情報（未確認）：浅水再循環流における摩擦長さスケールと閉じ込めの効果．
  \end{itemize}
  \item \textbf{Rajaratnam (1976)} \textit{Turbulent Jets}（Elsevier, 書籍）．
  \begin{itemize}
    \item 状況：未入手（書籍のため論文 DB になし）．2026-06-19 時点．
    \item 想定セクション：§1/§5（ポート吐出＝潜没・平面噴流の古典的基盤）．
    \item 入手の代替：大学図書館の電子書籍/OPAC・相互貸借（ILL）．噴流基礎は Dracos (1992)・Giger et al.\ (1991) で部分的に代替可．
  \end{itemize}
  \item \textbf{Raffel, Willert, Wereley \& Kompenhans} \textit{Particle Image Velocimetry: A Practical Guide}（Springer, 書籍）．
  \begin{itemize}
    \item 状況：未入手（書籍）．2026-06-19 時点．
    \item 想定セクション：§4（PIV 実務の標準リファレンス）．
    \item 入手の代替：大学図書館の電子書籍/OPAC．PIV 手法の出典は Adrian (1991)・Thielicke \& Stamhuis (2014) で代替可．
  \end{itemize}
  \item \textbf{Fischer, List, Koh, Imberger \& Brooks (1979)} \textit{Mixing in Inland and Coastal Waters}（Academic Press, 書籍）．
  \begin{itemize}
    \item 状況：未入手（書籍）．2026-07-26 時点．
    \item 想定セクション：§1（吸込み流＝シンク流の基礎）．
    \item \textbf{残る穴（2026-07-26 に一部解消）}：\texttt{discussion.tex} の看板主張のうち「理想シンクは渦度供給源をもたない」の部分は Imberger et al.\ (1976) で引けるようになった（§1，浮力・分布シンクの極限としての言明）．\textbf{未解決なのは「実際の吸込み流が距離とともにどう減衰し，鉛直方向にどんな構造をもつか」}の側で，49 本にこれを扱う文献がない（M\"uller の「吸込みは局所的」は深い系の話で，本論文が対比の相手にしているため支えにならない）．
    \item 入手の代替：大学図書館の電子書籍/OPAC．噴流側は Dracos (1992)・Giger (1991) で代替できるが\textbf{吸込み側の減衰構造の代替は現時点でない}．なお Imberger (1980) ``Selective withdrawal: a review'' は\textbf{成層取水の枠組みなので均質流体の本研究には届かない}（Imberger et al.\ 1976 のメモ参照）．主張(a)を厚くするだけなら流体力学の教科書（Batchelor, Lamb 等）のポテンシャルシンクで足りる．
  \end{itemize}
  \item \textbf{Knauss (ed., 1987)} \textit{Swirling Flow Problems at Intakes}（IAHR Hydraulic Structures Design Manual 1, Balkema, 書籍）．
  \begin{itemize}
    \item 状況：未入手（書籍）．2026-07-26 時点．
    \item 想定セクション：§1（取水口渦・限界潜没の標準リファレンス）．
    \item 入手の代替：\textbf{M\"oller et al.\ (2015) で相当部分を代替済み}（同論文が Knauss の限界潜没則を Fig.~8 で図示・比較しており，二次情報として参照できる）．ただし Knauss の原典の設計基準そのものを引く必要が生じた場合は要入手．
  \end{itemize}
  \item \textbf{Celik, Cehreli \& Yavuz (2005)} ``Index of resolution quality for large eddy simulations'', \textit{J. Fluids Eng.} 127(5), 949--958．
  \begin{itemize}
    \item 状況：未入手．2026-07-27 時点．
    \item 想定セクション：§9（LES 解像度の妥当性基準）．
    \item 入手の代替：\textbf{Pope (2004) で代替可能}．Pope の $M = k_r/(K+k_r)$（$\epsilon_M = 0.2$ で運動エネルギーの 80~\% を解像）が同じ役割を果たす．Celik の LES\_IQ は Pope の $M$ を実装しやすくした指数で，本質的な内容は重複する．本研究では \textbf{laminar 設定のため $k_r$ を直接出せない}という別の問題があり，どちらの指数を使うにせよ同じ困難に直面する．
  \end{itemize}
  \item \textbf{Jain, Ranga Raju \& Garde (1978)} ``Vortex formation at vertical pipe intakes'', \textit{J. Hydraul. Div.} ASCE 104(HY10), 1429--1445．
  \begin{itemize}
    \item 状況：未入手．2026-07-26 時点．
    \item 想定セクション：§1（鉛直管取水口の渦形成・限界潜没）．
    \item 入手の代替：M\"oller et al.\ (2015) の参考文献に収録されており，同論文経由で二次的に参照可．本研究のポートは水平矩形断面なので鉛直管の知見の適用範囲は限定的．
  \end{itemize}
\end{itemize}

\end{document}
